% Lutte contre les problèmes liés à la santé reproductive des adolescent(e)s — SVT 1ère D — Cours signé OnBuch+
% Programme officiel MINESEC. Compiler avec : tectonic ND09-sante-reproductive-adolescents-d.tex
\newcommand{\DOCMATIERE}{SVT}
\newcommand{\DOCNIVEAU}{1ère D}
\newcommand{\DOCMODULE}{Module 2 : L'Éducation à la Santé}
\newcommand{\DOCLECON}{ND09}
\newcommand{\DOCTITRE}{Lutte contre les problèmes liés à la santé reproductive des adolescent(e)s}
% OnBuch+ — préambule « cours signés » ENRICHI (compilé avec tectonic / XeLaTeX)
% Système visuel « Pop » d'OnBuch+ : fond crème (jamais blanc), bordures encre
% épaisses, ombres « dures » décalées, titres Archivo Black en capitales.
% Version MATHÉMATIQUES : graphes (pgfplots/tikz), boîtes pédagogiques
% (définition, propriété, méthode, attention, le savais-tu, exercice résolu,
% exercice, TP, bilan), page de garde et signature OnBuch+ ancrée en bas.
%
% Macros attendues dans main.tex AVANT \input{preamble} :
%   \DOCMATIERE  \DOCNIVEAU  \DOCMODULE  \DOCLECON (numéro)  \DOCTITRE
\documentclass[11pt]{article}

\usepackage{fontspec}
\usepackage[french]{babel}
\usepackage[a4paper,margin=2cm,top=3.4cm,bottom=2.8cm,headheight=1.4cm,footskip=1.3cm]{geometry}
\usepackage{amsmath,amssymb,amsfonts}
\usepackage{mathtools} % \xmapsto, \coloneqq... souvent utilisés par les modèles
\usepackage{cancel} % simplifications barrées (bilans de forces)
% \abs{z} (module, valeur absolue) et \norm{u} (norme), utilisés par les modèles
\providecommand{\abs}{}\let\abs\relax\DeclarePairedDelimiter{\abs}{\lvert}{\rvert}
\providecommand{\norm}{}\let\norm\relax\DeclarePairedDelimiter{\norm}{\lVert}{\rVert}
\usepackage[table]{xcolor} % [table] : fournit aussi \rowcolors (alternance de lignes)
\usepackage{pagecolor}
\usepackage{tikz}
\usetikzlibrary{arrows.meta,positioning,calc,shapes.geometric,shapes.misc,shapes.arrows,decorations.pathmorphing,decorations.pathreplacing,decorations.markings,patterns,shadows,mindmap,trees,fit,backgrounds,matrix,intersections,angles,quotes}
\usepackage{pgfplots}
\usepackage[version=4]{mhchem} % équations nucléaires \ce{^{235}_{92}U}
\usepackage[european,siunitx]{circuitikz} % schémas électriques
\pgfplotsset{compat=1.18}
\usepgfplotslibrary{fillbetween,groupplots} % aires entre courbes (calcul intégral)
\usepackage{enumitem}
\usepackage{fancyhdr}
% [most] charge toutes les bibliothèques tcolorbox (listings, minted, raster,
% documentation...) et gonfle inutilement la mémoire TeX sur un document long
% (~300 boîtes) : on ne charge que ce qui est réellement utilisé.
\usepackage[skins,breakable]{tcolorbox}
\usepackage{graphicx}
\usepackage{colortbl}
\usepackage{booktabs}
\usepackage{tabularx}
\usepackage{diagbox} % case d'en-tête diagonale des tableaux à double entrée
% Colonnes extensibles centrée / gauche / droite (C, L, R), souvent utilisées par les modèles
\newcolumntype{C}{>{\centering\arraybackslash}X}
\newcolumntype{L}{>{\raggedright\arraybackslash}X}
\newcolumntype{R}{>{\raggedleft\arraybackslash}X}
\usepackage{array}
\usepackage{multicol}
\usepackage{siunitx}
% parse-numbers=false : accepte aussi \SI{2,5 \times 10^{-3}}{...}
\sisetup{locale=FR,output-decimal-marker={,},group-minimum-digits=4,parse-numbers=false}
\DeclareSIUnit\ohm{\ensuremath{\symup{\Omega}}} % U+2126 absent de la police
\DeclareSIUnit\division{div} % réglages d'oscilloscope (ms/div, V/div)
\DeclareSIUnit\tr{tr} % tours (vitesse de rotation, ex. \SI{1500}{\tr\per\minute})
\DeclareSIUnit\ch{ch} % cheval-vapeur (puissance moteur, unité usuelle hors SI)
\DeclareSIUnit\franccfa{FCFA} % franc CFA (coûts, contexte camerounais)
\DeclareSIUnit\dioptre{\ensuremath{\delta}} % vergence d'une lentille (dioptrie)
\usepackage{qrcode}
% \degree : commande fréquemment utilisée par les modèles pour un angle ou une
% température en dehors de \SI{}{} (non fournie par les paquets chargés).
\providecommand{\degree}{^{\circ}}
% \qed (fin de démonstration), utilisé par les modèles sans amsthm
\providecommand{\qed}{\hfill$\square$}
% \barre : double barre des valeurs interdites dans un tableau de variations (tkz-tab)
\providecommand{\barre}{\ensuremath{\|}}
% environnement proof (amsthm non chargé) utilisé par les modèles
\ifdefined\proof\else\newenvironment{proof}[1][Démonstration]{\par\noindent\textit{#1.}\ }{\qed\par}\fi
\usepackage{lastpage}
\usepackage{hyperref}
\hypersetup{hidelinks}

% ============================================================
% Palette « Pop » OnBuch+ — mêmes hex que la classe P (plus_ui.dart)
% ============================================================
\definecolor{poporange}{HTML}{F59321}
\definecolor{popdark}{HTML}{E07A0C}
\definecolor{popcream}{HTML}{FFF6E8}
\definecolor{popink}{HTML}{1C1714}
\definecolor{poppurple}{HTML}{7A5AE0}
\definecolor{popgreen}{HTML}{2E9E6A}
\definecolor{poppink}{HTML}{E0457B}
\definecolor{popblue}{HTML}{2F7DE1}
\definecolor{popgold}{HTML}{C98A12}
\definecolor{popmuted}{HTML}{8A7F76}
\definecolor{popline}{HTML}{EADFD2}
\definecolor{popgray}{HTML}{8A7F76}
\colorlet{popgrey}{popgray}
% Alias tolérants (le modèle nomme parfois les couleurs différemment)
\colorlet{popwhite}{white}
\colorlet{popblack}{popink}
\colorlet{popred}{poppink}
\colorlet{popyellow}{popgold}
% Teintes claires pour fonds de tableaux / zones
\colorlet{popblueL}{popblue!10!white}
\colorlet{poporangeL}{poporange!14!white}
\colorlet{popgreenL}{popgreen!12!white}
\colorlet{poppurpleL}{poppurple!12!white}
\colorlet{poppinkL}{poppink!12!white}
\colorlet{popgoldL}{popgold!14!white}

\pagecolor{popcream}

% ============================================================
% Typographie
% ============================================================
\defaultfontfeatures{Ligatures=TeX}
\setmainfont{PlusJakartaSans-Medium.ttf}[Path=../../../fonts/,BoldFont=PlusJakartaSans-Bold.ttf]
% Maths en sans-serif (Fira Math) avec chiffres et lettres droites de Plus
% Jakarta, pour que les formules se fondent dans le texte.
\usepackage[math-style=ISO,bold-style=ISO]{unicode-math}
\setmathfont{FiraMath-Regular.otf}[Path=../../../fonts/]
\setmathfont{PlusJakartaSans-Medium.ttf}[Path=../../../fonts/,range={up/num,up/{Latin,latin},\mathup}]
\usepackage{icomma} % virgule décimale sans espace en mode math
% Caractères mathématiques Unicode absents des polices de texte (Plus Jakarta,
% Archivo Black) : rendus par leur équivalent mathématique, en texte comme en
% formule — sinon ils s'affichent en carré vide (ex. « Arithmétique dans ℤ »).
\usepackage{newunicodechar}
\newunicodechar{ℝ}{\ensuremath{\mathbb{R}}}
\newunicodechar{ℤ}{\ensuremath{\mathbb{Z}}}
\newunicodechar{ℂ}{\ensuremath{\mathbb{C}}}
\newunicodechar{ℕ}{\ensuremath{\mathbb{N}}}
\newunicodechar{ℚ}{\ensuremath{\mathbb{Q}}}
\newunicodechar{𝔻}{\ensuremath{\mathbb{D}}}
\newunicodechar{α}{\ensuremath{\alpha}}
\newunicodechar{β}{\ensuremath{\beta}}
\newunicodechar{γ}{\ensuremath{\gamma}}
\newunicodechar{δ}{\ensuremath{\delta}}
\newunicodechar{ε}{\ensuremath{\varepsilon}}
\newunicodechar{θ}{\ensuremath{\theta}}
\newunicodechar{λ}{\ensuremath{\lambda}}
\newunicodechar{μ}{\ensuremath{\mu}}
\newunicodechar{σ}{\ensuremath{\sigma}}
\newunicodechar{τ}{\ensuremath{\tau}}
\newunicodechar{φ}{\ensuremath{\varphi}}
\newunicodechar{ψ}{\ensuremath{\psi}}
\newunicodechar{ω}{\ensuremath{\omega}}
\newunicodechar{Σ}{\ensuremath{\Sigma}}
% majuscules produites par \MakeUppercase dans les titres (α-aminés -> Α)
\newunicodechar{Α}{\ensuremath{\alpha}}
\newunicodechar{Β}{\ensuremath{\beta}}
\newunicodechar{Γ}{\ensuremath{\Gamma}}
\newunicodechar{Δ}{\ensuremath{\Delta}}
\newunicodechar{Ε}{\ensuremath{\varepsilon}}
\newunicodechar{Θ}{\ensuremath{\Theta}}
\newunicodechar{Λ}{\ensuremath{\Lambda}}
\newunicodechar{Μ}{\ensuremath{\mu}}
\newunicodechar{Τ}{\ensuremath{\tau}}
\newunicodechar{Φ}{\ensuremath{\Phi}}
\newunicodechar{Ψ}{\ensuremath{\Psi}}
\newunicodechar{Ω}{\ensuremath{\Omega}}
\newunicodechar{↦}{\ensuremath{\mapsto}}
\newunicodechar{∈}{\ensuremath{\in}}
\newunicodechar{∉}{\ensuremath{\notin}}
\newunicodechar{∧}{\ensuremath{\wedge}}
\newunicodechar{⇔}{\ensuremath{\Leftrightarrow}}
\newunicodechar{⟺}{\ensuremath{\Longleftrightarrow}}
\newunicodechar{⇒}{\ensuremath{\Rightarrow}}
\newunicodechar{⟹}{\ensuremath{\Longrightarrow}}
\newunicodechar{∘}{\ensuremath{\circ}}
\newunicodechar{∪}{\ensuremath{\cup}}
\newunicodechar{∩}{\ensuremath{\cap}}
\newunicodechar{≡}{\ensuremath{\equiv}}
\newunicodechar{⊂}{\ensuremath{\subset}}
\newunicodechar{⊥}{\ensuremath{\perp}}
\newunicodechar{∃}{\ensuremath{\exists}}
\newunicodechar{∀}{\ensuremath{\forall}}
\newunicodechar{∛}{\ensuremath{\sqrt[3]{\,}}}
\newunicodechar{∜}{\ensuremath{\sqrt[4]{\,}}}
\newunicodechar{⃗}{\ensuremath{{}^{\to}}}
\newunicodechar{ⁿ}{\ifmmode^{n}\else\textsuperscript{\ensuremath{n}}\fi}
\newunicodechar{ˣ}{\ifmmode^{x}\else\textsuperscript{\ensuremath{x}}\fi}
\newunicodechar{ᵇ}{\ifmmode^{b}\else\textsuperscript{\ensuremath{b}}\fi}
\newunicodechar{ʳ}{\ifmmode^{r}\else\textsuperscript{\ensuremath{r}}\fi}
\newunicodechar{ᵘ}{\ifmmode^{u}\else\textsuperscript{\ensuremath{u}}\fi}
\newunicodechar{ʸ}{\ifmmode^{y}\else\textsuperscript{\ensuremath{y}}\fi}
\newunicodechar{ᵃ}{\ifmmode^{a}\else\textsuperscript{\ensuremath{a}}\fi}
\newunicodechar{ᵉ}{\ifmmode^{e}\else\textsuperscript{\ensuremath{e}}\fi}
\newunicodechar{ᵏ}{\ifmmode^{k}\else\textsuperscript{\ensuremath{k}}\fi}
\newunicodechar{ᵖ}{\ifmmode^{p}\else\textsuperscript{\ensuremath{p}}\fi}
\newunicodechar{ᴺ}{\ifmmode^{N}\else\textsuperscript{\ensuremath{N}}\fi}
\newunicodechar{ᶜ}{\ifmmode^{c}\else\textsuperscript{\ensuremath{c}}\fi}
\newunicodechar{ʰ}{\ifmmode^{h}\else\textsuperscript{\ensuremath{h}}\fi}
\newunicodechar{ᵠ}{\ifmmode^{q}\else\textsuperscript{\ensuremath{q}}\fi}
\newunicodechar{ₙ}{\ifmmode_{n}\else\textsubscript{\ensuremath{n}}\fi}
\newunicodechar{ₐ}{\ifmmode_{a}\else\textsubscript{\ensuremath{a}}\fi}
\newunicodechar{ᵢ}{\ifmmode_{i}\else\textsubscript{\ensuremath{i}}\fi}
\newunicodechar{ᵣ}{\ifmmode_{r}\else\textsubscript{\ensuremath{r}}\fi}
\newunicodechar{ₖ}{\ifmmode_{k}\else\textsubscript{\ensuremath{k}}\fi}
\newunicodechar{ᵦ}{\ifmmode_{\beta}\else\textsubscript{\ensuremath{\beta}}\fi}
\newunicodechar{ₓ}{\ifmmode_{x}\else\textsubscript{\ensuremath{x}}\fi}
\newfontfamily\popdisplay{ArchivoBlack-Regular.ttf}[Path=../../../fonts/]
\newfontfamily\poplogo{SpaceGrotesk-Bold.ttf}[Path=../../../fonts/]
\newfontfamily\popsemi{PlusJakartaSans-SemiBold.ttf}[Path=../../../fonts/]
\newfontfamily\popxbold{PlusJakartaSans-ExtraBold.ttf}[Path=../../../fonts/]
\newfontfamily\popxboldls{PlusJakartaSans-ExtraBold.ttf}[Path=../../../fonts/,LetterSpace=3.0]

\setlist[enumerate]{leftmargin=1.7em,itemsep=5pt,topsep=4pt}
\setlist[itemize]{leftmargin=1.7em,itemsep=4pt,topsep=4pt,label={\color{poporange}\textbullet}}
\setlength{\parindent}{0pt}
\setlength{\parskip}{5pt}
\renewcommand{\arraystretch}{1.25}

% pgfplots : style Pop par défaut
\pgfplotsset{
  every axis/.append style={axis line style={popink,line width=1pt},tick style={popink},
    grid style={popline,line width=0.5pt},label style={font=\small},tick label style={font=\footnotesize}},
  popcurve/.style={poporange,line width=1.8pt,no markers,smooth},
}
% Même style disponible dans un \draw TikZ simple (hors pgfplots)
\tikzset{popcurve/.style={poporange,line width=1.8pt}}
\tikzset{popgrid/.style={popline,very thin}}
\colorlet{popcurve}{poporange} % le nom du style est parfois utilisé comme couleur

% ============================================================
% Pill / squiggle
% ============================================================
\newcommand{\poppill}[3]{%
  \tikz[baseline=-0.55ex]\node[draw=popink,line width=1.2pt,rounded corners=2.6pt,%
    fill=#2,inner xsep=6.5pt,inner ysep=3pt]{%
    \popxboldls\fontsize{7}{7}\selectfont\color{#3}\MakeUppercase{#1}};%
}
\newcommand{\popsquiggle}[1]{%
  \tikz[baseline=-0.4ex]{
    \draw[#1,line width=2.6pt,line cap=round]
      plot [smooth] coordinates {(0,0.09) (0.34,0.01) (0.68,0.21) (1.02,0.01) (1.36,0.10)};
  }%
}
% Mot-clé surligné (marqueur orange sous le texte)
\newcommand{\cle}[1]{{\popxbold\color{popdark}#1}}

% ============================================================
% En-tête / pied
% ============================================================
\pagestyle{fancy}
\fancyhf{}
\renewcommand{\headrulewidth}{0pt}
\renewcommand{\footrulewidth}{0pt}
\lhead{\raisebox{-0.15ex}{\poplogo\fontsize{13}{13}\selectfont\color{popink}OnBuch\color{poporange}+}}
\rhead{\poppill{\DOCMATIERE\ \textbullet\ \DOCNIVEAU\ \textbullet\ Leçon \DOCLECON}{white}{popink}}
\lfoot{\popxbold\fontsize{7.6}{7.6}\selectfont\color{popmuted}\MakeUppercase{Cours signé OnBuch+}\ \textbullet\ {\popsemi\DOCTITRE}}
\rfoot{\tikz[baseline=-0.6ex]\node[rounded corners=8.5pt,fill=popink,minimum height=17pt,minimum width=17pt,inner xsep=5pt,inner sep=0]{\popxbold\fontsize{8}{8}\selectfont\color{popcream}\thepage\,/\,\pageref*{LastPage}};}

% ============================================================
% Titres
% ============================================================
\newcommand{\chaptitre}[2]{%
  \begin{tcolorbox}[enhanced,colback=poporange,colframe=popink,boxrule=2.6pt,arc=11pt,
      drop shadow=popink,left=15pt,right=15pt,top=12pt,bottom=11pt,before skip=3pt,after skip=18pt]
    \poppill{#2}{popink}{popcream}\par\vspace{9pt}
    {\raggedright\hyphenpenalty=10000\exhyphenpenalty=10000\popdisplay\fontsize{22}{24}\selectfont\color{popcream}\MakeUppercase{#1}\par}
  \end{tcolorbox}
}
% Section numérotée : pastille orange avec le numéro + titre Archivo
\newcounter{popsec}
\newcommand{\coursec}[1]{%
  \stepcounter{popsec}\setcounter{popsub}{0}%
  \par\vspace{16pt}\noindent
  \begin{minipage}{\linewidth}
  \tikz[baseline=-0.7ex]\node[draw=popink,line width=1.6pt,fill=poporange,rounded corners=4pt,minimum size=22pt,inner sep=0]{\popdisplay\fontsize{12}{12}\selectfont\color{popcream}\thepopsec};%
  \hspace{8pt}{\popdisplay\fontsize{15}{17}\selectfont\color{popink}\MakeUppercase{#1}}\par
  \vspace{4pt}\noindent{\color{poporange}\rule{36pt}{3.6pt}}
  \end{minipage}\par\nopagebreak\vspace{8pt}
}
\newcounter{popsub}
\newcommand{\courssub}[1]{%
  \stepcounter{popsub}%
  \par\vspace{9pt}\noindent{\popxbold\fontsize{11.5}{13}\selectfont\color{popink}{\color{poporange}\thepopsec.\thepopsub}\hspace{6pt}#1}\par\nopagebreak\vspace{4pt}
}

% ============================================================
% Boîtes pédagogiques — même recette : fond blanc, bordure épaisse colorée,
% ombre dure encre, badge plein. Titre optionnel [#1].
% ============================================================
\tcbset{popbase/.style={breakable,enhanced,colback=white,boxrule=2.3pt,arc=9pt,
  shadow={3pt}{-3pt}{0pt}{popink},left=14pt,right=14pt,top=11pt,bottom=11pt,before skip=11pt,after skip=15pt,
  fontupper=\popsemi\fontsize{10.6}{14.5}\selectfont\color{popink},
  fontlower=\popsemi\fontsize{10.6}{14.5}\selectfont\color{popink}}}
% Coche dessinée (le glyphe ✓ n'existe pas dans les polices du document)
\newcommand{\popcheck}[1]{\tikz[baseline=-0.5ex]\draw[#1,line width=1.6pt,line cap=round,line join=round](-0.13,0.0)--(-0.04,-0.09)--(0.13,0.1);}
\providecommand{\coche}{\popcheck{popgreen}}
\providecommand{\resultat}[1]{\par\smallskip\noindent\fcolorbox{popgreen}{popgreenL}{\parbox{\dimexpr\linewidth-2\fboxsep-2\fboxrule\relax}{\textbf{Résultat.} #1}}\par\smallskip}
\newcommand{\popboxhead}[3]{\poppill{#1}{#2}{white}\ifx\relax#3\relax\else\hspace{6pt}{\popxbold\fontsize{10.6}{12}\selectfont\color{popink}#3}\fi\par\vspace{7pt}}

\newtcolorbox{aretenir}[1][]{popbase,colframe=popgreen,before upper={\popboxhead{À retenir}{popgreen}{#1}}}
\newtcolorbox{exemplebox}[1][]{popbase,colframe=poporange,before upper={\popboxhead{Exemple}{poporange}{#1}}}
\newtcolorbox{definition}[1][]{popbase,colframe=popblue,before upper={\popboxhead{Définition}{popblue}{#1}}}
\newtcolorbox{propriete}[1][]{popbase,colframe=poppurple,before upper={\popboxhead{Propriété}{poppurple}{#1}}}
\newtcolorbox{methode}[1][]{popbase,colframe=popgold,before upper={\popboxhead{Méthode}{popgold}{#1}}}
\newtcolorbox{attention}[1][]{popbase,colframe=poppink,before upper={\popboxhead{Attention}{poppink}{#1}}}
\newtcolorbox{experience}[1][]{popbase,colframe=popblue,colback=popblueL,before upper={\popboxhead{Expérience}{popblue}{#1}}}
% « Le savais-tu ? » : carte violette pleine (contexte, culture, Cameroun)
\newtcolorbox{savaistu}[1][]{popbase,colback=poppurple,colframe=popink,
  fontupper=\popsemi\fontsize{10.4}{14.2}\selectfont\color{white},
  before upper={\poppill{Le savais-tu\,?}{white}{poppurple}\ifx\relax#1\relax\else\hspace{6pt}{\popxbold\color{white}#1}\fi\par\vspace{7pt}}}
% Exercice résolu : énoncé en haut, \tcblower puis la solution sur fond crème
\newcounter{popexo}\newcounter{popexr}
\newtcolorbox{exoresolu}[1][]{popbase,colframe=poporange,colbacklower=popcream,
  segmentation style={popink,line width=1.2pt,dashed},
  before upper={\stepcounter{popexr}\popboxhead{Exercice résolu \thepopexr}{poporange}{#1}},
  before lower={\poppill{Solution}{popink}{popcream}\par\vspace{6pt}}}
% Exercice (non résolu en place) — la correction est dans la partie Corrigés
\newtcolorbox{exercice}[1][]{popbase,colframe=popink,
  before upper={\stepcounter{popexo}\popboxhead{Exercice \thepopexo}{popink}{#1}}}
% Corrigé
\newtcolorbox{corrige}[1][]{popbase,colframe=popgreen,colback=popgreenL,
  before upper={\popboxhead{Corrigé}{popgreen}{#1}}}
% Objectifs / prérequis / bilan
\newtcolorbox{objectifs}{popbase,colframe=popink,colback=poporangeL,
  before upper={\popboxhead{Objectifs de la leçon}{popink}{}}}
\newtcolorbox{prerequis}{popbase,colframe=popink,colback=popblueL,
  before upper={\popboxhead{Prérequis}{popblue}{}}}
\newtcolorbox{bilan}{popbase,colframe=popink,colback=white,boxrule=2.8pt,
  before upper={\popboxhead{Fiche bilan}{poporange}{L'essentiel à connaître pour l'examen}}}
% Figure encadrée avec légende
\newtcolorbox{popfigure}[1][]{enhanced,breakable=false,colback=white,colframe=popline,boxrule=1.4pt,arc=8pt,
  left=10pt,right=10pt,top=10pt,bottom=8pt,before skip=10pt,after skip=12pt,halign=center,
  lower separated=false,halign lower=center,
  fontlower=\popsemi\fontsize{9}{11.5}\selectfont\color{popmuted},
  #1}
\newcommand{\legende}[1]{\par\vspace{6pt}{\popsemi\fontsize{9}{11.5}\selectfont\color{popmuted}#1}}

% ============================================================
% Signature OnBuch+
% ============================================================
\newcommand{\poplogoinline}[1]{{\poplogo\fontsize{#1}{#1}\selectfont\color{popink}OnBuch\color{poporange}+}}
\newcommand{\signatureonbuch}{%
  \begin{tcolorbox}[enhanced,colback=popink,colframe=popink,boxrule=2.2pt,arc=10pt,
      drop shadow=poporange,left=14pt,right=14pt,top=10pt,bottom=10pt,before skip=0pt,after skip=0pt]
    \tikz[baseline=-0.7ex]\node[circle,fill=popgreen,draw=popcream,line width=1.3pt,minimum size=22pt,inner sep=0]{\popcheck{white}};%
    \hspace{9pt}\begin{minipage}[c]{0.62\linewidth}
      {\popxbold\fontsize{12}{13}\selectfont\color{popcream}Signé\ \textbullet\ {\poplogo OnBuch\color{poporange}+}}\par\vspace{2pt}
      {\raggedright\popsemi\fontsize{8}{10.5}\selectfont\color{popline}\MakeUppercase{Équipe pédagogique OnBuch+}\\\MakeUppercase{Conforme au programme officiel MINESEC}\par}
    \end{minipage}\hfill
    {\poplogo\fontsize{22}{22}\selectfont\color{popcream}OnBuch\color{poporange}+}
  \end{tcolorbox}%
}

% ============================================================
% Page de garde
% #1 titre · #2 matière · #3 niveau · #4 module · #5 n° leçon · #6 durée ·
% #7 description / objectifs (texte ou itemize)
% ============================================================
\newcommand{\pagegarde}[7]{%
  \thispagestyle{empty}
  \newgeometry{a4paper,margin=1.7cm,top=1.6cm,bottom=1.4cm}%
  \begin{tikzpicture}[remember picture,overlay]
    \fill[poporange!18] ([xshift=-3.2cm,yshift=-3.6cm]current page.north east) circle (4.6cm);
    \fill[poppurple!14] ([xshift=2.4cm,yshift=7.6cm]current page.south west) circle (3.8cm);
  \end{tikzpicture}%
  {\poplogo\fontsize{30}{30}\selectfont\color{popink}OnBuch\color{poporange}+}\hfill
  \raisebox{0.6ex}{\poppill{Cours signé \textbullet\ Premium}{popink}{popcream}}\par
  \vspace{1pt}\popsquiggle{poppurple}\par
  \vspace{8pt}
  \begin{tcolorbox}[enhanced,colback=poporange,colframe=popink,boxrule=2.8pt,arc=13pt,
      drop shadow=popink,left=18pt,right=18pt,top=12pt,bottom=13pt,after skip=10pt]
    \poppill{#2 \textbullet\ #3}{popink}{popcream}\hspace{5pt}\poppill{#4}{white}{popink}\par\vspace{8pt}
    {\popdisplay\fontsize{14}{14}\selectfont\color{popink}LEÇON #5}\par\vspace{4pt}
    {\raggedright\hyphenpenalty=10000\exhyphenpenalty=10000\popdisplay\fontsize{25}{27}\selectfont\color{popcream}\MakeUppercase{#1}\par}
    \vspace{8pt}
    {\popxbold\fontsize{9.5}{9.5}\selectfont\color{popink}\MakeUppercase{Durée conseillée : #6}}
  \end{tcolorbox}
  \begin{tcolorbox}[enhanced,colback=white,colframe=popink,boxrule=2.2pt,arc=10pt,
      drop shadow=popink,left=16pt,right=16pt,top=10pt,bottom=10pt,after skip=8pt]
    \poppill{Dans cette leçon}{poppurple}{white}\par\vspace{5pt}
    \popsemi\fontsize{9.3}{12.6}\selectfont\color{popink}#7
  \end{tcolorbox}
  \noindent\begin{minipage}[t]{0.487\linewidth}
    \begin{tcolorbox}[enhanced,colback=popgreen,colframe=popink,boxrule=2.2pt,arc=10pt,
        drop shadow=popink,left=11pt,right=11pt,top=10pt,bottom=10pt,height=3.1cm,valign=center]
      \tcbox[colback=white,colframe=popink,boxrule=1.3pt,arc=5pt,boxsep=0pt,nobeforeafter,
        left=3pt,right=3pt,top=3pt,bottom=3pt,valign=center]{\qrcode[height=1.9cm]{https://play.google.com/store/apps/details?id=com.luvvix.onbuch&hl=fr}}%
      \hspace{8pt}\begin{minipage}[c]{0.5\linewidth}
        {\popdisplay\fontsize{11}{12.5}\selectfont\color{white}\MakeUppercase{Télécharge OnBuch}}\par\vspace{4pt}
        {\raggedright\popsemi\fontsize{8.4}{11}\selectfont\color{white}Gratuit \textbullet\ Google Play\\Tuteur IA, annales, concours, résultats\par}
      \end{minipage}
    \end{tcolorbox}
  \end{minipage}\hfill
  \begin{minipage}[t]{0.487\linewidth}
    \begin{tcolorbox}[enhanced,colback=poppurple,colframe=popink,boxrule=2.2pt,arc=10pt,
        drop shadow=popink,left=11pt,right=11pt,top=10pt,bottom=10pt,height=3.1cm,valign=center]
      \tikz[baseline=-0.7ex]\node[circle,fill=white,draw=popink,line width=1.3pt,minimum size=1.6cm,inner sep=0]{\popdisplay\fontsize{24}{24}\selectfont\color{poppurple}+};%
      \hspace{8pt}\begin{minipage}[c]{0.6\linewidth}
        {\popdisplay\fontsize{11}{12.5}\selectfont\color{white}\MakeUppercase{Passe à OnBuch+}}\par\vspace{4pt}
        {\raggedright\popsemi\fontsize{8.4}{11}\selectfont\color{white}Tous les cours signés, Tuteur IA illimité, fiches PDF — onglet OnBuch+ de l'app\par}
      \end{minipage}
    \end{tcolorbox}
  \end{minipage}
  \vspace{8pt}
  \popcontenu
  \vfill
  \signatureonbuch
  \restoregeometry
  \newpage
}

% Rangée « Ce cours contient » (page de garde)
\newcommand{\poptile}[3]{% #1 couleur #2 gros texte #3 légende
  \begin{tcolorbox}[enhanced,colback=#1,colframe=popink,boxrule=2pt,arc=9pt,shadow={2.5pt}{-2.5pt}{0pt}{popink},
      left=6pt,right=6pt,top=9pt,bottom=9pt,halign=center,valign=center,height=1.75cm,width=\linewidth,nobeforeafter]
    {\popdisplay\fontsize{12.5}{13}\selectfont\color{white}\MakeUppercase{#2}}\par\vspace{3pt}
    {\centering\popsemi\fontsize{7.8}{9.5}\selectfont\color{white}#3\par}
  \end{tcolorbox}}
\newcommand{\popcontenu}{%
  \par\noindent{\popxboldls\fontsize{7.5}{7.5}\selectfont\color{popmuted}CE COURS SIGNÉ CONTIENT}\par\vspace{7pt}
  \noindent\begin{minipage}[t]{0.235\linewidth}\poptile{popblue}{Cours}{complet, illustré, pas à pas}\end{minipage}\hfill
  \begin{minipage}[t]{0.235\linewidth}\poptile{poporange}{Méthodes}{et exercices résolus}\end{minipage}\hfill
  \begin{minipage}[t]{0.235\linewidth}\poptile{poppink}{Exercices}{type examen + corrigés}\end{minipage}\hfill
  \begin{minipage}[t]{0.235\linewidth}\poptile{popgreen}{Fiche bilan}{l'essentiel à retenir}\end{minipage}\par}

% Sommaire
\newtcolorbox{sommaire}{enhanced,colback=white,colframe=popink,boxrule=2.2pt,arc=10pt,
  drop shadow=popink,left=18pt,right=18pt,top=13pt,bottom=13pt,before skip=6pt,after skip=20pt,
  before upper={\poppill{Sommaire}{poppurple}{white}\par\vspace{9pt}\popsemi\fontsize{10.6}{14.5}\selectfont\color{popink}}}

% Signature de fin de cours — poussée tout en bas de la dernière page
\newcommand{\findecours}{%
  \par\vfill\nopagebreak
  \begin{center}\popsquiggle{poporange}\end{center}\vspace{4pt}
  \signatureonbuch
}
\AtBeginDocument{\def\llbracket{\mathopen{[\mkern-3.5mu[}}\def\rrbracket{\mathclose{]\mkern-3.5mu]}}} % intervalles d'entiers (glyphe ⟦ absent de la police)
% Glyphes absents de Fira Math : redéfinis à partir de glyphes disponibles
\AtBeginDocument{%
  \def\setminus{\mathbin{\backslash}}\let\smallsetminus\setminus
  \def\vdots{\mathord{\vbox{\baselineskip3.2pt\lineskiplimit0pt\kern4pt\hbox{.}\hbox{.}\hbox{.}}}}%
  \def\ddots{\mathinner{\mkern1mu\raise6.5pt\hbox{.}\mkern2mu\raise3.6pt\hbox{.}\mkern2mu\raise0.7pt\hbox{.}\mkern1mu}}%
  \def\triangle{\mathord{\scalebox{0.9}{$\symup{\Delta}$}}}%
  \def\nabla{\mathord{\raisebox{\depth}{\scalebox{1}[-1]{$\symup{\Delta}$}}}}%
  \def\checkmark{\popcheck{popdark}}%
  \def\star{\mathbin{\ast}}\def\bigstar{\mathord{\ast}}%
  \def\diamond{\mathbin{\scalebox{0.6}{$\Diamond$}}}%
  \def\bigcup{\mathop{\mathchoice{\raisebox{-0.3em}{\scalebox{1.7}{$\cup$}}}{\scalebox{1.25}{$\cup$}}{\cup}{\cup}}\displaylimits}%
  \def\bigcap{\mathop{\mathchoice{\raisebox{-0.3em}{\scalebox{1.7}{$\cap$}}}{\scalebox{1.25}{$\cap$}}{\cap}{\cap}}\displaylimits}%
  \def\lesseqqgtr{\mathrel{\lessgtr}}\def\bigoplus{\mathop{\oplus}\displaylimits}%
}
\providecommand{\ssi}{si et seulement si} % abréviation fréquente
\AtBeginDocument{\providecommand{\wideparen}{\overparen}} % arc de cercle

%% FIGFIX-BEGIN v5 — correctifs globaux des figures (docs/onbuch-generation/tools/figfix)
\usepackage{adjustbox}\usepackage{etoolbox}\tcbuselibrary{raster}
% toute figure TikZ/pgfplots plus large que la ligne est réduite à la largeur disponible
\BeforeBeginEnvironment{tikzpicture}{\begin{adjustbox}{max width=\linewidth}}
\AfterEndEnvironment{tikzpicture}{\end{adjustbox}}
% légendes pgfplots lisibles par-dessus les courbes
\pgfplotsset{every axis/.append style={legend style={fill=white,fill opacity=0.92,text opacity=1,draw=black!15,inner sep=3pt}}}
% étiquettes d'arêtes (syntaxe edge["..."]) avec fond blanc
\tikzset{every edge quotes/.append style={fill=white,inner sep=1.2pt}}
%% FIGFIX-END
\begin{document}

\pagegarde{Lutte contre les problèmes liés à la santé reproductive des adolescent(e)s}{SVT}{1ère D}{Module 2 : L'Éducation à la Santé}{ND09}{8 h}{Cette leçon étudie la santé reproductive des adolescents : fonctionnement de l'appareil reproducteur et puberté, causes et conséquences des grossesses précoces, infections sexuellement transmissibles et leur prévention, méthodes de contraception et planification familiale. Elle vise à doter l'élève de connaissances scientifiques et de comportements responsables pour préserver sa santé et réussir sa scolarité.
\vspace{4pt}
\begin{multicols}{2}\small
\begin{itemize}
\item À la fin de cette leçon, je sais décrire les transformations de la puberté chez le garçon et chez la fille.
\item À la fin de cette leçon, je sais expliquer les mécanismes de la fécondation et du cycle menstruel.
\item À la fin de cette leçon, je sais expliquer les risques liés aux grossesses précoces et non désirées à partir d'un cas concret.
\item À la fin de cette leçon, je sais identifier les principales IST, leurs modes de transmission et leurs conséquences.
\item À la fin de cette leçon, je sais justifier l'utilisation des moyens de prévention des IST (abstinence, fidélité, préservatif).
\item À la fin de cette leçon, je sais comparer les principales méthodes de contraception (naturelles, mécaniques, chimiques, hormonales).
\end{itemize}\end{multicols}}

\begin{sommaire}
\begin{multicols}{2}
\begin{enumerate}
\item Physiologie de la reproduction humaine et puberté
\item Grossesses précoces et non désirées
\item Infections sexuellement transmissibles (IST)
\item Méthodes de contraception
\item Planification familiale et comportement responsable
\item Activité d'intégration
\item Méthodes et exercices résolus
\item Exercices
\item Corrigés des exercices
\item Fiche bilan
\end{enumerate}
\end{multicols}
\end{sommaire}

\begin{objectifs}
\begin{itemize}
\item À la fin de cette leçon, je sais décrire les transformations de la puberté chez le garçon et chez la fille.
\item À la fin de cette leçon, je sais expliquer les mécanismes de la fécondation et du cycle menstruel.
\item À la fin de cette leçon, je sais expliquer les risques liés aux grossesses précoces et non désirées à partir d'un cas concret.
\item À la fin de cette leçon, je sais identifier les principales IST, leurs modes de transmission et leurs conséquences.
\item À la fin de cette leçon, je sais justifier l'utilisation des moyens de prévention des IST (abstinence, fidélité, préservatif).
\item À la fin de cette leçon, je sais comparer les principales méthodes de contraception (naturelles, mécaniques, chimiques, hormonales).
\item À la fin de cette leçon, je sais définir la planification familiale et expliquer son intérêt pour la famille et la société.
\item À la fin de cette leçon, je sais adopter et promouvoir un comportement responsable en matière de santé reproductive.
\end{itemize}
\end{objectifs}

\begin{prerequis}
\begin{itemize}
\item Notions sur les cellules reproductrices (gamètes) vues en classe de Seconde
\item Notion d'hormone et de régulation hormonale (début d'année de Première)
\item Notions sur les microbes pathogènes et les modes de contamination (hygiène, Seconde)
\item Anatomie générale de l'appareil génital humain (Seconde)
\end{itemize}
\end{prerequis}

\newpage

\coursec{Physiologie de la reproduction humaine et puberté}

Dans un quartier de Yaoundé, Amina, 16 ans, élève en classe de Première, découvre qu'elle est enceinte. Sa scolarité est menacée, sa santé aussi, et son entourage se déchire. Au même moment, son frère aîné apprend qu'il a contracté une infection sexuellement transmissible qu'il aurait pu éviter. Ces deux drames, hélas trop fréquents dans notre pays, ne sont pas des fatalités : ils reposent souvent sur une méconnaissance du fonctionnement du corps humain. Comment ces situations auraient-elles pu être prévenues ? Quelles connaissances scientifiques et quels comportements responsables permettent à un adolescent de protéger sa santé reproductive et son avenir ?

Pour répondre à ces questions, il faut d'abord comprendre \cle{comment fonctionne la reproduction humaine} : c'est l'objet de cette première section. Nous verrons ensuite, dans les sections suivantes, comment ces connaissances permettent de prévenir les grossesses précoces et les infections sexuellement transmissibles.

\courssub{Anatomie des appareils génitaux masculin et féminin}

\textbf{Observation préalable.} Pourquoi une grossesse débute-t-elle toujours dans le corps de la femme, alors que l'homme et la femme participent tous deux à la reproduction ? Parce que les deux sexes possèdent des organes spécialisés et complémentaires : les uns produisent les \cle{gamètes} (cellules reproductrices), les autres accueillent leur rencontre et le développement de l'embryon.

\begin{definition}[Appareil génital]
L'\cle{appareil génital} est l'ensemble des organes qui assurent la fonction de reproduction. Il comprend des \cle{gonades} (glandes produisant les gamètes et les hormones sexuelles), des \cle{voies génitales} (conduits d'évacuation ou de transport des gamètes) et des organes annexes.
\end{definition}

\textbf{L'appareil génital masculin} comprend :
\begin{itemize}
\item les \cle{testicules} (les gonades mâles), logés dans une poche cutanée appelée \emph{scrotum} (bourses) : ils produisent les \cle{spermatozoïdes} et sécrètent la \cle{testostérone}. Leur position externe maintient une température d'environ \SI{35}{\celsius}, inférieure à celle du corps (\SI{37}{\celsius}), indispensable à la spermatogenèse ;
\item les \cle{épididymes}, canaux enroulés accolés aux testicules, où les spermatozoïdes achèvent leur maturation et acquièrent leur mobilité ;
\item les \cle{canaux déférents}, qui transportent les spermatozoïdes depuis l'épididyme vers l'urètre ;
\item les glandes annexes (\cle{vésicules séminales} et \cle{prostate}) qui sécrètent le liquide séminal ; le mélange spermatozoïdes + liquide séminal forme le \cle{sperme} ;
\item l'\cle{urètre}, canal traversant le \cle{pénis}, qui assure l'évacuation du sperme lors de l'éjaculation (et de l'urine lors de la miction, mais jamais en même temps).
\end{itemize}

\textbf{L'appareil génital féminin} comprend :
\begin{itemize}
\item les \cle{ovaires} (les gonades femelles), situés dans la cavité abdominale : ils contiennent la réserve de follicules ovariens et sécrètent les \cle{œstrogènes} et la \cle{progestérone} ;
\item les \cle{trompes de Fallope} (ou trompes utérines), deux conduits reliant chaque ovaire à l'utérus : leur pavillon frangé capte l'ovule ; c'est \textbf{dans le tiers externe de la trompe (ampoule)} qu'a lieu la fécondation ;
\item l'\cle{utérus}, organe musculeux creux dont la paroi interne, la \cle{muqueuse utérine} (ou endomètre), s'épaissit cycliquement pour accueillir l'embryon ; sa partie inférieure rétrécie est le \cle{col de l'utérus} ;
\item le \cle{vagin}, conduit musculeux reliant le col de l'utérus à la vulve : il reçoit le sperme lors du rapport sexuel et constitue la voie d'expulsion du fœtus lors de l'accouchement.
\end{itemize}

\begin{popfigure}
\begin{center}
\begin{tikzpicture}[scale=0.9, every node/.style={font=\small}]
% ===== APPAREIL MASCULIN (vue sagittale médiane) =====
\begin{scope}[xshift=0cm]
\node[font=\small\bfseries, popblue] at (2.5,5.8) {Appareil génital masculin (coupe sagittale)};
% Vessie
\draw[fill=popblue!20, draw=popblue, thick] (1.5,4.3) ellipse (0.8 and 0.5);
\node[font=\scriptsize] at (1.5,4.3) {Vessie};
% Vésicules séminales
\draw[fill=popgold!30, draw=popgold, thick] (2.5,4.3) ellipse (0.4 and 0.6);
\node[font=\scriptsize] at (2.5,4.3) {Vés. sém.};
% Prostate
\draw[fill=popgold!50, draw=popgold, thick] (1.8,3.0) circle (0.35);
\node[font=\scriptsize] at (1.8,3.0) {Prostate};
% Canal déférent (gauche)
\draw[thick, poppurple] (2.9,4.3) .. controls (3.5,4.2) and (3.5,3.2) .. (3.0,2.5)
.. controls (2.8,2.0) and (2.5,1.8) .. (2.2,1.6);
% Urètre
\draw[thick, popdark] (1.5,3.0) -- (1.5,1.5) -- (0.8,1.2) -- (0.5,0.5);
% Pénis (corps caverneux/spongieux simplifié)
\draw[fill=popgreen!20, draw=popgreen, thick] (0.5,1.5) -- (0.5,0.5) -- (1.5,0.5) -- (1.5,1.5) -- cycle;
% Testicule + Épididyme
\draw[fill=poppink!20, draw=poppink, thick] (2.5,0.8) ellipse (0.5 and 0.7);
\node[font=\scriptsize] at (2.5,0.8) {Testicule};
\draw[thick, poppink] (3.0,1.0) .. controls (3.4,1.2) and (3.4,0.4) .. (3.0,0.6);
\node[font=\scriptsize, poppink] at (3.6,0.8) {Épididyme};
% Légendes avec lignes de rappel
\draw[thin, popmuted, dashed] (1.5,3.8) -- (-0.5,4.5) node[left, font=\scriptsize, align=right] {Vessie};
\draw[thin, popmuted, dashed] (2.5,4.3) -- (4.8,4.8) node[right, font=\scriptsize] {Vésicules séminales};
\draw[thin, popmuted, dashed] (1.8,3.0) -- (4.8,3.5) node[right, font=\scriptsize] {Prostate};
\draw[thin, popmuted, dashed] (3.0,3.5) -- (4.8,2.5) node[right, font=\scriptsize] {Canal déférent};
\draw[thin, popmuted, dashed] (1.0,1.0) -- (-0.5,0.2) node[left, font=\scriptsize, align=right] {Urètre \\ dans le pénis};
\draw[thin, popmuted, dashed] (2.5,0.8) -- (4.8,0.5) node[right, font=\scriptsize, align=center]{Testicule \\(spermatozoïdes, testostérone)};
\end{scope}
% ===== APPAREIL FÉMININ (vue sagittale médiane) =====
\begin{scope}[xshift=9.5cm]
\node[font=\small\bfseries, poppink] at (2.5,5.8) {Appareil génital féminin (coupe sagittale)};
% Utérus (paroi + endomètre)
\draw[fill=poppink!20, draw=poppink, thick] (2.2,3.8) ellipse (0.9 and 1.1);
\draw[fill=poppink!40, draw=poppink, thick] (2.2,3.8) ellipse (0.5 and 0.7);
\node[font=\scriptsize] at (2.2,3.8) {Endomètre};
% Col de l'utérus
\draw[thick, poppink] (1.3,2.7) -- (1.3,1.8); \draw[thick, poppink] (3.1,2.7) -- (3.1,1.8);
% Vagin
\draw[fill=popgreen!20, draw=popgreen, thick] (1.3,1.8) -- (0.9,0.2) -- (3.5,0.2) -- (3.1,1.8) -- cycle;
% Trompe (gauche)
\draw[thick, poppurple] (3.1,4.5) .. controls (4.2,5.0) and (5.0,4.8) .. (5.2,4.2);
% Pavillon
\draw[thick, poppurple, decorate, decoration={snake, amplitude=0.5pt, segment length=2pt}] (5.2,4.2) -- (5.5,4.2);
% Ovaire
\draw[fill=popgold!30, draw=popgold, thick] (5.0,3.2) ellipse (0.4 and 0.55);
\node[font=\scriptsize] at (5.0,3.2) {Ovaire};
% Légendes
\draw[thin, popmuted, dashed] (2.2,4.9) -- (-0.5,5.5) node[left, font=\scriptsize, align=right] {Utérus \\(corps \& endomètre)};
\draw[thin, popmuted, dashed] (2.2,2.7) -- (-0.5,2.5) node[left, font=\scriptsize, align=right] {Col de \\ l'utérus};
\draw[thin, popmuted, dashed] (1.1,1.0) -- (-0.5,0.5) node[left, font=\scriptsize, align=right] {Vagin};
\draw[thin, popmuted, dashed] (4.0,5.0) -- (6.5,5.5) node[right, font=\scriptsize, align=center]{Trompe de Fallope \\(site de la fécondation)};
\draw[thin, popmuted, dashed] (5.0,3.2) -- (6.5,3.0) node[right, font=\scriptsize, align=center]{Ovaire \\(ovules, hormones)};
\end{scope}
\end{tikzpicture}
\end{center}
\legende{Schémas simplifiés des appareils génitaux masculin (à gauche) et féminin (à droite), en coupe sagittale médiane. Les gonades (testicules, ovaires) produisent les gamètes et les hormones sexuelles ; les voies génitales assurent le transport des gamètes jusqu'à leur rencontre dans l'ampoule de la trompe.}
\end{popfigure}

\begin{attention}[Ne pas confondre -- Pièges classiques au Probatoire]
\begin{itemize}
\item Chez la femme, l'orifice urinaire (méat urinaire) et l'orifice génital (vagin) sont \textbf{distincts} (la femme urine et a des rapports par des orifices différents) ; chez l'homme, l'urètre est un conduit commun mais l'urine et le sperme ne passent \textbf{jamais en même temps} (le sphincter vésical interne ferme la vessie pendant l'éjaculation).
\item La fécondation n'a \textbf{pas} lieu dans l'utérus, mais dans le \textbf{tiers externe de la trompe de Fallope (ampoule)}. L'utérus est le site de la nidation et de la grossesse.
\item L'ovaire n'est \textbf{pas} soudé à la trompe : le pavillon frangé « aspire » l'ovule libéré dans la cavité péritonéale.
\end{itemize}
\end{attention}

\courssub{La puberté : la maturation sexuelle}

\textbf{Situation concrète.} Entre la classe de CM2 et la classe de Seconde, les élèves d'un même âge se transforment de façon spectaculaire : la voix des garçons mue, des poils apparaissent, les filles voient leurs premières règles survenir, parfois en pleine classe, source de gêne quand personne ne les y a préparées. Toutes ces transformations ont une cause unique : la mise en route, par le cerveau, de la fabrication des \cle{hormones sexuelles}.

\begin{definition}[Puberté]
La \cle{puberté} est la période de transition entre l'enfance et l'âge adulte, au cours de laquelle les gonades deviennent fonctionnelles sous l'action des hormones hypophysaires. Elle se traduit par l'apparition des \cle{caractères sexuels secondaires} et par l'acquisition de la \cle{capacité de se reproduire}. Elle survient en moyenne entre \textbf{10 et 14 ans chez la fille} et entre \textbf{12 et 16 ans chez le garçon}, avec de grandes variations individuelles normales.
\end{definition}

\textbf{Le déclencheur : une commande cérébrale (axe hypothalamo-hypophysaire).} Au niveau de la base du cerveau, deux structures jouent un rôle de « chef d'orchestre » :
\begin{itemize}
\item l'\cle{hypothalamus} sécrète la \cle{GnRH} (\emph{Gonadotropin-Releasing Hormone}), qui stimule l'hypophyse antérieure ;
\item l'\cle{hypophyse} (glande située sous l'hypothalamus) libère alors deux hormones dites \cle{gonadotrophines} : la \cle{FSH} (\emph{hormone folliculo-stimulante}) et la \cle{LH} (\emph{hormone lutéinisante}).
\end{itemize}
Ces deux hormones agissent sur les gonades : chez le garçon, la FSH stimule les cellules de Sertoli (spermatogenèse) et la LH stimule les cellules de Leydig (sécrétion de \cle{testostérone}) ; chez la fille, la FSH stimule la croissance folliculaire (sécrétion d'\cle{œstrogènes}) et la LH déclenche l'ovulation et la formation du corps jaune (sécrétion de \cle{progestérone}).

\begin{center}
\begin{tabularx}{\linewidth}{|l|X|X|}
\hline
\rowcolor{popblueL} \textbf{Hormone} & \textbf{Origine} & \textbf{Rôle principal} \\
\hline
GnRH & Hypothalamus & Déclenche la sécrétion de FSH et LH par l'hypophyse \\
\hline
FSH & Hypophyse antérieure & Fille : croissance des follicules ovariens (sécrétion d'œstrogènes) ; Garçon : soutien de la spermatogenèse (via cellules de Sertoli) \\
\hline
LH & Hypophyse antérieure & Fille : déclenche l'ovulation, luteinisation (corps jaune $\to$ progestérone) ; Garçon : stimulation des cellules de Leydig $\to$ sécrétion de testostérone \\
\hline
Testostérone & Testicules (cellules de Leydig) & Caractères sexuels secondaires masculins, entretien de la spermatogenèse, libido \\
\hline
Œstrogènes & Ovaires (follicules) & Caractères sexuels secondaires féminins, reconstruction de l'endomètre (phase folliculaire) \\
\hline
Progestérone & Ovaires (corps jaune) & Prépare et maintient l'endomètre pour une éventuelle nidation (phase lutéale), effet thermogène \\
\hline
\end{tabularx}
\end{center}

\textbf{Caractères sexuels primaires et secondaires.} Les \cle{caractères sexuels primaires} sont les organes génitaux eux-mêmes, présents dès la naissance (testicules, pénis, canaux déférents ; ovaires, utérus, trompes, vagin). Les \cle{caractères sexuels secondaires} apparaissent à la puberté sous l'action des hormones sexuelles :
\begin{itemize}
\item chez le garçon (sous l'action de la testostérone) : mue de la voix (allongement des cordes vocales), pousse de la barbe et des poils (visage, thorax, pubis), élargissement des épaules, développement de la masse musculaire et squelettique ;
\item chez la fille (sous l'action des œstrogènes) : développement des seins (thélarche), élargissement du bassin, apparition des poils pubiens et axillaires (pubarche), survenue des \cle{premières règles} ou \cle{ménarche}.
\end{itemize}

\begin{savaistu}[Et au Cameroun ?]
Au Cameroun, l'âge moyen des premières règles (ménarche) se situe entre \textbf{12 et 13 ans}, avec des variations individuelles normales entre 10 et 15 ans selon la nutrition, l'ethnie et l'environnement. Cela signifie qu'une fille peut devenir capable de concevoir un enfant dès l'école primaire ou le collège : d'où l'importance capitale d'une \textbf{information précoce}, donnée sans tabou par la famille, l'école (cours de SVT, clubs santé) et les centres de santé (CSPC, centres « Ado »), bien avant que les situations à risque ne se présentent.
\end{savaistu}

\courssub{La production des gamètes : spermatogenèse et ovogenèse}

\begin{definition}[Gamétogenèse]
La \cle{gamétogenèse} est l'ensemble des divisions cellulaires (mitoses de multiplication, \emph{méiose} de réduction, différenciation) qui produisent les gamètes haploïdes ($n$) à partir des cellules souches diploïdes ($2n$) des gonades. Chez l'homme, c'est la \cle{spermatogenèse} ; chez la femme, l'\cle{ovogenèse}.
\end{definition}

\textbf{Chez l'homme : une production continue et massive.} À partir de la puberté et jusqu'à la fin de la vie, les tubules séminifères des testicules produisent en continu des spermatozoïdes, à raison de \textbf{plusieurs dizaines de millions par jour} (environ 1000 par battement de cœur). Chaque spermatozoïde est une cellule mobile, hautement différenciée, dotée d'un flagelle, contenant un noyau haploïde ($n = 23$ chromosomes). Un éjaculat (volume 2 à 5 mL) contient typiquement \textbf{40 à 300 millions de spermatozoïdes}.

\textbf{Chez la femme : un stock limité et une libération cyclique.} Durant la vie fœtale, les ovaires d'une fillette constituent une réserve de follicules (environ 1 à 2 millions à la naissance, 400 000 à la puberté). Au cours de la vie reproductive, seulement \textbf{400 à 500} arriveront à maturité et ovuleront. À chaque cycle, \textbf{un seul ovule} (cellule volumineuse, immobile, riche en réserves nutritives, haploïde $n = 23$ chromosomes) est en général libéré : c'est l'\cle{ovulation}. Les autres follicules en croissance subissent une atrésie (dégénérescence).

\begin{attention}[Différence essentielle -- Source d'erreurs fréquentes]
\begin{itemize}
\item \textbf{Continuité :} L'homme produit des gamètes en permanence (spermatogenèse continue) ; la femme libère un ovule par cycle (ovogenèse cyclique).
\item \textbf{Durée de vie :} L'ovule ne reste fécondable que \textbf{12 à 24 heures} après l'ovulation. Les spermatozoïdes survivent \textbf{3 à 5 jours} (parfois jusqu'à 7 jours) dans les voies génitales féminines grâce au mucus cervical.
\item \textbf{Conséquence :} C'est ce qui explique l'existence d'une « fenêtre de fertilité » de plusieurs jours \textbf{avant} l'ovulation. Un rapport sexuel 3 à 5 jours avant l'ovulation peut féconder l'ovule du jour J.
\end{itemize}
\end{attention}

\courssub{Le cycle menstruel (ou cycle ovarien-utérin)}

\textbf{Observation.} Une femme non enceinte et en bonne santé a des règles à intervalles réguliers. Pourquoi ce rythme ? Parce que l'ovaire et l'utérus se préparent chaque mois, de façon coordonnée par les hormones, à une éventuelle grossesse.

\begin{definition}[Cycle menstruel]
Le \cle{cycle menstruel} est l'ensemble des transformations cycliques et synchrones des ovaires (cycle ovarien) et de l'utérus (cycle utérin), commandées par les hormones (FSH, LH, œstrogènes, progestérone). Sa durée moyenne est de \textbf{28 jours} (variabilité physiologique de 21 à 35 jours). Par convention, le \textbf{premier jour des règles marque le début du cycle (Jour 1)}.
\end{definition}

Le cycle se déroule en plusieurs phases synchrones :

\begin{enumerate}
\item \textbf{Phase des règles / Menstrues (Jours 1 à 5 environ)} : En l'absence de fécondation au cycle précédent, le corps jaune a régressé, les taux d'œstrogènes et de progestérone chutent brutalement. La muqueuse utérine (endomètre fonctionnel), privée de soutien hormonal, se nécrose, se détache et s'évacue par le vagin avec du sang et du mucus : ce sont les \cle{règles}. Simultanément, sous l'effet de la FSH (légèrement élevée car libérée de l'inhibition), une cohorte de follicules commence à croître dans un ovaire.
\item \textbf{Phase folliculaire / Pré-ovulatoire (Jours 1 à 14)} : Parmi la cohorte, un follicule devient \textbf{dominant} ; il sécrète de plus en plus d'\cle{œstrogènes} (estradiol). Les œstrogènes provoquent la \textbf{prolifération et la reconstruction de l'endomètre} (phase proliférative utérine). Le taux élevé d'œstrogènes (seuil critique maintenu ~36-48h) exerce un \textbf{rétrocontrôle positif} sur l'hypophyse : il déclenche un pic brutal et massif de \cle{LH} (et un pic moindre de FSH).
\item \textbf{Ovulation (Vers le 14\textsuperscript{e} jour pour un cycle de 28 jours)} : Le pic de LH provoque la rupture de la paroi du follicule dominant (ovulation) et la libération de l'ovocyte II (ovule), aspiré par le pavillon de la trompe. L'ovule est viable 12 à 24 h.
\item \textbf{Phase lutéale / Post-ovulatoire (Jours 14 à 28)} : Le follicule rompu se transforme en \cle{corps jaune} (corpus luteum), structure endocrine transitoire qui sécrète de grandes quantités de \cle{progestérone} (et des œstrogènes). Sous l'action de la progestérone, l'endomètre s'épaissit, se vascularise et se glandulaire : c'est la \textbf{phase sécrétoire}, prête à accueillir un embryon (fenêtre d'implantation J20-J24). \textbf{S'il n'y a pas fécondation}, le corps jaune régresse vers le 24\textsuperscript{e}-26\textsuperscript{e} jour (luteolyse), les taux d'hormones chutent $\to$ nouvelles règles $\to$ nouveau cycle.
\end{enumerate}

\begin{popfigure}
\begin{center}
\begin{tikzpicture}
\begin{axis}[
width=\linewidth, height=7cm,
xlabel={Jours du cycle}, ylabel={Taux hormonal relatif (u.a.)},
xmin=0, xmax=28, ymin=0, ymax=10.5,
xtick={0,5,10,14,20,28}, ytick=\empty,
legend style={at={(0.02,0.98)}, anchor=north west, font=\scriptsize, draw=popline, fill=popcream},
axis lines=left, grid=major, grid style={popline!30},
clip=false
]
% FSH : petit pic J1, bas, pic ovulatoire J14
\addplot[popblue, thick, domain=0:28, samples=300] { 
    1.5 + 1.5*exp(-((x-2)^2)/3) + 2.5*exp(-((x-13.8)^2)/4) 
};
\addlegendentry{FSH}
% LH : bas, pic énorme J14
\addplot[poppurple, very thick, domain=0:28, samples=300] { 
    1.0 + 8.5*exp(-((x-13.8)^2)/1.5) 
};
\addlegendentry{LH}
% Œstrogènes : montée J5-J13, pic J13, baisse, second pic J21
\addplot[poporange, thick, domain=0:28, samples=300] { 
    1.0 + 5.5*exp(-((x-12.5)^2)/12) + 2.5*exp(-((x-21)^2)/10) 
};
\addlegendentry{Œstrogènes}
% Progestérone : bas, montée J15-J22, pic J21, chute
\addplot[popgreen, thick, domain=0:28, samples=300] { 
    0.5 + 6.5*exp(-((x-21)^2)/14) 
};
\addlegendentry{Progestérone}
% Ligne ovulation
\draw[dashed, popdark, thick] (axis cs:13.8,0) -- (axis cs:13.8,10) node[pos=0.95, anchor=south east, font=\scriptsize\itshape, text=popdark]{Pic LH $\to$ Ovulation};
% Barre règles
\draw[popred, thick] (axis cs:1,9.8) -- (axis cs:5,9.8) node[midway, above, font=\scriptsize, text=popred]{Règles};
% Barre phase lutéale
\draw[poppurple, thick] (axis cs:14,10.2) -- (axis cs:28,10.2) node[midway, above, font=\scriptsize, text=poppurple]{Phase lutéale (Corps jaune)};
\end{axis}
\end{tikzpicture}
\end{center}
\legende{Évolution schématique des taux sanguins de FSH, LH, œstrogènes et progestérone au cours d'un cycle standard de 28 jours. Le pic de LH (vers J14) déclenche l'ovulation ; la progestérone domine la phase lutéale grâce au corps jaune. La chute des hormones (J28) entraîne les règles.}
\end{popfigure}

\begin{propriete}[Ovulation et fenêtre de fertilité (connaissance admise, base du raisonnement)]
Dans un cycle de 28 jours, l'\cle{ovulation} a lieu en moyenne au \textbf{14\textsuperscript{e} jour} (J14). L'ovule restant fécondable 12 à 24 h et les spermatozoïdes survivant 3 à 5 jours dans le mucus cervical, la \cle{fécondation} est possible pendant une \textbf{fenêtre d'environ 5 à 6 jours} : les 4-5 jours \textbf{avant} l'ovulation et le jour qui la suit. Comme la durée des cycles varie (phase folliculaire variable) et que l'ovulation n'est pas toujours exactement au J14, cette fenêtre est \textbf{imprévisible avec certitude} chez une femme donnée sans méthodes de repérage (température, tests d'ovulation, observation du mucus).
\end{propriete}

\begin{definition}[Période fertile]
La \cle{période fertile} est l'intervalle du cycle pendant lequel un rapport sexuel non protégé peut aboutir à une grossesse. Elle correspond à la fenêtre entourant l'ovulation. Une fille est fertile \textbf{dès ses premières règles} (la première ovulation précède en réalité les toutes premières règles), et une grossesse reste possible même lors de rapports pendant les règles, car les cycles des adolescentes sont souvent irréguliers et courts.
\end{definition}

\begin{attention}[Erreur fréquente et dangereuse -- « Mythes » à déconstruire]
Croire qu'une fille \textbf{ne peut pas tomber enceinte dès ses premières règles}, ou qu'un rapport \textbf{pendant les règles} est sans risque, ou encore que le \textbf{retrait (coït interrompu)} est une méthode fiable, sont des erreurs aux conséquences graves. En réalité :
\begin{enumerate}
\item La première ovulation survient \emph{avant} les premières règles (la ménarche est la conséquence de la première chute hormonale post-ovulatoire sans fécondation).
\item Chez l'adolescente, les cycles sont souvent anovulatoires ou irréguliers (phase folliculaire imprévisible) : l'ovulation peut survenir très tôt (J8-J10) ou très tard.
\item Les spermatozoïdes survivent plusieurs jours dans le mucus cervical.
\item Le liquide pré-éjaculatoire peut contenir des spermatozoïdes.
\end{enumerate}
\textbf{Conclusion : Tout rapport sexuel non protégé, à n'importe quel moment du cycle, peut entraîner une grossesse.}
\end{attention}

\courssub{La fécondation et la nidation}

\textbf{La fécondation.} Lors d'un rapport sexuel avec éjaculation intravaginale, le sperme est déposé au fond du vagin (fornix postérieur). Les spermatozoïdes, mobiles grâce à leur flagelle, remontent alors activement : ils traversent le col de l'utérus (filtrés par le mucus cervical), l'utérus et pénètrent dans les trompes. Sur les centaines de millions de départ, \textbf{seules quelques centaines (environ 200-300)} atteignent l'ampoule de la trompe où attend l'ovule. La \cle{fécondation} est la fusion de la membrane plasmique d'\textbf{un seul} spermatozoïde ($n = 23$ chromosomes, porteur soit d'un chromosome X soit Y) avec celle de l'ovule ($n = 23$ chromosomes, porteur obligatoirement d'un chromosome X), dans le \textbf{tiers externe de la trompe (ampoule)}. Il se forme une cellule unique, la \cle{cellule-œuf} (ou zygote), diploïde ($2n = 46$ chromosomes) : le patrimoine génétique du nouvel individu est fixé, y compris son sexe chromosomique (XX = fille, XY = garçon). Le blocage de la polyspermie (entrée de plusieurs spermatozoïdes) est immédiat (réaction corticale, modification de la zone pellucide).

\textbf{La nidation (implantation).} La cellule-œuf se divise aussitôt par mitoses successives (clivage) en 2, 4, 8 cellules (blastomères)... tout en migrant lentement dans la trompe vers l'utérus grâce aux cils de l'épithélium tubaire et aux contractions musculaires (trajet de 3 à 4 jours). Arrivé dans la cavité utérine au stade de \textbf{blastocyste} (petite sphère creuse d'environ 100-150 cellules, J5-J6 post-fécondation), l'embryon « éclot » de sa zone pellucide et se fixe dans l'endomètre épaissi et vascularisé (phase sécrétoire) vers le \textbf{7\textsuperscript{e} jour après la fécondation} : c'est la \cle{nidation} (ou implantation). L'embryon (plus précisément le trophoblaste, future partie fœtale du placenta) sécrète alors une hormone spécifique, la \cle{hCG} (\emph{gonadotrophine chorionique humaine}), qui « sauve » le corps jaune maternel : la progestérone continue d'être produite, l'endomètre est maintenu, et \textbf{les règles ne surviennent pas}. L'absence de règles (aménorrhée) est ainsi le premier signe clinique d'une grossesse. C'est aussi le principe des tests de grossesse urinaires (bandelettes) vendus en pharmacie : ils détectent la présence de l'hormone hCG dans les urines (seuil ~25-50 UI/L, positif dès le 1er jour de retard de règles).

\begin{exoresolu}[Localiser la fécondation et raisonner sur la fertilité]
\textbf{Énoncé.} Une jeune femme a des cycles réguliers de 28 jours. Ses dernières règles ont débuté le 3 mars. 
\begin{enumerate}
\item À quelle date approximative aura lieu son ovulation ?
\item Elle a eu un rapport non protégé le 14 mars : une grossesse est-elle possible ? Justifier en utilisant la notion de fenêtre de fertilité.
\item Expliquer pourquoi, si une grossesse démarre, les règles prévues le 31 mars ne surviendront pas.
\end{enumerate}

\tcblower
\textbf{Solution détaillée.}
\begin{enumerate}
\item L'ovulation a lieu en moyenne au 14\textsuperscript{e} jour du cycle. Le 1\textsuperscript{er} jour du cycle est le premier jour des règles (3 mars). Le 14\textsuperscript{e} jour correspond donc au \textbf{16 mars} (3 mars + 13 jours).
\item Le 14 mars correspond au 12\textsuperscript{e} jour du cycle (3 mars = J1 $\to$ 14 mars = J12), soit environ \textbf{2 jours avant l'ovulation} prévue (J14). Or, les spermatozoïdes survivent 3 à 5 jours (parfois plus) dans les voies génitales féminines, notamment dans le mucus cervical favorable de la phase pré-ovulatoire. Ils seront donc encore vivants et fécondants lors de la libération de l'ovule le 16 mars. \textbf{Une grossesse est donc tout à fait possible}, car le rapport a eu lieu en plein dans la fenêtre de fertilité (J10 à J15 environ).
\item Si la fécondation a lieu (vers le 16 mars), l'embryon migre vers l'utérus et réalise la nidation vers le 7\textsuperscript{e} jour post-fécondation (autour du 23 mars). Dès la nidation, le trophoblaste sécrète l'hormone \cle{hCG}. Cette hormone maintient le corps jaune ovarien en vie et fonctionnel : celui-ci continue de produire de la progestérone (et des œstrogènes) au-delà de sa durée de vie habituelle (14 jours). L'endomètre, sous l'action continue de la progestérone, est conservé au lieu de se nécroser et de se détacher : \textbf{les règles ne surviennent pas} à la date prévue (31 mars = J29). Leur absence (aménorrhée) constitue le premier signe de grossesse et doit conduire à consulter un centre de santé pour confirmation (test hCG, échographie).
\end{enumerate}
\end{exoresolu}

\begin{aretenir}[L'essentiel de la section -- À maîtriser pour le Probatoire]
\begin{itemize}
\item Les \cle{gonades} (testicules, ovaires) ont une double fonction : production des gamètes (spermatogenèse continue / ovogenèse cyclique) et sécrétion des hormones sexuelles (testostérone / œstrogènes-progestérone).
\item La \cle{fécondation} a lieu dans l'\textbf{ampoule de la trompe} (tiers externe) ; la grossesse se déroule dans l'\textbf{utérus}.
\item La \cle{puberté} (10-14 ans fille, 12-16 ans garçon) est déclenchée par l'axe hypothalamo-hypophysaire (GnRH $\to$ FSH/LH) et marque l'acquisition de la capacité de reproduction.
\item Le \cle{cycle menstruel} (moyenne 28 jours) est piloté par FSH, LH, œstrogènes, progestérone. L'ovulation (J14) est déclenchée par le pic de LH.
\item La \cle{fenêtre de fertilité} s'étend sur ~5-6 jours autour de l'ovulation (survie spermatozoïdes 3-5 j + survie ovule 12-24 h). \textbf{Une fille est fertile dès ses premières règles.}
\item La \cle{nidation} (J7 post-fécondation) entraîne la sécrétion de \cle{hCG}, le maintien du corps jaune et de la progestérone, donc l'arrêt des règles (aménorrhée).
\end{itemize}
\end{aretenir}

\begin{exercice}[Pour vérifier tes connaissances -- Type Probatoire]
\begin{enumerate}
\item Définis la puberté et cite deux caractères sexuels secondaires propres à chaque sexe en précisant l'hormone responsable.
\item Schématise, sous forme de frise chronologique annotée, les quatre phases du cycle menstruel (J1 à J28) en précisant pour chacune : l'événement ovarien majeur, l'événement utérin majeur et l'hormone(s) dominante(s).
\item Une adolescente de 14 ans affirme : « Je ne risque pas de tomber enceinte, j'ai eu mes premières règles il y a seulement deux mois et mes cycles ne sont pas réguliers. » Corrige cette affirmation en t'appuyant sur \textbf{trois} arguments scientifiques précis.
\item \textbf{Calcul.} Un couple a un rapport non protégé le 10 du mois. La femme a des cycles réguliers de 30 jours. Ses dernières règles datent du 1er du mois. L'ovulation a-t-elle lieu vers le 15\textsuperscript{e} jour ? Une grossesse est-elle possible ? Justifier par un calcul de dates.
\end{enumerate}
\end{exercice}

\coursec{Grossesses précoces et non désirées}

Dans la section précédente, nous avons étudié la physiologie de la reproduction humaine et les transformations de la puberté : nous savons désormais que dès les premières règles (ménarche), une adolescente possède des ovaires fonctionnels capables de libérer des ovocytes, et qu'un rapport sexuel non protégé pendant la période fertile peut aboutir à une fécondation. La biologie, donc, rend la grossesse possible très tôt. Mais le corps d'une adolescente est-il prêt à mener une grossesse à terme ? La société, la famille, l'école sont-elles prêtes ? C'est tout l'objet de cette section : comprendre, avec des données et une démarche scientifique, pourquoi les grossesses précoces constituent un problème majeur de santé publique au Cameroun, et comment agir de façon responsable.

\courssub{Définition et situation concrète}

\textbf{Une situation d'introduction.} Dans un quartier de Douala, Amina, 16 ans, élève en classe de Seconde, s'absente de plus en plus souvent des cours. Elle se sent fatiguée, a des nausées le matin, et ses règles ont cessé depuis trois mois. Le test de grossesse réalisé au centre de santé intégré est positif. Amina n'avait ni prévu ni souhaité cette grossesse ; elle ignorait qu'elle se trouvait en période fertile le jour du rapport non protégé. Sa scolarité, sa santé et son avenir sont désormais profondément bouleversés.

Cette situation, loin d'être exceptionnelle, illustre la notion centrale de cette section.

\begin{definition}[Grossesse précoce et non désirée]
On appelle \cle{grossesse précoce} (ou grossesse chez l'adolescente) toute grossesse survenant chez une fille âgée de moins de 18 à 19 ans, c'est-à-dire avant la fin complète de la croissance et de la maturation du bassin et de l'appareil génital. Elle est dite \cle{non désirée} lorsqu'elle n'était ni voulue ni planifiée par la jeune fille : elle résulte le plus souvent d'un rapport sexuel non protégé, d'une méconnaissance du cycle menstruel ou d'un rapport imposé.
\end{definition}

\begin{exemplebox}[Une réalité statistique camerounaise]
Au Cameroun, selon les données indicatives des Enquêtes Démographiques et de Santé (EDS), environ \textbf{1 adolescente sur 4} (environ 25 \%) a déjà commencé une vie de mère avant l'âge de 20 ans : elle est soit déjà mère, soit enceinte de son premier enfant. Ce taux est nettement plus élevé en milieu rural et dans les régions où la scolarisation des filles est faible (Extrême-Nord, Nord, Adamaoua) qu'à Yaoundé ou Douala. Cette proportion place les grossesses précoces parmi les premiers problèmes de santé reproductive des adolescent(e)s dans notre pays.
\end{exemplebox}

\courssub{Les causes des grossesses précoces}

Pourquoi une grossesse survient-elle si tôt, alors que la jeune fille ne la désirait pas ? La démarche scientifique consiste à ne pas se contenter d'une explication unique (« c'est de sa faute »), mais à rechercher l'ensemble des facteurs, souvent combinés.

\begin{itemize}
\item \textbf{Les rapports sexuels précoces et non protégés.} C'est la cause biologique directe : tout rapport sexuel sans contraception pendant la période fertile (quelques jours avant et le jour de l'ovulation, vers le 14\textsuperscript{ème} jour d'un cycle de 28 jours) expose à la fécondation. L'absence d'utilisation du préservatif ou de toute autre méthode contraceptive rend ce risque maximal.
\item \textbf{L'ignorance de la période fertile et de la physiologie de la reproduction.} Beaucoup d'adolescentes croient à tort qu'on ne peut pas tomber enceinte « la première fois », ou juste après les règles. Or un ovocyte peut être libéré à des dates variables, et les spermatozoïdes survivent 3 à 5 jours dans les voies génitales féminines : la fécondation est possible presque à tout moment du cycle.
\item \textbf{La pauvreté.} Certaines adolescentes échangent des relations sexuelles contre de l'argent, des cadeaux ou le paiement de leurs fournitures scolaires (phénomène parfois appelé « rapport transactionnel »). La précarité économique les prive du pouvoir de négocier l'usage du préservatif.
\item \textbf{La pression sociale et l'influence du groupe.} Le désir de « faire comme les autres », la pression d'un partenaire plus âgé, ou l'idée fausse qu'un enfant prouve la fécondité, poussent certaines jeunes filles à des rapports non protégés.
\item \textbf{Les violences et les abus sexuels.} Une partie des grossesses précoces résulte de viols ou d'abus, parfois au sein même de l'entourage. La victime n'a alors aucun choix : c'est un crime, et non une faute de la jeune fille.
\item \textbf{Les mariages précoces.} Dans certaines localités, des filles sont mariées dès 13 ou 14 ans. La grossesse qui suit est « socialement acceptée » mais reste médicalement précoce, avec les mêmes risques sanitaires.
\end{itemize}

\begin{attention}[Prévention, pas jugement moral]
Une erreur fréquente, en classe comme dans la société, consiste à confondre \textbf{prévention} et \textbf{jugement moral}. L'approche du scientifique et de l'éducateur n'est pas de condamner la jeune fille enceinte, mais d'\emph{analyser les causes} pour \emph{prévenir} le phénomène et \emph{soutenir} la personne concernée. Stigmatiser une adolescente enceinte ne résout rien ; comprendre les mécanismes biologiques et sociaux, si. Dans tes rédactions au Probatoire, adopte toujours un ton scientifique et humain : on explique, on justifie, on propose — on ne juge pas.
\end{attention}

\courssub{Les conséquences des grossesses précoces}

\courssub{Conséquences sur la santé de la mère et de l'enfant}

Le corps d'une adolescente n'a pas terminé sa croissance : le bassin osseux, en particulier, n'atteint ses dimensions définitives que vers 18-20 ans. Une grossesse avant cet âge expose donc à des complications graves, d'autant plus sévères que l'âge est jeune (risque maximal avant 15 ans).

\begin{itemize}
\item \textbf{Des accouchements difficiles (dystocie).} Le bassin trop étroit peut empêcher le passage du fœtus : le travail se prolonge (dystocie mécanique), nécessitant souvent une césarienne en urgence, qui n'est pas toujours accessible en zone rurale.
\item \textbf{Les fistules obstétricales.} Lors d'un travail prolongé non pris en charge, la tête du fœtus comprime durablement les tissus mous du petit bassin. Ceux-ci se nécrosent par ischémie, créant une communication anormale entre le vagin et la vessie (fistule vésico-vaginale) ou le rectum (fistule recto-vaginale). La jeune mère souffre alors de fuites permanentes d'urines ou de selles, source de souffrance physique, psychologique et d'exclusion sociale. C'est une complication fréquente des accouchements précoces non assistés, notamment dans le septentrion camerounais.
\item \textbf{Une mortalité maternelle élevée.} Hémorragies de la délivrance, infections (septicémie puerpérale), éclampsie (crises convulsives liées à l'hypertension artérielle gravidique) : les complications de la grossesse et de l'accouchement tuent davantage les adolescentes (surtout < 15 ans) que les femmes de 20 à 30 ans.
\item \textbf{Une mortalité infantile élevée.} Les enfants nés de mères adolescentes présentent plus souvent une prématurité, une hypotrophie (faible poids de naissance, inférieur à 2,5 kg) et un risque accru de décès dans la première année de vie.
\end{itemize}

\begin{savaistu}[Une cause majeure de décès dans le monde]
Selon l'Organisation mondiale de la santé (OMS), les complications de la grossesse et de l'accouchement figurent parmi les \textbf{premières causes de décès des filles de 15 à 19 ans} dans le monde, en particulier dans les pays en développement. Chaque année, plusieurs millions d'adolescentes accouchent avant 19 ans, et des dizaines de milliers en meurent ou en gardent des séquelles à vie, comme les fistules obstétricales. Au Cameroun, des programmes de réparation chirurgicale des fistules existent dans certains hôpitaux de référence (ex. hôpital gynéco-obstétrique de Yaoundé, hôpital de district de Mokolo) : preuve que le problème est réel, mais aussi qu'il peut être pris en charge.
\end{savaistu}

\courssub{Conséquences scolaires}

La grossesse précoce interrompt presque toujours le parcours scolaire :
\begin{itemize}
\item \textbf{L'abandon des études} : la fatigue de la grossesse, l'accouchement puis la charge de l'enfant rendent la poursuite des études très difficile ; la plupart des jeunes mères ne retournent jamais en classe.
\item \textbf{L'exclusion scolaire} : par honte, par stigmatisation ou par pression de l'entourage, l'adolescente enceinte quitte l'établissement, alors que le partenaire, lui, poursuit généralement sa scolarité — une injustice qui aggrave les inégalités entre filles et garçons.
\item \textbf{La perte des perspectives d'avenir} : sans diplôme, l'accès à un emploi qualifié et rémunéré devient quasi impossible, enfermant la jeune mère dans la dépendance économique.
\end{itemize}

\courssub{Conséquences sociales et économiques}

\begin{itemize}
\item \textbf{Le rejet familial} : certaines familles chassent la jeune fille enceinte du domicile, la privant de tout soutien matériel et affectif au moment où elle en a le plus besoin.
\item \textbf{La stigmatisation} : moqueries, rumeurs, exclusion du groupe des pairs ; la jeune mère porte seule le poids du regard social.
\item \textbf{La précarité économique} : sans formation ni revenu, la jeune mère et son enfant entrent dans un cycle de pauvreté ; l'enfant lui-même grandit dans des conditions difficiles, ce qui perpétue le phénomène d'une génération à l'autre.
\end{itemize}

\begin{popfigure}
\begin{center}
\begin{tikzpicture}[
  font=\small,
  cause/.style={rectangle, rounded corners, draw=poporange, fill=poporangeL, text width=3.1cm, align=center, inner sep=3pt},
  central/.style={rectangle, rounded corners, draw=popdark, fill=popdark, text=white, text width=3.4cm, align=center, inner sep=4pt, font=\small\bfseries},
  cons/.style={rectangle, rounded corners, draw=popblue, fill=popblueL, text width=3.1cm, align=center, inner sep=3pt},
  lien/.style={-{Stealth[length=2.2mm]}, thick, popmuted}
]
\node[central] (G) at (0,0) {Grossesse précoce et non désirée};
% Causes (à gauche)
\node[cause] (c1) at (-6.2,3.2) {Rapports sexuels précoces non protégés};
\node[cause] (c2) at (-6.2,1.6) {Ignorance de la période fertile};
\node[cause] (c3) at (-6.2,0) {Pauvreté et rapports transactionnels};
\node[cause] (c4) at (-6.2,-1.6) {Violences, abus sexuels};
\node[cause] (c5) at (-6.2,-3.2) {Mariages précoces, pression sociale};
\draw[lien] (c1.east) -- (G.west);
\draw[lien] (c2.east) -- (G.west);
\draw[lien] (c3.east) -- (G.west);
\draw[lien] (c4.east) -- (G.west);
\draw[lien] (c5.east) -- (G.west);
% Conséquences (à droite)
\node[cons] (s1) at (6.2,3.2) {\textbf{Santé} : accouchement difficile, fistules, mortalité maternelle et infantile};
\node[cons] (s2) at (6.2,0.8) {\textbf{Scolarité} : abandon des études, exclusion};
\node[cons] (s3) at (6.2,-1.6) {\textbf{Social} : rejet familial, stigmatisation};
\node[cons] (s4) at (6.2,-3.6) {\textbf{Économie} : précarité, cycle de pauvreté};
\draw[lien] (G.east) -- (s1.west);
\draw[lien] (G.east) -- (s2.west);
\draw[lien] (G.east) -- (s3.west);
\draw[lien] (G.east) -- (s4.west);
\end{tikzpicture}
\end{center}
\legende{Arbre des causes (à gauche) et des conséquences (à droite) des grossesses précoces et non désirées. Les causes sont multiples et souvent combinées ; les conséquences touchent la santé, la scolarité, la vie sociale et la situation économique de l'adolescente.}
\end{popfigure}

\courssub{Une pratique aggravante : l'avortement clandestin}

Face à une grossesse non désirée, certaines adolescentes, par peur du rejet familial ou de l'exclusion scolaire, ont recours à l'\cle{avortement clandestin} : interruption de grossesse pratiquée en cachette, hors de tout cadre médical, par des personnes non qualifiées, avec des moyens dangereux (tiges, produits caustiques, médicaments détournés à doses incontrôlées, décoctions non dosées).

\begin{attention}[L'avortement clandestin tue]
L'avortement clandestin est une pratique \textbf{illégale} (sauf exceptions très restrictives prévues par la loi camerounaise : viol avéré ou danger pour la vie de la mère, sous contrôle médical strict) et \textbf{mortellement dangereuse}. Ses complications sont :
\begin{itemize}
\item les \textbf{hémorragies} massives, pouvant entraîner le décès en quelques heures ;
\item les \textbf{infections graves} (septicémie) par instruments non stérilisés ;
\item les \textbf{perforations de l'utérus} et les lésions des organes voisins (intestin, vessie) ;
\item la \textbf{stérilité définitive} par destruction de la muqueuse utérine (synéchies) ou obstruction des trompes.
\end{itemize}
Les complications des avortements clandestins constituent une cause importante de décès et de stérilité chez les adolescentes. La seule conduite sûre face à une grossesse non désirée est de \textbf{consulter rapidement un personnel de santé} et d'en parler à un adulte de confiance, jamais de recourir à des pratiques clandestines.
\end{attention}

\courssub{Données : le niveau d'instruction, un facteur protecteur}

Appliquons la démarche d'investigation. \textbf{Observation} : dans les statistiques de santé, toutes les adolescentes ne sont pas également touchées par les grossesses précoces. \textbf{Hypothèse} : le niveau d'instruction jouerait un rôle protecteur. \textbf{Données} : examinons les taux indicatifs de maternité précoce selon le niveau scolaire.

\begin{popfigure}
\begin{center}
\begin{tikzpicture}
\begin{axis}[
  ybar,
  width=0.85\linewidth,
  height=6.5cm,
  bar width=0.9cm,
  ymin=0, ymax=50,
  ylabel={Adolescentes déjà mères ou enceintes (\%)},
  symbolic x coords={Aucun niveau, Primaire, Secondaire, Supérieur},
  xtick=data,
  x tick label style={align=center, font=\small},
  ytick={0,10,20,30,40,50},
  nodes near coords,
  every node near coord/.append style={font=\small\bfseries, color=popdark},
  axis line style={popline},
  tick style={popline},
]
\addplot[fill=popblue, draw=popdark] coordinates {(Aucun niveau,42) (Primaire,32) (Secondaire,14) (Supérieur,5)};
\end{axis}
\end{tikzpicture}
\end{center}
\legende{Taux indicatifs de maternité précoce (adolescentes déjà mères ou enceintes avant 20 ans) selon le niveau d'instruction, d'après des données de type EDS. Plus le niveau d'instruction est élevé, plus le taux de grossesses précoces est faible.}
\end{popfigure}

\textbf{Interprétation} : le taux de maternité précoce diminue fortement lorsque le niveau d'instruction augmente : environ 42 \% chez les jeunes filles sans instruction contre environ 5 \% chez celles ayant atteint le supérieur. \textbf{Conclusion} : l'hypothèse est validée — la scolarisation des filles est un \textbf{facteur protecteur} majeur contre les grossesses précoces, car elle apporte des connaissances sur la santé reproductive, retarde l'âge au mariage et offre des perspectives économiques. Maintenir les filles à l'école est donc une stratégie de santé publique.

\courssub{Méthode d'analyse d'un cas concret}

\begin{methode}[Analyser un cas de grossesse précoce]
Face à un document relatant le cas d'une adolescente enceinte (exercice classique au Probatoire), procède en quatre étapes :
\begin{enumerate}
\item \textbf{Décrire les faits} : relever l'âge, la situation scolaire et familiale, les circonstances de la grossesse, sans jugement.
\item \textbf{Identifier les causes} : distinguer la cause biologique directe (rapport non protégé en période fertile) des causes favorisantes (ignorance, pauvreté, pression sociale, abus, mariage précoce).
\item \textbf{Dégager les conséquences} en les classant : sanitaires (risques pour la mère et l'enfant), scolaires (abandon, exclusion), sociales (rejet, stigmatisation) et économiques (précarité).
\item \textbf{Proposer des solutions} : prévention (information, éducation à la santé reproductive, contraception), conduite à tenir (dialogue avec les parents, consultation médicale, soutien à la jeune mère, poursuite ou reprise de la scolarité).
\end{enumerate}
\end{methode}

\begin{exoresolu}[Étude de cas : la situation de Béatrice]
\textbf{Énoncé.} Béatrice, 15 ans, élève en classe de Troisième dans un lycée de Bafoussam, vit chez sa grand-mère. Son père est décédé et sa mère, vendeuse au marché, peine à payer les fournitures scolaires. Un commerçant du quartier, âgé de 28 ans, lui offre régulièrement de l'argent et des cadeaux. À la suite de rapports sexuels non protégés, Béatrice tombe enceinte. Chassée par sa grand-mère, elle abandonne l'école. Un voisin lui conseille de « faire sortir le ventre » avec une décoction.
\begin{enumerate}
\item Identifie les causes de la grossesse de Béatrice.
\item Dégage les conséquences déjà visibles et celles qu'elle risque.
\item Propose une conduite à tenir adaptée.
\end{enumerate}
\tcblower
\textbf{Solution détaillée.}
\begin{enumerate}
\item \textbf{Causes.} Cause biologique directe : des rapports sexuels non protégés, alors que Béatrice, pubère, ovulait déjà. Causes favorisantes : la \emph{pauvreté} (rapports transactionnels avec un homme plus âgé contre de l'argent), l'\emph{ignorance} de la période fertile et des moyens de prévention, et le \emph{rapport de force} avec un adulte (abus de la vulnérabilité d'une mineure, ce qui relève de la loi).
\item \textbf{Conséquences.} Déjà visibles : abandon des études, rejet familial, précarité. Risques sanitaires : accouchement difficile (bassin immature à 15 ans), fistule obstétricale, mortalité maternelle et infantile accrues ; risque supplémentaire si elle suit le conseil du voisin : l'avortement clandestin peut provoquer hémorragie, infection, stérilité, voire le décès. Risque social : stigmatisation durable.
\item \textbf{Conduite à tenir.} Béatrice doit \emph{refuser toute pratique d'avortement clandestin}, se rendre au centre de santé le plus proche pour un suivi prénatal, en parler à un adulte de confiance (sa mère, un enseignant, une assistante sociale) et bénéficier d'un soutien psychologique et matériel. La famille doit être sensibilisée pour accueillir plutôt que rejeter. Après l'accouchement, une reprise de la scolarité doit être encouragée. En amont, la prévention passe par l'information sur la santé reproductive et la dénonciation des abus sur mineures.
\end{enumerate}
\end{exoresolu}

\courssub{Conduite à tenir : informer, dialoguer, soutenir}

Face au problème des grossesses précoces, chacun peut agir à son niveau :
\begin{itemize}
\item \textbf{S'informer et informer} : connaître la physiologie de la reproduction (période fertile, fécondation) et les moyens de prévention ; partager ces connaissances de manière responsable.
\item \textbf{Dialoguer} : entretenir un climat de confiance avec ses parents, ses enseignants ou les personnels de santé ; un adolescent bien informé et écouté prend de meilleures décisions. Les consultations des centres de santé sont confidentielles.
\item \textbf{Soutenir la jeune mère} : ne jamais rejeter ni stigmatiser une camarade enceinte ou jeune mère ; l'encourager à poursuivre le suivi médical et, après l'accouchement, à reprendre ses études. Le soutien de l'entourage améliore nettement la santé de la mère et de l'enfant.
\item \textbf{Promouvoir la scolarisation des filles} : comme le montrent les données, l'école est le meilleur bouclier contre les grossesses précoces.
\end{itemize}

\begin{aretenir}[L'essentiel de la section]
\begin{itemize}
\item Une \cle{grossesse précoce} survient avant 18-19 ans ; elle est le plus souvent \cle{non désirée}.
\item Ses causes sont multiples : rapports non protégés, ignorance de la période fertile, pauvreté, pression sociale, violences, mariages précoces.
\item Ses conséquences sont graves : risques sanitaires (accouchement difficile, fistules, mortalité maternelle et infantile), abandon scolaire, rejet familial, précarité.
\item L'\cle{avortement clandestin} est illégal et mortellement dangereux : il faut toujours consulter un personnel de santé.
\item La scolarisation des filles est un facteur protecteur majeur.
\item L'attitude responsable : informer, dialoguer, soutenir — jamais juger.
\end{itemize}
\end{aretenir}

\begin{exercice}[Pour t'entraîner]
Dans un village de la région de l'Extrême-Nord, Mariam, mariée à 14 ans, est enceinte à 15 ans. Son accouchement, sans assistance médicale, dure deux jours ; elle survit mais souffre désormais de fuites urinaires permanentes.
\begin{enumerate}
\item Quelle complication présente Mariam ? Explique son mécanisme.
\item Cette grossesse est-elle « socialement acceptée » ? Est-elle pour autant sans risque ? Justifie.
\item Propose deux mesures qui auraient pu prévenir cette situation.
\end{enumerate}
\end{exercice}

\begin{corrige}[Exercice]
\begin{enumerate}
\item Mariam présente une \textbf{fistule obstétricale} (vésico-vaginale). Lors du travail prolongé, la tête du fœtus a comprimé longuement les tissus entre le vagin et la vessie ; privés de sang (ischémie), ces tissus se sont nécrosés, créant une communication anormale responsable des fuites d'urines.
\item Oui, cette grossesse est socialement acceptée car elle s'inscrit dans un mariage. Mais elle reste médicalement \textbf{précoce} : à 15 ans, le bassin n'est pas mature, d'où la dystocie et la fistule. L'acceptation sociale ne supprime pas les risques biologiques.
\item Mesures de prévention : retarder l'âge au mariage et maintenir les filles à l'école ; assurer un suivi prénatal et un accouchement assisté par du personnel qualifié dans une formation sanitaire.
\end{enumerate}
\end{corrige}

Nous venons de voir que les grossesses précoces constituent le premier risque lié aux rapports sexuels non protégés. Mais ce n'est pas le seul : ces mêmes rapports exposent aussi aux \textbf{infections sexuellement transmissibles (IST)}, dont le VIH/Sida, que nous étudierons dans la section suivante.

\coursec{Infections sexuellement transmissibles (IST)}

Dans la section précédente, nous avons vu que les grossesses précoces et non désirées constituent un risque majeur des rapports sexuels non protégés chez les adolescent(e)s. Mais ce n'est pas le seul. Lors d'un rapport sexuel non protégé, des micro-organismes pathogènes peuvent également passer d'un partenaire à l'autre : ce sont les \cle{infections sexuellement transmissibles}, ou IST. Contrairement à une grossesse, une IST peut passer inaperçue pendant des mois, voire des années, tout en causant des dégâts irréversibles. Comprendre ces infections, leurs modes de transmission et les moyens de s'en protéger est donc un enjeu vital de santé publique, en particulier pour les jeunes.

\courssub{Qu'est-ce qu'une IST ?}

\textbf{Une situation concrète pour démarrer.} Dans un centre de santé de Yaoundé, une jeune femme de 19 ans consulte pour des brûlures en urinant et un écoulement anormal. Le médecin diagnostique une gonococcie. À côté d'elle, un jeune homme de 22 ans, en parfaite apparence, apprend lors d'une campagne de dépistage gratuite qu'il est porteur du VIH : il ne ressentait \emph{aucun} symptôme. Deux situations différentes, un même mécanisme : des micro-organismes transmis lors de rapports sexuels non protégés.

\begin{definition}[Infection sexuellement transmissible]
Une \cle{infection sexuellement transmissible} (IST), anciennement appelée maladie sexuellement transmissible (MST), est une infection causée par un agent pathogène (virus, bactérie, parasite ou champignon) transmis \textbf{principalement lors de rapports sexuels non protégés} (vaginaux, anaux ou bucco-génitaux). Certaines IST se transmettent aussi par le sang contaminé ou de la mère à l'enfant.
\end{definition}

\begin{exemplebox}[Exemples d'agents pathogènes responsables d'IST]
\begin{itemize}
\item \textbf{Virus} : VIH (virus de l'immunodéficience humaine), virus de l'hépatite B (VHB), virus de l'herpès (HSV), papillomavirus humains (HPV) ;
\item \textbf{Bactéries} : \emph{Neisseria gonorrhoeae} (gonococcie), \emph{Treponema pallidum} (syphilis), \emph{Chlamydia trachomatis} (chlamydiose) ;
\item \textbf{Parasites} : \emph{Trichomonas vaginalis} (trichomonose).
\end{itemize}
\end{exemplebox}

\begin{attention}[On ne reconnaît pas une personne infectée à son apparence]
Erreur fréquente et dangereuse : croire qu'une personne « propre », « belle » ou « en bonne santé » ne peut pas être infectée. \textbf{De nombreuses IST sont asymptomatiques} (silencieuses), surtout au début : la chlamydiose ne provoque aucun symptôme chez la majorité des femmes infectées, et une personne séropositive au VIH peut paraître en pleine forme pendant des années. Seul un \cle{dépistage} (prise de sang ou prélèvement) permet de savoir si l'on est infecté. Ne jamais se fier aux apparences.
\end{attention}

\courssub{Les principales IST : agents, symptômes, conséquences}

\textbf{Démarche d'investigation.} Comment les médecins ont-ils identifié les agents responsables ? Prenons l'exemple historique du SIDA.
\begin{enumerate}
\item \textbf{Observation} (1981) : des jeunes patients développent des infections graves inhabituelles (pneumonies, cancers rares), signe d'un effondrement des défenses immunitaires.
\item \textbf{Hypothèse} : un agent infectieux transmissible détruit les cellules immunitaires.
\item \textbf{Expérience/données} (1983) : les équipes de Luc Montagnier et Françoise Barré-Sinoussi (Institut Pasteur) isolent un rétrovirus, le VIH, dans les ganglions des malades. On montre ensuite qu'il infecte préférentiellement les lymphocytes T4 (cellules portant le récepteur CD4), chefs d'orchestre de l'immunité.
\item \textbf{Interprétation} : la destruction progressive des lymphocytes T4 explique l'affaiblissement immunitaire et l'apparition des maladies « opportunistes ».
\item \textbf{Conclusion} : le VIH est l'agent du SIDA ; il se transmet par le sang et les sécrétions sexuelles.
\end{enumerate}

\textbf{Symptômes généraux à connaître.} Même si chaque IST a ses particularités, certains signes doivent conduire immédiatement à consulter :
\begin{itemize}
\item \cle{écoulement} anormal par le vagin ou le pénis (purulent ou non) ;
\item \cle{brûlures} ou douleurs en urinant ;
\item \cle{ulcérations}, boutons, vésicules ou chancres sur les organes génitaux ;
\item douleurs du bas-ventre, démangeaisons, fièvre.
\end{itemize}
Mais rappelons-le : \textbf{l'absence de symptômes ne signifie pas l'absence d'infection}.

\begin{center}
\rowcolors{1}{popcream}{white}
\begin{tabularx}{\linewidth}{|l|X|X|X|}
\hline
\rowcolor{popblueL} \textbf{IST} & \textbf{Agent} & \textbf{Symptômes principaux} & \textbf{Conséquences / traitement} \\
\hline
\textbf{VIH/SIDA} & Virus (rétrovirus VIH) & Souvent silencieux des années ; puis infections opportunistes & Destruction des lymphocytes T4, immunodéficience ; pas de guérison, mais trithérapie à vie \\
\hline
\textbf{Gonococcie} & Bactérie (\emph{Neisseria gonorrhoeae}) & Écoulement purulent, brûlures mictionnelles ; parfois silencieuse chez la femme & Stérilité (homme et femme) ; guérit par antibiotiques \\
\hline
\textbf{Syphilis} & Bactérie (\emph{Treponema pallidum}) & Chancre indolore, puis éruptions ; phase latente longue & Atteintes cardiaques et nerveuses, malformations du fœtus ; guérit par antibiotiques (pénicilline) \\
\hline
\textbf{Chlamydiose} & Bactérie (\emph{Chlamydia trachomatis}) & Très souvent asymptomatique ; écoulements, brûlures & Salpingites, stérilité féminine, grossesse extra-utérine ; guérit par antibiotiques \\
\hline
\textbf{Herpès génital} & Virus (HSV) & Vésicules douloureuses récidivantes & Récidives à vie, transmission mère-enfant ; traitement antiviral (atténue, ne guérit pas) \\
\hline
\textbf{Hépatite B} & Virus (VHB) & Souvent silencieuse ; jaunisse, fatigue & Cirrhose, cancer du foie ; \textbf{vaccin efficace} \\
\hline
\textbf{Papillomavirus (HPV)} & Virus (HPV) & Souvent silencieux ; condylomes (verrues génitales) & \textbf{Cancer du col de l'utérus} ; \textbf{vaccin préventif} \\
\hline
\end{tabularx}
\end{center}

\begin{aretenir}[Deux grandes familles d'IST]
\begin{itemize}
\item Les IST \textbf{bactériennes} (gonococcie, syphilis, chlamydiose) \textbf{se guérissent} par des antibiotiques, à condition d'être traitées tôt ;
\item Les IST \textbf{virales} (VIH, herpès, hépatite B, HPV) \textbf{ne se guérissent pas} : on peut seulement contrôler l'infection (trithérapie antirétrovirale pour le VIH) ou la prévenir par la vaccination (hépatite B, HPV).
\end{itemize}
D'où l'importance capitale de la \textbf{prévention}.
\end{aretenir}

\courssub{Modes de transmission}

Les agents des IST circulent dans le \cle{sang}, le \cle{sperme}, les \cle{sécrétions vaginales} et le lait maternel des personnes infectées. Trois grandes voies de transmission :

\begin{enumerate}
\item \textbf{Voie sexuelle} (la principale) : rapports sexuels non protégés avec une personne infectée. Les micro-lésions des muqueuses génitales facilitent l'entrée des agents pathogènes ;
\item \textbf{Voie sanguine} : transfusion de sang contaminé non contrôlé, partage de seringues ou d'objets tranchants souillés (rasoirs, aiguilles de tatouage ou de scarification non stérilisées) ;
\item \textbf{Transmission mère-enfant} : pendant la grossesse, l'accouchement ou l'allaitement (VIH, syphilis, hépatite B). Un suivi médical de la femme enceinte permet de réduire fortement ce risque.
\end{enumerate}

\begin{attention}[Ce qui ne transmet PAS le VIH]
Le VIH ne se transmet \textbf{pas} par une poignée de main, un baiser sur la joue, le partage d'un repas ou de toilettes, ni par les piqûres de moustiques. Ces idées fausses nourrissent la stigmatisation injuste des personnes vivant avec le VIH.
\end{attention}

\courssub{Conséquences des IST non traitées}

Une IST négligée peut avoir des conséquences graves et parfois irréversibles :
\begin{itemize}
\item \textbf{Stérilité} : les chlamydioses et gonococcies non traitées provoquent chez la femme des infections des trompes de Fallope (salpingites) qui les bouchent, et chez l'homme des inflammations des épididymes ;
\item \textbf{Complications de la grossesse} : grossesse extra-utérine, fausses couches, accouchement prématuré, infection ou malformation du nouveau-né (syphilis congénitale) ;
\item \textbf{Cancers} : certains papillomavirus (HPV) sont responsables de la quasi-totalité des cancers du col de l'utérus, première cause de mortalité par cancer chez la femme dans de nombreux pays africains ; l'hépatite B chronique peut évoluer vers le cancer du foie ;
\item \textbf{Affaiblissement immunitaire} : le VIH détruit les lymphocytes T4 ; sans traitement, l'organisme devient vulnérable aux infections opportunistes (tuberculose, pneumonies) : c'est le stade du \cle{SIDA} (syndrome d'immunodéficience acquise) ;
\item \textbf{Effet amplificateur} : une IST non traitée (ulcérations, inflammations) multiplie le risque de contracter ou de transmettre le VIH.
\end{itemize}

\begin{exemplebox}[Un enjeu de santé publique pour les jeunes]
Selon l'ONUSIDA, la majorité des nouvelles infections à VIH chez les jeunes survient entre \textbf{15 et 24 ans}, et l'Afrique subsaharienne reste la région la plus touchée. Au Cameroun, les campagnes de dépistage gratuit dans les lycées, universités et centres de santé visent précisément cette tranche d'âge. Cela montre que les adolescent(e)s et jeunes adultes sont en première ligne : l'information et la prévention sont leurs meilleures armes.
\end{exemplebox}

\courssub{Les moyens de prévention : la stratégie ABC et au-delà}

\begin{propriete}[Double protection du préservatif (admise)]
Le \cle{préservatif} (masculin ou féminin), \textbf{correctement utilisé à chaque rapport}, est le \textbf{seul} moyen de contraception qui protège \emph{à la fois} des grossesses non désirées \emph{et} des IST. Aucun autre contraceptif (pilule, implant, stérilet, injectable) ne protège des IST.
\end{propriete}

La prévention des IST repose sur une stratégie simple, résumée par le sigle \cle{ABC}, complétée par le dépistage et la vaccination :

\begin{itemize}
\item \textbf{A — Abstinence} : l'abstinence sexuelle est la protection la plus sûre, car sans rapport sexuel, il n'y a pas de transmission sexuelle possible ;
\item \textbf{B — « Be faithful » : fidélité réciproque} : une relation stable avec un partenaire unique, non infecté, et mutuellement fidèle, réduit le risque au minimum. Justification : si aucun des deux partenaires n'est infecté et qu'aucun n'a de contact à risque à l'extérieur, l'agent pathogène ne peut pas entrer dans le couple ;
\item \textbf{C — Condom (préservatif)} : le préservatif constitue une \textbf{barrière mécanique} imperméable qui empêche le contact entre le sperme, les sécrétions vaginales et les muqueuses du partenaire : les agents pathogènes ne peuvent plus traverser ;
\item \textbf{Dépistage} : un test régulier (et avant tout nouveau couple) permet de connaître son statut, de se traiter tôt et de ne pas transmettre à son insu ;
\item \textbf{Vaccination} : il existe des vaccins très efficaces contre l'\textbf{hépatite B} et contre les \textbf{papillomavirus (HPV)}. Ils apprennent au système immunitaire à reconnaître le virus avant toute exposition.
\end{itemize}

\begin{methode}[Justifier scientifiquement un moyen de prévention]
Pour justifier un moyen de prévention lors d'un devoir ou d'un exposé, procède en trois étapes :
\begin{enumerate}
\item \textbf{Identifier le mode de transmission} de l'agent (voie sexuelle, sanguine, mère-enfant) ;
\item \textbf{Décrire le mécanisme d'action} de la barrière ou du moyen proposé (barrière mécanique, absence d'exposition, immunité vaccinale) ;
\item \textbf{Relier les deux} : montrer que le moyen \emph{coupe la chaîne de transmission} à un maillon précis.
\end{enumerate}
Exemple : le préservatif (barrière mécanique) bloque la voie sexuelle en empêchant le contact des sécrétions contaminées avec les muqueuses.
\end{methode}

\begin{popfigure}
\begin{tikzpicture}[font=\small, >=Stealth, node distance=4mm]
% Personnes
\node[circle, draw=popdark, fill=poppinkL, minimum size=1.1cm, label=above:{\textbf{Personne infectée}}] (A) at (0,0) {};
\node[circle, draw=popdark, fill=popgreenL, minimum size=1.1cm, label=above:{\textbf{Partenaire sain}}] (B) at (7,0) {};
% Voies de transmission
\draw[->, very thick, poppink] (A) -- node[above, align=center]{Voie sexuelle\\ (sperme, sécrétions)} (B);
\draw[->, very thick, poporange] (A) to[out=-35,in=-145] node[below, align=center]{Voie sanguine\\ (sang contaminé)} (B);
\draw[->, very thick, poppurple] (A) to[out=135,in=225, looseness=6] node[left, align=center, xshift=-3mm]{Mère $\rightarrow$ enfant\\ (grossesse, accouchement,\\ allaitement)} (A);
% Barrières ABC
\node[draw=popblue, fill=popblueL, rounded corners, align=center, minimum width=3.4cm] (abc) at (3.5,-3.4)
{\textbf{Barrières de prévention (ABC)}\\ A : Abstinence\\ B : Fidélité réciproque\\ C : Condom (préservatif)};
\draw[->, very thick, popblue, dashed] (abc.north) -- (3.5,-0.9) node[midway, right, align=left]{coupent la chaîne\\ de transmission};
\node[draw=popgreen, fill=popgreenL, rounded corners, align=center] (vac) at (9.2,-3.4)
{\textbf{Compléments}\\ Dépistage régulier\\ Vaccins : hépatite B, HPV};
\end{tikzpicture}
\legende{Modes de transmission des IST (flèches) et moyens de prévention qui coupent la chaîne de transmission : stratégie ABC, dépistage et vaccination.}
\end{popfigure}

\begin{exoresolu}[Justifier un moyen de prévention]
\textbf{Énoncé.} Lors d'une séance de sensibilisation dans un lycée de Douala, un élève affirme : « La pilule suffit, puisqu'elle évite la grossesse, on est protégé de tout. » À l'aide de tes connaissances, corrige cette affirmation en justifiant scientifiquement l'intérêt du préservatif.
\tcblower
\textbf{Solution détaillée.}
\begin{enumerate}
\item \textbf{Mode de transmission} : les agents des IST (VIH, gonocoques, chlamydiae, etc.) se transmettent lors des rapports sexuels par le contact entre le sperme ou les sécrétions vaginales contaminés et les muqueuses génitales du partenaire.
\item \textbf{Mécanisme de la pilule} : la pilule contraceptive agit par voie hormonale : elle bloque l'ovulation. Elle n'interpose \textbf{aucune barrière} entre les sécrétions et les muqueuses. Elle protège donc de la grossesse mais \textbf{pas des IST}.
\item \textbf{Mécanisme du préservatif} : le préservatif est une gaine imperméable qui recueille le sperme et isole les muqueuses des deux partenaires : il coupe la voie sexuelle de transmission des agents pathogènes \emph{et} empêche la fécondation.
\item \textbf{Conclusion} : l'affirmation est fausse. Seul le préservatif, correctement utilisé à chaque rapport, assure une double protection (grossesses + IST). La pilule peut être associée au préservatif, mais ne peut pas le remplacer pour la prévention des IST.
\end{enumerate}
\end{exoresolu}

\begin{savaistu}[Le vaccin anti-HPV : une victoire contre le cancer]
Le cancer du col de l'utérus est causé dans la quasi-totalité des cas par une infection persistante à certains papillomavirus (HPV). Or un \textbf{vaccin} existe : administré aux jeunes filles (et de plus en plus aux jeunes garçons) idéalement avant les premiers rapports sexuels, il protège contre les souches de HPV responsables de la plupart de ces cancers. Le Cameroun a intégré ce vaccin dans son Programme Élargi de Vaccination. C'est l'un des rares cas où un vaccin prévient directement un cancer !
\end{savaistu}

\courssub{Dépistage, traitement précoce et lutte contre la stigmatisation}

\textbf{Pourquoi se faire dépister ?} Parce que beaucoup d'IST sont silencieuses, le dépistage est le seul moyen de connaître son statut. Il est gratuit ou peu coûteux dans de nombreux centres de santé camerounais, et \textbf{confidentiel}. Un diagnostic précoce permet :
\begin{itemize}
\item de \textbf{guérir} les IST bactériennes avant les complications (stérilité) ;
\item de \textbf{contrôler} les IST virales : sous trithérapie bien suivie, une personne séropositive a une charge virale indétectable et ne transmet plus le VIH ;
\item de \textbf{protéger ses partenaires} et, pour une femme enceinte, son enfant à naître.
\end{itemize}

\textbf{Lutter contre la stigmatisation.} Rejeter, insulter ou isoler une personne atteinte d'une IST (en particulier du VIH) est à la fois injuste et contre-productif : la peur du rejet décourage le dépistage et favorise la propagation silencieuse des infections. Une personne vivant avec le VIH sous traitement peut étudier, travailler, fonder une famille. Adopter un comportement responsable, c'est donc à la fois \emph{se protéger}, \emph{se faire dépister}, \emph{se traiter} si besoin, et \emph{respecter} les personnes concernées.

\begin{aretenir}[L'essentiel de la section]
\begin{itemize}
\item Une IST est une infection transmise principalement lors de rapports sexuels non protégés ; beaucoup sont \textbf{asymptomatiques} ;
\item Les IST bactériennes se guérissent, les IST virales non (mais vaccins contre l'hépatite B et le HPV) ;
\item Conséquences graves : stérilité, cancers (col de l'utérus, foie), immunodéficience (SIDA) ;
\item Prévention : stratégie \textbf{ABC} (Abstinence, Fidélité réciproque, Condom), dépistage, vaccination ;
\item Le préservatif est le \textbf{seul} moyen protégeant à la fois des IST et des grossesses.
\end{itemize}
\end{aretenir}

\begin{exercice}[Pour t'entraîner]
\begin{enumerate}
\item Définis une IST et cite deux agents viraux et deux agents bactériens responsables d'IST.
\item Explique pourquoi une personne apparemment en bonne santé peut transmettre une IST.
\item Justifie scientifiquement, en trois étapes, l'efficacité de la fidélité réciproque comme moyen de prévention.
\item Pourquoi dit-on que le dépistage est un acte de protection à la fois individuel et collectif ?
\end{enumerate}
\end{exercice}

\coursec{Méthodes de contraception}

Dans la section précédente, nous avons étudié les infections sexuellement transmissibles (IST) et constaté que le préservatif est le seul moyen qui protège à la fois contre les IST et contre une grossesse non désirée. Mais les IST ne sont pas le seul risque lié à une sexualité précoce ou non protégée : il y a aussi la \cle{grossesse}. Une question se pose alors naturellement : existe-t-il des moyens fiables d'éviter une grossesse lorsqu'on ne désire pas (ou pas encore) d'enfant ? C'est toute la question de la \cle{contraception}, que nous allons maintenant explorer de manière rigoureuse.

\courssub{La contraception : définition et notion d'efficacité}

\textbf{Situation concrète.} Dans un centre de santé intégré de Yaoundé, une jeune femme de 19 ans, déjà mère d'un enfant, demande à l'infirmière : « Comment puis-je terminer mes études avant d'avoir un deuxième enfant ? » La réponse du personnel de santé repose sur un ensemble de techniques bien codifiées : les méthodes contraceptives.

\begin{definition}[Contraception]
La \cle{contraception} est l'ensemble des moyens, naturels ou artificiels, temporaires ou définitifs, permettant d'\textbf{éviter une grossesse} en empêchant la rencontre entre le spermatozoïde et l'ovocyte, en bloquant l'ovulation, ou en empêchant la nidation de l'œuf dans l'utérus.
\end{definition}

\begin{definition}[Efficacité d'une méthode contraceptive]
L'\cle{efficacité} d'une méthode contraceptive se mesure par le pourcentage de femmes \textbf{n'ayant pas} de grossesse après une année d'utilisation. On distingue :
\begin{itemize}
\item l'\textbf{efficacité théorique} (utilisation parfaite, sans erreur, indice de Pearl bas) ;
\item l'\textbf{efficacité pratique} (utilisation réelle, avec oublis, erreurs de manipulation, interactions médicamenteuses), toujours plus faible que l'efficacité théorique.
\end{itemize}
Une méthode est d'autant plus fiable que son efficacité pratique est proche de 100 \%. L'indice de Pearl (nombre de grossesses pour 100 femmes-an d'utilisation) est l'indicateur standard : plus il est bas, plus la méthode est efficace.
\end{definition}

\begin{attention}[Ce que la contraception ne fait pas]
Aucune méthode contraceptive, \textbf{sauf le préservatif (masculin ou féminin)}, ne protège contre les IST. Pilule, implant, stérilet, spermicides : tous empêchent (plus ou moins bien) la grossesse, mais \textbf{aucun ne bloque les agents pathogènes} (virus, bactéries, parasites). Ne jamais confondre « contraception » (anti-grossesse) et « prophylaxie » (protection contre les IST).
\end{attention}

\courssub{Les méthodes naturelles (basées sur l'observation du cycle)}

\textbf{Observation.} Chez la femme, la fécondation n'est possible que quelques jours par cycle, autour de l'ovulation (vers le 14\textsuperscript{e} jour d'un cycle standard de 28 jours). La période fertile s'étend de \textbf{3 à 5 jours avant l'ovulation} (durée de vie des spermatozoïdes dans les voies génitales féminines) \textbf{jusqu'à 24--48 h après} (durée de vie de l'ovocyte). D'où l'idée : éviter les rapports sexuels non protégés pendant cette « fenêtre de fertilité ».

\begin{itemize}
\item \textbf{L'abstinence périodique (méthode du calendrier ou d'Ogino-Knaus)} : le couple s'abstient de rapports pénis-vagin pendant la période estimée fertile, calculée à partir de la durée des 6 à 12 cycles précédents (formule : jour le plus court $-19$ à jour le plus long $-10$). \textbf{Limites majeures} : la date d'ovulation varie d'un cycle à l'autre (stress, fatigue, maladie, adolescence) ; le calcul est donc très incertain.
\item \textbf{La méthode de la température basale} : détection du léger pic de température (0,3 à 0,5 °C) post-ovulatoire (effet thermogène de la progestérone). \textbf{Limite} : ne signale l'ovulation qu'\emph{après} qu'elle a eu lieu (trop tard pour l'éviter), cycles irréguliers, fièvre, sommeil perturbé.
\item \textbf{La méthode de Billings (observation de la glaire cervicale)} : reconnaissance de la glaire filante, claire, abondante (type « blanc d'œuf ») signe d'imminence de l'ovulation. \textbf{Limite} : apprentissage long, subjectivité, infections vaginales modifient la glaire.
\item \textbf{Le retrait (coït interrompu)} : l'homme retire son pénis du vagin avant l'éjaculation. \textbf{Limites} : présence possible de spermatozoïdes dans le liquide pré-éjaculatoire (glandes de Cowper) ; exige un contrôle parfait, rarement réalisé chez l'adolescent.
\end{itemize}

Ces méthodes sont \textbf{gratuites, sans matériel ni effets secondaires médicamenteux}, mais leur efficacité pratique est \textbf{faible} (indice de Pearl élevé : 15 à 25 grossesses pour 100 femmes-an, soit 75 à 85 \% d'efficacité). Elles sont \textbf{fortement déconseillées chez l'adolescente} dont les cycles sont irréguliers.

\begin{savaistu}[Pourquoi la méthode du calendrier échoue-t-elle si souvent chez les adolescentes ?]
La méthode d'Ogino suppose des cycles \textbf{réguliers} (variation $<$ 8 jours entre le cycle le plus court et le plus long). Or, pendant l'adolescence, l'axe hypothalamo-hypophysaire qui commande l'ovulation n'est pas encore mature : les cycles sont souvent \textbf{anovulatoires ou irréguliers} (25 jours un mois, 40 jours le mois suivant). L'ovulation devient imprévisible et le « calcul » du calendrier ne vaut presque rien. C'est pourquoi de nombreuses grossesses précoces surviennent chez des adolescentes qui pensaient « ne pas être dans la période dangereuse ».
\end{savaistu}

\courssub{Les méthodes mécaniques (méthodes barrières)}

Le principe est simple : interposer un \textbf{obstacle physique} entre les spermatozoïdes déposés dans le vagin et l'ovocyte dans la trompe.

\begin{itemize}
\item \textbf{Le préservatif masculin} : gaine fine (latex ou polyuréthane pour allergiques au latex) placée sur le pénis en érection \emph{avant} tout contact génital, qui recueille le sperme. Efficacité pratique : environ 85 à 98 \% (indice de Pearl 2 à 15). \textbf{C'est le seul moyen qui protège aussi contre les IST, dont le VIH.} Il est peu coûteux, distribué gratuitement dans de nombreux centres de santé, lycées (clubs santé) et associations de lutte contre le sida au Cameroun.
\item \textbf{Le préservatif féminin} : gaine en nitrile (polyuréthane) souple, placée dans le vagin avant le rapport (anneau interne au fond du vagin, anneau externe recouvrant la vulve). Même double avantage (grossesse + IST), mais moins disponible, plus coûteux et demande un apprentissage.
\item \textbf{Le diaphragme (ou cape cervicale)} : coupelle en silicone ou latex que la femme place devant le col de l'utérus avant le rapport, \textbf{obligatoirement associé à un spermicide}. Efficacité pratique moyenne (environ 85 \%, indice de Pearl $\approx$ 15). Nécessite un ajustement de taille par un professionnel.
\item \textbf{Le stérilet ou DIU (Dispositif Intra-Utérin)} : petit objet en plastique en forme de T, muni de fils de retrait, placé \textbf{dans la cavité utérine} par un personnel de santé qualifié. Deux types :
  \begin{itemize}
  \item \textbf{DIU au cuivre} : ions cuivre toxiques pour les spermatozoïdes et l'ovocyte, inflammatoire local empêchant la nidation. Durée 5 à 10 ans.
  \item \textbf{DIU hormonal (LNG-IUS)} : diffuse une progestative (lévonorgestrel) localement : épaississement de la glaire cervicale (blocage spermatozoïdes), atrophie endométriale (anti-nidation), parfois inhibition ovulatoire. Durée 3 à 5 ans. Règles souvent moins abondantes.
  \end{itemize}
  Très efficace (efficacité pratique $>$ 99 \%, indice de Pearl $<$ 1), réversible à l'ablation, mais \textbf{ne protège pas des IST}. Contre-indiqué en IST active non traitée ou malformation utérine.
\end{itemize}

\courssub{Les méthodes hormonales}

\textbf{Problème scientifique.} On sait (section 1, physiologie de la reproduction) que l'ovulation est déclenchée par les hormones hypophysaires \cle{FSH} (folliculostimulante) et \cle{LH} (lutéinisante), elles-mêmes contrôlées par un rétrocontrôle exercé par les hormones ovariennes (œstrogènes et progestérone). Peut-on utiliser ce rétrocontrôle pour \emph{bloquer} l'ovulation ?

\begin{propriete}[Principe des contraceptifs hormonaux (admis au Probatoire)]
Les méthodes hormonales apportent dans le sang des \textbf{hormones de synthèse} (analogues des œstrogènes \emph{éthinylestradiol} et/ou de la progestérone \emph{progestatifs}). Ces hormones exercent un \textbf{rétrocontrôle négatif} permanent sur l'hypothalamus et l'hypophyse antérieure : la sécrétion pulsatile de \cle{GnRH}, puis de \cle{FSH} et de \cle{LH}, est inhibée. Sans FSH, le follicule ne mûrit pas ; sans pic de LH, \textbf{l'ovulation n'a pas lieu}. Sans ovocyte libéré, la fécondation est impossible.
\textbf{Effets secondaires associés} : épaississement de la glaire cervicale (barrière aux spermatozoïdes) et modification de l'endomètre (défavorable à la nidation), qui renforcent l'effet contraceptif.
\end{propriete}

\begin{experience}[Démarche historique : mise en évidence du rétrocontrôle]
\begin{enumerate}
\item \textbf{Observation} : pendant la grossesse, les taux d'œstrogènes et de progestérone (produits par le corps jaune puis le placenta) sont très élevés et \textbf{aucune ovulation ne se produit} (pas de règles, pas de nouveau follicule).
\item \textbf{Hypothèse} : ces hormones ovariennes inhibent l'hypophyse et empêchent l'ovulation (rétrocontrôle négatif).
\item \textbf{Expérience (années 1950, travaux de Pincus, Rock, Djerassi)} : on administre à des femmes volontaires des doses contrôlées d'hormones de synthèse (progestatifs puis association œstro-progestative) ; on mesure FSH et LH dans le sang et on suit les ovaires.
\item \textbf{Résultats} : les taux de FSH et LH s'effondrent au niveau basal ; aucun follicule n'arrive à maturité ; aucune ovulation n'est observée.
\item \textbf{Conclusion} : l'apport hormonal exogène bloque l'ovulation par rétrocontrôle négatif — c'est le principe de la pilule contraceptive orale combinée.
\end{enumerate}
\end{experience}

Les principales formes de contraception hormonale (toutes sur prescription médicale après bilan) :

\begin{itemize}
\item \textbf{La pilule combinée (œstro-progestative)} : comprimé pris chaque jour, à heure fixe (fenêtre de 12 h), pendant 21 jours suivis de 7 jours d'arrêt (règles de privation) ou en continu. Très efficace si prise correctement (efficacité théorique $>$ 99,7 \%, pratique $\approx$ 91 \% à cause des oublis, vomissements, diarrhées, interactions antibiotiques/anticonvulsivants). Contre-indications : thrombose veineuse, hypertension artérielle sévère, tabagisme $>$ 35 ans, migraine avec aura, cancer hormono-dépendant.
\item \textbf{La pilule progestative seule (micro-pilule)} : sans œstrogène. Prise \textbf{chaque jour à heure fixe} (fenêtre de 3 h pour les anciennes, 12 h pour les nouvelles désogestrel), sans arrêt. Indiquée si contre-indication aux œstrogènes (allaitement, thrombose, tabagisme). Efficacité pratique $\approx$ 91--93 \%. Règles souvent irrégulières (spotting).
\item \textbf{L'implant contraceptif (sous-cutané)} : petit bâtonnet souple (4 cm $\times$ 2 mm) inséré sous la peau de la face interne du bras (non dominant) par un personnel de santé formé, sous anesthésie locale. Diffuse un progestatif (etonogestrel ou lévonorgestrel) pendant \textbf{3 à 5 ans}. Efficacité pratique $>$ 99,9 \% (indice de Pearl $\approx$ 0,05), \textbf{sans risque d'oubli quotidien}. Règles imprévisibles (aménorrhée fréquente ou spotting). Retrait possible à tout moment.
\item \textbf{L'injectable (progestatif seul)} : injection intramusculaire (DMPA -- acétate de médroxyprogestérone) ou sous-cutanée tous les \textbf{3 mois} (13 semaines) au centre de santé. Efficacité pratique $>$ 99 \% si le calendrier est respecté. Retard de retour de fertilité possible (3 à 10 mois après l'arrêt). Perte de densité minérale osseuse réversible (surveillance si usage prolongé $>$ 2 ans).
\end{itemize}

\begin{attention}[Erreur fréquente au Probatoire : « la pilule protège de tout »]
\textbf{Faux !} La pilule (comme l'implant, l'injectable ou le stérilet) ne protège \textbf{que} contre la grossesse. Elle ne bloque ni le VIH, ni la gonococcie, ni la chlamydiose, ni l'hépatite B, ni la syphilis. \textbf{Seul le préservatif (masculin ou féminin) assure la double protection} (grossesse + IST). D'où la recommandation officielle de la \textbf{double protection} : préservatif à chaque rapport + une méthode contraceptive efficace (pilule, implant, injectable, DIU).
\end{attention}

\courssub{Les méthodes chimiques et les méthodes définitives}

\textbf{Les spermicides} sont des substances chimiques (crèmes, gels, ovules, éponges, films) placées dans le vagin \textbf{10 à 30 minutes avant} le rapport, qui \textbf{détruisent la membrane des spermatozoïdes ou les immobilisent}. Utilisés \textbf{seuls}, ils sont peu fiables (efficacité pratique $\approx$ 70 à 80 \%, indice de Pearl 20--30) ; ils servent surtout en \textbf{complément} d'une méthode barrière (diaphragme, cape cervicale) ou du préservatif (lubrifiant spermatocide). Non irritants pour la muqueuse saine, mais usage fréquent peut favoriser les vaginoses.

\textbf{Les méthodes définitives (stérilisation chirurgicale)} concernent les adultes (généralement $>$ 30 ans ou ayant déjà des enfants) qui ne désirent \textbf{plus aucun enfant} (projet parental achevé). Ce sont des méthodes quasi irréversibles (réparation chirurgicale complexe, succès non garanti), d'efficacité supérieure à 99,5 \%.
\begin{itemize}
\item \textbf{La ligature des trompes (salpingectomie ou occlusion tubaire)} chez la femme : les trompes utérines sont sectionnées, ligaturées, coagulées ou obturées par clips/anneaux (cœlioscopie ou mini-laparotomie). L'ovocyte ne peut plus rencontrer les spermatozoïdes ; l'ovulation et les règles persistent.
\item \textbf{La vasectomie} chez l'homme : les canaux déférents sont sectionnés et ligaturés au niveau du scrotum (anesthésie locale, 15 min). Le sperme éjaculé ne contient plus de spermatozoïdes (azoospermie) après $\approx$ 20 éjaculations / 3 mois (contrôle spermogramme obligatoire). La production de testostérone et l'éjaculation persistent.
\end{itemize}

\courssub{La contraception d'urgence : un filet de sécurité, pas une méthode régulière}

\begin{savaistu}[Contraception d'urgence (pilule du lendemain / DIU d'urgence)]
En cas de rapport sexuel non ou mal protégé (préservatif rompu, oubli de pilule $\ge$ 12 h, viol, rapport non consenti), il existe une \textbf{contraception d'urgence} :
\begin{itemize}
\item \textbf{Pilule progestative (lévonorgestrel 1,5 mg en dose unique)} : à prendre \textbf{le plus vite possible}, efficacité maximale dans les 12--24 h, possible jusqu'à 72 h (3 jours) après le rapport. Retarde l'ovulation si elle n'a pas encore eu lieu. Disponible \textbf{sans ordonnance} en pharmacie et centres de santé au Cameroun.
\item \textbf{Pilule anti-progestative (ulipristal acétate 30 mg)} : efficace jusqu'à \textbf{120 h (5 jours)} après le rapport. Sur ordonnance.
\item \textbf{DIU au cuivre (d'urgence)} : posé jusqu'à 5 jours après le rapport (ou 5 jours après la date présumée d'ovulation). C'est la \textbf{plus efficace} (quasi 100 \%) et assure une contraception durable ensuite.
\end{itemize}
\textbf{Attention} : ce n'est \textbf{pas} une méthode contraceptive régulière (efficacité moindre, troubles du cycle, coût). Elle ne protège pas des IST. Un test de grossesse est recommandé 3 semaines plus tard si pas de règles.
\end{savaistu}

\courssub{Comparaison de l'efficacité des méthodes}

\begin{center}
\begin{tabularx}{\linewidth}{|l|X|l|}
\hline
\rowcolor{popblueL} \textbf{Catégorie} & \textbf{Exemples principaux} & \textbf{Efficacité pratique indicative (1 an)} \\
\hline
Naturelles & Calendrier, température, Billings, retrait & 75 à 85 \% (Indice Pearl 15--25) \\
\hline
Chimiques & Spermicides seuls & 70 à 80 \% (Indice Pearl 20--30) \\
\hline
Mécaniques (barrières) & Préservatif masculin, féminin, diaphragme + spermicide & Préservatif : 85 à 98 \% ; Diaphragme : $\approx$ 85 \% \\
\hline
Hormonales (utilisateur-dépendantes) & Pilule combinée, micro-pilule & $\approx$ 91 à 93 \% (Indice Pearl 7--9) \\
\hline
Hormonales (longue durée, \textbf{indépendantes de l'utilisateur}) & Implant, Injectable, DIU hormonal & $>$ 99,5 \% (Indice Pearl $<$ 0,5) \\
\hline
Mécaniques (longue durée) & DIU au cuivre & $>$ 99,5 \% (Indice Pearl $<$ 0,8) \\
\hline
Définitives & Ligature des trompes, Vasectomie & $>$ 99,5 \% \\
\hline
\end{tabularx}
\end{center}

\begin{popfigure}
\begin{tikzpicture}
\begin{axis}[
    ybar, bar width=0.55cm,
    width=\linewidth, height=6.5cm,
    ymin=0, ymax=105,
    ylabel={Efficacité pratique (\%)},
    symbolic x coords={Spermicides, Calendrier, Retrait, Diaphragme, Préservatif, Pilule, Injectable, Implant, DIU Cuivre, DIU Hormonal},
    xtick=data,
    x tick label style={rotate=35, anchor=east, font=\small},
    nodes near coords, nodes near coords style={font=\small},
    every node near coord/.append style={yshift=2pt},
    axis lines=left,
    ytick={0,20,40,60,80,100},
    title style={font=\small, align=center},
    title={Efficacité pratique : méthodes courtes vs longues durées}
]
\addplot[fill=poporange, draw=popdark] coordinates {
    (Spermicides,72) (Calendrier,76) (Retrait,80) (Diaphragme,84) (Préservatif,87)
    (Pilule,91) (Injectable,99.5) (Implant,99.9) (DIU Cuivre,99.4) (DIU Hormonal,99.8)
};
\end{axis}
\end{tikzpicture}
\legende{Comparaison des taux d'efficacité pratique (utilisation réelle sur un an) des principales méthodes contraceptives. Les méthodes « utilisateur-dépendantes » (naturelles, chimiques, barrières, pilule) ont une efficacité plus variable. Les méthodes « longue durée » (implant, injectable, DIU) et définitives sont les plus efficaces car elles ne dépendent pas de l'observance quotidienne. \textbf{Rappel crucial : seul le préservatif (masculin ou féminin) protège aussi contre les IST.}}
\end{popfigure}

\begin{exoresolu}[Choisir et justifier une méthode adaptée]
\textbf{Énoncé.} Amina, 20 ans, étudiante à Douala, a un partenaire dont elle ne connaît pas le statut sérologique vis-à-vis du VIH. Elle ne souhaite pas de grossesse avant la fin de ses études (dans 3 ans). Son amie lui conseille la pilule combinée seule. Ce conseil est-il suffisant ? Justifie ta réponse en t'appuyant sur les connaissances du cours et propose une stratégie adaptée.
\tcblower
\textbf{Solution détaillée.}
\begin{enumerate}
\item \textbf{Analyse des besoins d'Amina} : double risque à couvrir -- (1) grossesse non désirée (projet d'études sur 3 ans) ; (2) IST/VIH (statut du partenaire inconnu, risque réel).
\item \textbf{Analyse du conseil « pilule seule »} : la pilule combinée a une efficacité pratique élevée contre la grossesse ($\approx$ 91 \%) \emph{si} elle est prise rigoureusement chaque jour. Cependant, elle \textbf{n'offre aucune protection contre les IST ni le VIH}. Le conseil est donc \textbf{insuffisant et dangereux} pour le risque infectieux.
\item \textbf{Stratégie adaptée -- La double protection} : utiliser le \textbf{préservatif masculin (ou féminin)} à \textbf{chaque rapport} (protection IST/VIH) \textbf{associé} à une méthode contraceptive très efficace contre la grossesse.
\item \textbf{Choix de la contraception pour Amina} : compte tenu de la durée (3 ans) et du risque d'oubli (vie étudiante chargée), une méthode \textbf{longue durée indépendante de l'utilisateur} est idéale : \textbf{l'implant (3--5 ans) ou le DIU (cuivre ou hormonal)}. La pilule reste possible si elle est certaine de sa régularité, mais l'implant/DIU évite le stress de l'oubli.
\item \textbf{Démarche pratique} : Amina doit consulter un \textbf{personnel de santé qualifié} (centre de santé intégré, service de planification familiale, hôpital de district) pour : bilan pré-thérapeutique (tension, antécédents), choix éclairé de la méthode, pose/prescription, et proposer à son partenaire un \textbf{dépistage VIH/IST} (gratuit et confidentiel au Cameroun).
\end{enumerate}
\textbf{Conclusion} : face à un double risque (grossesse + IST), la réponse correcte et responsable est la \textbf{double protection} : préservatif + contraception efficace (de préférence longue durée).
\end{exoresolu}

\courssub{Choisir une contraception et accéder aux services : la démarche responsable}

\begin{methode}[Bien choisir sa méthode contraceptive : la démarche en 5 étapes]
\begin{enumerate}
\item \textbf{Consulter un personnel de santé qualifié} (médecin, infirmier(ère) d'État, sage-femme) dans un centre de santé intégré, un hôpital de district/régional, ou un service de \cle{planification familiale} (PF) : \textbf{jamais d'automédication hormonale} sur conseil d'une amie, d'un vendeur de rue ou d'internet.
\item \textbf{Bénéficier d'un conseil pré-contraceptif} : entretien confidentiel (antécédents médicaux, mode de vie, tabagisme, projet d'enfant, risque IST), examen clinique si nécessaire (tension artérielle, examen gynécologique pour DIU/implant), information sur les avantages, inconvénients, effets secondaires et signes d'alerte de chaque méthode.
\item \textbf{Choisir en toute autonomie éclairée} : tenir compte de l'efficacité pratique (préférer méthodes longue durée si risque d'oubli), de la tolérance, de la réversibilité, du coût (souvent subventionné au Cameroun), et de la protection IST (ajout systématique du préservatif si risque).
\item \textbf{Se faire suivre} : consultation de contrôle à 3 mois (puis annuel), signaler tout effet indésirable inhabituel (céphalées sévères, douleurs abdominales, saignements anormaux, signe de phlébite), possibilité de changer de méthode si insatisfaction.
\item \textbf{Promouvoir la responsabilité partagée} : impliquer le partenaire (port du préservatif, accompagnement aux consultations, dépistage couple, soutien en cas d'effet secondaire).
\end{enumerate}
\end{methode}

Au Cameroun, les services de \cle{planification familiale} des hôpitaux (ex: Hôpital Central, Hôpital Gynéco-Obstétrique et Pédiatrique de Yaoundé, Hôpital Laquintinie de Douala) et des centres de santé intégrés (CSI), ainsi que des associations spécialisées (CAMNAFAW, ACMS, UNFPA partenaires), informent, conseillent et délivrent les méthodes contraceptives, souvent à coût réduit ou gratuit pour les jeunes. Les campagnes de sensibilisation (radios communautaires, clubs de santé scolaire, journées mondiales : contraception 26 sept, sida 1er déc, population 11 juil) jouent un rôle essentiel pour briser les tabous.

\begin{aretenir}[L'essentiel de la section -- Contraception]
\begin{itemize}
\item La \cle{contraception} regroupe tous les moyens d'éviter une grossesse : méthodes naturelles (peu fiables, déconseillées à l'adolescence), mécaniques/barrières, hormonales (blocage de l'ovulation par rétrocontrôle négatif sur FSH/LH + effets sur glaire/endomètre), chimiques (spermicides, adjuvant seulement) et définitives (stérilisation).
\item L'\cle{efficacité pratique} varie fortement : de 70--85 \% (naturelles, chimiques, barrières simples) à $>$ 99,5 \% (implant, injectable, DIU, stérilisation). Les méthodes \textbf{longue durée (LARC)} sont les plus efficaces car indépendantes de l'observance quotidienne.
\item \textbf{Seul le préservatif (masculin ou féminin)} protège à la fois de la grossesse et des IST (dont VIH) : c'est le pilier de la \textbf{double protection} (préservatif + contraception efficace).
\item La \textbf{contraception d'urgence} (pilule lévonorgestrel jusqu'à 72 h, ulipristal jusqu'à 120 h, DIU cuivre jusqu'à 120 h) est un secours, \textbf{pas une méthode régulière}.
\item Le choix d'une méthode se fait \textbf{avec un personnel de santé qualifié}, dans un centre de santé ou service de planification familiale, après conseil personnalisé. C'est un droit (santé reproductive) et un acte de responsabilité.
\end{itemize}
\end{aretenir}

\begin{exercice}[Pour t'entraîner -- Type Probatoire]
\begin{enumerate}
\item Classe les méthodes suivantes par ordre d'efficacité pratique \textbf{croissante} (de la moins efficace à la plus efficace) : Implant, Retrait, Pilule combinée, Préservatif masculin, Spermicide, DIU au cuivre.
\item Pour chacune de ces six méthodes, indique si elle protège (\textbf{OUI} / \textbf{NON}) contre les infections sexuellement transmissibles (IST).
\item Une élève de Première D utilise la méthode du calendrier (Ogino). Elle a des cycles de durée variable : 26, 31, 28, 35 jours. Explique scientifiquement pourquoi cette méthode est particulièrement risquée pour elle.
\item Cite les deux composantes de la « double protection » et justifie pourquoi elle est recommandée pour un adolescent sexuellement actif.
\end{enumerate}
\end{exercice}

\begin{corrige}[Exercice 1 -- Contraception]
\begin{enumerate}
\item \textbf{Ordre croissant d'efficacité pratique} (valeurs indicatives) :
Spermicide ($\approx$ 72 \%) $<$ Retrait ($\approx$ 80 \%) $<$ Calendrier ($\approx$ 76 \%, très variable) $<$ Préservatif masculin ($\approx$ 87 \%) $<$ Pilule combinée ($\approx$ 91 \%) $<$ DIU au cuivre ($>$ 99,4 \%) $<$ Implant ($>$ 99,9 \%).
(Note : Calendrier et Retrait sont très proches et peu fiables ; l'ordre exact dépend des études, mais tous deux $<$ Préservatif $<$ Pilule $<$ DIU/Implant).
\item \textbf{Protection IST} :
Spermicide : \textbf{NON} ; Retrait : \textbf{NON} ; Calendrier : \textbf{NON} ; Préservatif masculin : \textbf{OUI} (seul de la liste) ; Pilule combinée : \textbf{NON} ; DIU au cuivre : \textbf{NON} ; Implant : \textbf{NON}.
\item \textbf{Explication} : La méthode d'Ogino suppose des cycles réguliers pour prédire la période fertile (calcul basé sur les cycles les plus courts et les plus longs). Chez cette élève, les cycles varient de 26 à 35 jours (amplitude 9 jours $>$ 8 jours seuil de régularité). L'ovulation est imprévisible (peut survenir vers J12 comme vers J21). De plus, la survie des spermatozoïdes (3--5 j) et de l'ovocyte (24--48 h) élargit la fenêtre de fertilité réelle bien au-delà du calcul théorique. Risque de grossesse très élevé.
\item \textbf{Double protection} = \textbf{Préservatif} (protection IST/VIH) \textbf{+ Contraception efficace} (pilule, implant, DIU, injectable... protection grossesse). Recommandée car l'adolescent sexuellement actif est exposé aux \textbf{deux risques simultanés} : grossesse non désirée (perturbation scolarité, santé maternelle) et IST/VIH (séquelles, mortalité). Aucune méthode unique (sauf préservatif) ne couvre les deux, et le préservatif seul a une efficacité contraceptive pratique inférieure aux méthodes hormonales/DIU.
\end{enumerate}
\end{corrige}

Maintenant que nous connaissons les moyens d'éviter une grossesse non désirée et de se protéger des IST, il reste à situer ces outils dans une démarche plus large : la \textbf{planification familiale} et l'adoption d'un \textbf{comportement responsable} en matière de santé reproductive, objet de la dernière section de cette leçon.

\coursec{Planification familiale et comportement responsable}

Dans la section précédente, nous avons étudié les \cle{méthodes de contraception} : préservatif, pilule, injectables, implants, stérilet\ldots{} Ces outils permettent d'éviter une grossesse non désirée. Mais une question demeure : \emph{à quoi servent-ils concrètement dans la vie d'un couple et d'une société ?} C'est toute la question de la \cle{planification familiale}, qui n'est pas une simple technique, mais un véritable \textbf{projet de vie}. Cette dernière section va aussi nous permettre de boucler la leçon : de la puberté à la contraception, tout converge vers un objectif unique — permettre à chacun, et d'abord à toi, adolescent(e), de faire des \textbf{choix libres, éclairés et responsables} concernant sa santé reproductive.

\courssub{Qu'est-ce que la planification familiale ?}

\textbf{Une situation concrète pour commencer.}\par

Dans un quartier de Yaoundé, deux familles voisines ont des revenus comparables. Chez les Mbarga, quatre enfants sont nés en cinq ans ; la mère, épuisée, a dû interrompre ses études, le dernier bébé est né prématurément et l'aînée a été retirée de l'école pour aider à la maison. Chez les Etoundi, les trois enfants sont nés à trois ans d'intervalle ; la mère a pu poursuivre son activité commerciale, chaque enfant est vacciné, suivi médicalement et scolarisé. Même revenu, même quartier : la différence tient à une décision — \emph{choisir quand avoir ses enfants et combien en avoir}.

\begin{definition}[Planification familiale]
La \cle{planification familiale} est l'ensemble des moyens et des démarches qui permettent à un individu ou à un couple de \textbf{choisir librement et de manière éclairée} :
\begin{itemize}
\item le \textbf{nombre d'enfants} qu'il souhaite avoir ;
\item le \textbf{moment} des naissances et leur \textbf{espacement} dans le temps.
\end{itemize}
Elle repose sur l'information, le dialogue au sein du couple, l'accès aux services de santé et l'utilisation éventuelle des méthodes de contraception étudiées précédemment.
\end{definition}

\begin{attention}[Erreur fréquente à éviter]
Beaucoup croient que « planification familiale » signifie « \textbf{ne pas avoir d'enfants} » ou « limiter les naissances de force ». C'est \textbf{faux}. La planification familiale signifie \textbf{choisir quand et combien} : un couple peut décider d'avoir cinq enfants — c'est aussi de la planification, dès lors que ce choix est libre, éclairé et compatible avec la santé de la mère et les ressources de la famille. Elle ne s'impose jamais : elle se décide.
\end{attention}

\courssub{Pourquoi planifier les naissances ? Les intérêts de la planification familiale}

\textbf{Observation et données.}\par

Les données des enquêtes démographiques et de santé menées au Cameroun montrent un fait frappant : les enfants nés moins de deux ans après leur aîné présentent un risque de décès avant l'âge de 5 ans nettement plus élevé que ceux nés après un intervalle de trois ans ou plus. De même, les grossesses trop rapprochées épuisent les réserves de la mère (fer, calcium, protéines) et multiplient les complications : anémie, accouchement prématuré, bébé de faible poids.

\textbf{Hypothèse.} L'intervalle entre deux naissances influence directement la santé de la mère et celle de l'enfant.

\textbf{Interprétation.} La grossesse et l'allaitement mobilisent énormément de ressources chez la mère. Si une nouvelle grossesse survient avant que l'organisme maternel ne soit reconstitué, le fœtus et le nourrisson déjà né entrent en « compétition » pour des réserves insuffisantes. Un intervalle suffisant laisse le temps au corps de la mère de récupérer et à l'aîné d'être sevré et suivi.

\textbf{Conclusion.} Espacer les naissances est une mesure de santé publique efficace, simple et gratuite.

\begin{propriete}[Espacement des naissances — propriété admise]
Un \cle{espacement des naissances} d'\textbf{au moins 2 ans} (idéalement 3 ans) entre deux accouchements réduit significativement :
\begin{itemize}
\item la \textbf{mortalité maternelle} (moins d'hémorragies, d'anémies sévères, de complications de l'accouchement) ;
\item la \textbf{mortalité infantile} (moins de prématurité, de faible poids de naissance et de malnutrition) ;
\item la charge physique et psychologique de la mère.
\end{itemize}
Cette propriété, établie par de vastes études épidémiologiques, est admise : elle n'est pas à démontrer au Probatoire, mais tu dois savoir l'\textbf{expliquer} et l'\textbf{argumenter}.
\end{propriete}

\begin{popfigure}
\begin{center}
\begin{tikzpicture}
\begin{axis}[
    width=0.85\linewidth, height=6.5cm,
    ybar, bar width=0.9cm,
    symbolic x coords={$<$ 2 ans, 2--3 ans, $>$ 3 ans},
    xtick=data,
    ymin=0, ymax=100,
    ylabel={Risque relatif de décès avant 5 ans (\%)},
    xlabel={Intervalle entre deux naissances},
    ylabel style={font=\small}, xlabel style={font=\small},
    nodes near coords, every node near coord/.append style={font=\small\bfseries},
    axis line style={popline}, tick style={popline},
    ymajorgrids=true, grid style={popline!40},
]
\addplot[fill=poporange, draw=popdark] coordinates {($<$ 2 ans,100) (2--3 ans,60) ($>$ 3 ans,45)};
\end{axis}
\end{tikzpicture}
\end{center}
\legende{Risque relatif de mortalité infantile (avant 5 ans) en fonction de l'intervalle entre deux naissances (valeurs indicatives, risque rapporté à 100 pour un intervalle inférieur à 2 ans). Plus l'intervalle est long, plus le risque diminue : c'est l'argument chiffré en faveur de l'espacement des naissances.}
\end{popfigure}

\textbf{Les trois grands intérêts de la planification familiale.}\par

\begin{enumerate}
\item \textbf{Santé de la mère et de l'enfant.} Des grossesses espacées et désirées réduisent la mortalité maternelle et infantile, la prématurité et la malnutrition. La mère peut allaiter chaque enfant suffisamment longtemps et se remettre de chaque accouchement.
\item \textbf{Équilibre économique familial.} Chaque enfant représente des dépenses : alimentation, soins, habillement, fournitures scolaires. Des naissances planifiées permettent à la famille d'assumer dignement ces charges au lieu de s'appauvrir progressivement.
\item \textbf{Scolarisation et épanouissement des enfants.} Une famille dont les ressources ne sont pas diluées peut envoyer \emph{tous} ses enfants — filles comprises — à l'école, plus longtemps. La planification familiale est donc aussi un investissement dans l'avenir de toute une génération.
\end{enumerate}

\begin{exemplebox}[Deux scénarios comparés]
Une jeune femme devient mère à 16 ans, puis à 17 ans et demi, puis à 19 ans : trois grossesses rapprochées, scolarité interrompue, enfants fragiles. Une autre devient mère à 24 ans après ses études, puis à 27 ans : deux grossesses désirées, suivies, espacées de 3 ans. La planification familiale n'a pas changé le \emph{nombre} d'enfants de manière autoritaire ; elle a changé le \emph{moment} et les \emph{conditions}, donc le destin de toute la famille.
\end{exemplebox}

\courssub{Les acteurs de l'éducation à la santé reproductive}

La santé reproductive ne se construit pas seul(e) dans sa chambre : c'est un effort collectif où chacun a un rôle précis.

\begin{itemize}
\item \textbf{La famille} est le premier cadre d'éducation. Des parents qui parlent sans tabou de la puberté, du respect du corps et des responsabilités donnent à leurs enfants des repères solides. Le silence familial laisse le champ libre aux rumeurs de la rue.
\item \textbf{L'école} transmet des connaissances scientifiques rigoureuses (comme dans cette leçon !), développe l'esprit critique face aux idées reçues et forme aux compétences de vie : savoir dire non, savoir négocier, savoir chercher de l'aide.
\item \textbf{Les services de santé} (centres de santé intégrés, hôpitaux de district, unités de planification familiale) offrent information confidentielle, conseil, dépistage des IST, contraception et suivi des grossesses.
\item \textbf{Les médias, les associations et les chefs communautaires} relaient les campagnes de sensibilisation (lutte contre le VIH/Sida, contre les grossesses précoces).
\end{itemize}

\begin{savaistu}[La planification familiale au Cameroun]
Le Cameroun dispose d'un réseau de \textbf{centres de santé intégrés} et d'\textbf{unités de planification familiale} présents dans les hôpitaux et centres de santé, y compris en zone rurale. Ces structures accueillent aussi les \textbf{jeunes et adolescents} pour des consultations confidentielles : information sur la puberté, dépistage, conseil contraceptif. Des campagnes nationales, souvent appuyées par des partenaires comme l'UNFPA, distribuent gratuitement des préservatifs et sensibilisent dans les écoles et les quartiers. Se rendre dans une formation sanitaire pour s'informer est un \textbf{droit}, pas une honte.
\end{savaistu}

\courssub{Droits et responsabilités des adolescents}

\begin{definition}[Droits en santé reproductive]
Tout adolescent et toute adolescente possède des \cle{droits} en matière de santé reproductive :
\begin{itemize}
\item le droit à l'\textbf{information} exacte et complète sur son corps et la reproduction ;
\item le droit à la \textbf{santé} et à l'accès aux soins et au conseil ;
\item le droit de \textbf{décider librement} de sa vie affective et sexuelle, sans contrainte ni violence ;
\item le droit de \textbf{refuser} tout rapport sexuel, y compris forcé, et tout mariage imposé ;
\item le droit à la \textbf{confidentialité} dans les services de santé.
\end{itemize}
À ces droits correspondent des \textbf{responsabilités} : respecter le corps et le consentement de l'autre, se protéger et protéger son ou sa partenaire, assumer les conséquences de ses choix.
\end{definition}

\begin{attention}[Rapports forcés : un refus absolu]
Aucun rapport sexuel ne doit jamais être subi : ni sous la pression d'un petit ami (« si tu m'aimes, tu dois\ldots{} »), ni sous la menace, ni en échange d'argent ou de cadeaux, ni dans le cadre d'un mariage forcé. Le \textbf{rapport forcé est une violence} (et un crime) : il faut le refuser, s'en éloigner et \textbf{en parler immédiatement} à un adulte de confiance, à un enseignant ou à un service de santé. Le silence protège l'agresseur, jamais la victime.
\end{attention}

\courssub{Promouvoir un comportement responsable}

Être responsable en santé reproductive, ce n'est pas seulement « éviter les problèmes » : c'est construire une vie affective fondée sur le respect de soi et de l'autre. Plusieurs attitudes concrètes s'offrent à l'adolescent(e) :

\begin{itemize}
\item \textbf{L'abstinence responsable} : choisir d'attendre est un choix légitime, efficace à 100 \% contre les grossesses et les IST, qui permet de se concentrer sur ses études et sa maturité affective.
\item \textbf{La fidélité} : dans un couple, la fidélité mutuelle réduit fortement le risque d'IST.
\item \textbf{La protection} : le préservatif, correctement utilisé, protège à la fois des grossesses non désirées et des IST, dont le VIH.
\item \textbf{La communication} : parler franchement avec son ou sa partenaire, avec ses parents, avec un professionnel de santé.
\item \textbf{Le respect de l'autre} : le consentement libre et éclairé est la condition de toute relation.
\end{itemize}

\begin{methode}[Adopter et promouvoir un comportement responsable : 5 étapes]
\begin{enumerate}
\item \textbf{S'informer} : chercher des informations fiables (cours de SVT, personnel de santé, documents officiels) et rejeter les rumeurs.
\item \textbf{Dialoguer} : parler avec ses parents, ses amis, son ou sa partenaire ; ne jamais décider sous pression.
\item \textbf{Se protéger} : choisir l'abstinence ou utiliser systématiquement un moyen de protection adapté.
\item \textbf{Se faire dépister} : en cas de rapport à risque, consulter rapidement un centre de santé pour un dépistage des IST.
\item \textbf{Respecter autrui} : exiger et offrir le consentement, refuser toute forme de pression ou de violence, et encourager ses proches à faire de même.
\end{enumerate}
\end{methode}

\textbf{Lutter contre les idées reçues et les pressions sociales.}\par

De nombreuses rumeurs circulent : « le préservatif rend stérile », « la pilule fait grossir et détruit le ventre », « une fille qui se protège est une fille facile », « un vrai homme a des enfants tôt ». Toutes sont \textbf{fausses} et dangereuses. La démarche scientifique que tu as apprise — vérifier la source, confronter aux données, demander l'avis d'un professionnel — est ta meilleure arme. De même, la pression du groupe (« tout le monde le fait ») n'est jamais un argument : un choix responsable est un choix \emph{personnel}, assumé et éclairé.

\begin{exoresolu}[Analyser un cas et proposer un comportement responsable]
\textbf{Énoncé.} Sandrine, 17 ans, élève en Première D à Douala, fréquente un jeune homme de 22 ans qui lui demande d'avoir des rapports non protégés « pour prouver son amour », en affirmant que « les préservatifs rendent malade ». Une amie lui conseille de céder pour ne pas le perdre. En t'appuyant sur les notions de la leçon, analyse la situation et propose à Sandrine une conduite responsable.
\tcblower
\textbf{Solution détaillée.}
\begin{enumerate}
\item \textbf{Identification des risques.} Un rapport non protégé expose Sandrine à un \textbf{double risque} : une grossesse précoce non désirée (interruption probable de sa scolarité, complications liées à son jeune âge) et une IST, dont le VIH/Sida.
\item \textbf{Démantèlement des fausses affirmations.} L'affirmation « les préservatifs rendent malade » est une \textbf{idée reçue} sans fondement scientifique : le préservatif est au contraire le seul moyen qui protège simultanément des grossesses et des IST. L'argument « prouver son amour » est une \textbf{pression psychologique} : l'amour n'exige jamais de prendre des risques pour sa santé.
\item \textbf{Droits de Sandrine.} Elle a le droit de \textbf{refuser} tout rapport, d'exiger le respect de son consentement, de s'informer et de consulter en confidentialité un centre de santé.
\item \textbf{Conduite responsable proposée.} Sandrine peut : choisir l'\textbf{abstinence} et se concentrer sur ses études ; ou, si elle se sent prête, \textbf{dialoguer} avec son partenaire pour imposer l'usage systématique du \textbf{préservatif} ; se rendre dans une \textbf{unité de planification familiale} pour un conseil adapté et confidentiel ; et refuser tout rapport si le partenaire persiste dans son chantage — un partenaire qui ne respecte pas un refus ne respecte pas la personne.
\item \textbf{Conclusion.} Le comportement responsable consiste à décider librement, en étant informée, sans céder aux pressions, et à protéger sa santé et son avenir.
\end{enumerate}
\end{exoresolu}

\courssub{Synthèse de la leçon : une démarche cohérente}

Reprenons le fil de toute la leçon. La \textbf{puberté} transforme le corps de l'adolescent(e) et le rend \emph{capable de procréer} — sans lui donner encore la maturité pour assumer une parentalité. Cette capacité nouvelle expose à deux risques majeurs : les \textbf{grossesses précoces et non désirées}, aux lourdes conséquences sanitaires, scolaires et sociales, et les \textbf{infections sexuellement transmissibles}, parfois gravissimes comme le VIH/Sida. Face à ces risques, la science et la société offrent des réponses : les \textbf{moyens de prévention} (abstinence, fidélité, préservatif) et les \textbf{méthodes de contraception}. Enfin, la \textbf{planification familiale} inscrit ces outils dans un projet de vie : choisir librement le nombre et le moment des naissances, pour la santé de la mère et des enfants, l'équilibre de la famille et l'avenir de la société.

\begin{popfigure}
\begin{center}
\begin{tikzpicture}[
  mindmap, grow cyclic,
  every node/.style={concept, text=white, font=\small\bfseries},
  root concept/.append style={concept color=popblue, font=\bfseries},
  level 1/.append style={level distance=4.0cm, sibling angle=72, text width=2.8cm, align=center},
  concept color=popblue,
]
\node[root concept]{Santé reproductive de l'adolescent(e)}
  child[concept color=popgreen] { node {Puberté et physiologie de la reproduction} }
  child[concept color=poporange] { node {Grossesses précoces non désirées} }
  child[concept color=poppink] { node {IST et prévention} }
  child[concept color=poppurple] { node {Méthodes de contraception} }
  child[concept color=popgold] { node {Planification familiale et comportement responsable} };
\end{tikzpicture}
\end{center}
\legende{Carte mentale récapitulative de la leçon : les cinq thèmes s'articulent autour d'un objectif central — permettre à l'adolescent(e) de vivre une santé reproductive libre, informée et responsable.}
\end{popfigure}

\begin{aretenir}[L'essentiel de la section]
\begin{itemize}
\item La \textbf{planification familiale} est le choix libre et éclairé du \textbf{nombre d'enfants} et de l'\textbf{espacement des naissances} — elle ne signifie pas « ne pas avoir d'enfants ».
\item Un espacement d'\textbf{au moins 2 ans} entre deux naissances réduit significativement la mortalité maternelle et infantile.
\item Ses intérêts : santé de la mère et de l'enfant, équilibre économique familial, scolarisation des enfants.
\item L'école, la famille et les services de santé partagent la responsabilité de l'éducation à la santé reproductive.
\item L'adolescent(e) a des \textbf{droits} (information, soins, consentement, refus des rapports forcés) et des \textbf{responsabilités} (respect de l'autre, protection).
\item Comportement responsable : s'informer, dialoguer, se protéger, se faire dépister, respecter autrui — et combattre les idées reçues.
\end{itemize}
\end{aretenir}

\begin{exercice}[Pour t'entraîner]
\begin{enumerate}
\item Définis la planification familiale et corrige l'affirmation suivante : « La planification familiale, c'est l'interdiction d'avoir beaucoup d'enfants. »
\item Explique, en t'appuyant sur la physiologie de la grossesse et de l'allaitement, pourquoi un intervalle d'au moins 2 ans entre deux naissances protège la santé de la mère et de l'enfant.
\item Ton jeune frère affirme : « Seuls les adultes mariés ont le droit d'entrer dans une unité de planification familiale. » Que lui réponds-tu ?
\item Rédige un court message de sensibilisation (5 lignes maximum) destiné aux élèves de ta classe pour promouvoir un comportement responsable en santé reproductive.
\end{enumerate}
\end{exercice}

\begin{corrige}[Exercice]
\begin{enumerate}
\item La planification familiale est l'ensemble des moyens permettant à un individu ou un couple de choisir librement et en connaissance de cause le nombre d'enfants et l'espacement des naissances. L'affirmation est fausse : il n'y a aucune interdiction ; un couple peut choisir d'avoir beaucoup d'enfants, l'essentiel étant que le choix soit libre, éclairé et compatible avec la santé de la mère et les ressources familiales.
\item La grossesse et l'allaitement épuisent les réserves de la mère (fer, calcium, protéines). Si une nouvelle grossesse survient trop tôt, la mère n'a pas reconstitué ses réserves : risques d'anémie, d'accouchement prématuré, de bébé de faible poids, et l'aîné risque le sevrage précoce et la malnutrition. Un intervalle d'au moins 2 ans laisse le corps maternel récupérer et l'aîné grandir suffisamment, ce qui réduit la mortalité maternelle et infantile.
\item C'est faux : les unités de planification familiale et les centres de santé intégrés accueillent aussi les jeunes et adolescents, de manière confidentielle, pour s'informer, se faire conseiller et dépister. S'y rendre est un droit.
\item Exemple de message : « Notre avenir se construit aujourd'hui : informons-nous sur notre santé, respectons-nous mutuellement, refusons les pressions, protégeons-nous et n'hésitons pas à consulter un centre de santé. Un choix éclairé est un choix libre ! »
\end{enumerate}
\end{corrige}

\newpage

\coursec{Activité d'intégration}

\coursec{Lutte contre les problèmes liés à la santé reproductive des adolescent(e)s}

\courssub{Objectifs de l'activité}

À l'issue de cette activité d'intégration, tu seras capable de :

\begin{itemize}
\item \cle{analyser} un cas concret de risque de grossesse précoce en exploitant un graphique du cycle menstruel (repérage de la période fertile) ;
\item \cle{expliquer} les conséquences sanitaires, scolaires et sociales d'une grossesse précoce et non désirée ;
\item \cle{identifier} les infections sexuellement transmissibles (IST) encourues lors d'un rapport non protégé et \cle{justifier} les moyens de prévention adaptés ;
\item \cle{proposer} et \cle{promouvoir} un comportement responsable en matière de santé reproductive (dialogue, consultation, dépistage, double protection, planification familiale).
\end{itemize}

Cette activité mobilise l'ensemble des savoir-faire de la leçon dans une situation unique et réaliste, comme au Probatoire où l'on te demande souvent de résoudre un problème de santé à partir de documents.

\courssub{Situation}

\textbf{Dossier : le parcours d'Amina.}

Amina est une élève de Première D dans un lycée de Maroua (région de l'Extrême-Nord). Elle a 16 ans. Depuis ses premières règles à 13 ans, son cycle menstruel s'est peu à peu régularisé : il dure en moyenne \textbf{28 jours}, avec des règles de 4 à 5 jours. Le médecin du centre de santé intégré de son quartier lui a remis, lors d'une séance de sensibilisation organisée par le club santé du lycée, un graphique type du cycle menstruel de 28 jours (reproduit ci-dessous), montrant l'évolution des hormones (LH, FSH, progestérone) et de la température corporelle, ainsi que la période d'ovulation.

Un soir de fête de fin d'année scolaire, Amina a eu un \textbf{rapport sexuel non protégé} avec son ami Ibrahim, 19 ans. Ce rapport a eu lieu au \textbf{12\textsuperscript{ème} jour de son cycle} (le jour 1 étant le premier jour des règles). Ibrahim avait \textbf{refusé d'utiliser le préservatif}, affirmant que « ce n'est pas naturel » et qu'« il n'a jamais été malade ». Amina n'utilise \textbf{aucune méthode de contraception}. Elle ne connaît pas les antécédents médicaux d'Ibrahim, qui a déjà eu plusieurs partenaires.

Trois semaines plus tard, Amina s'inquiète : ses règles n'arrivent pas, elle ressent des nausées matinales et une tension des seins. Elle confie sa détresse à sa grande sœur, étudiante en médecine, qui décide de l'accompagner au centre de santé.

\begin{popfigure}
\begin{center}
\begin{tikzpicture}
\begin{axis}[
width=0.95\linewidth, height=6.5cm,
xlabel={Jours du cycle (jour 1 = premier jour des règles)},
ylabel={Concentration hormonale (u. a.)},
xmin=0, xmax=29, ymin=0, ymax=10,
xtick={1,5,10,12,14,16,20,25,28},
ytick=\empty,
legend style={at={(0.02,0.97)},anchor=north west,font=\small},
axis lines=left, grid=both, grid style={popline!40}
]
\addplot[popcurve,popblue,domain=1:28,samples=150]{3*exp(-((x-14)^2)/2) + 0.5};
\addlegendentry{LH}
\addplot[popcurve,poporange,domain=1:28,samples=150]{1.5*exp(-((x-13)^2)/6) + 0.4};
\addlegendentry{FSH}
\addplot[popcurve,popgreen,domain=1:28,samples=150]{4*exp(-((x-21)^2)/18) + 0.3};
\addlegendentry{Progestérone}
\draw[dashed,popink,thick] (axis cs:14,0) -- (axis cs:14,9.5) node[above,font=\small]{Ovulation (J14)};
\fill[poporange!25] (axis cs:9,0) rectangle (axis cs:16,10);
\node[font=\small\bfseries,popdark] at (axis cs:12.5,9.2) {Période fertile};
\draw[<-,thick,popink] (axis cs:12,1.2) -- (axis cs:12,3.5) node[above,font=\small\itshape]{Rapport non protégé (J12)};
\end{axis}
\end{tikzpicture}
\end{center}
\legende{Cycle menstruel de 28 jours : variations de la LH, de la FSH et de la progestérone. Le pic de LH (J14) déclenche l'ovulation. La période fertile (zone colorée, environ J9 à J16) tient compte de la durée de vie des spermatozoïdes (3 à 5 jours) et de l'ovocyte (environ 24 h).}
\end{popfigure}

\courssub{Consignes}

Par groupes de 4 à 6 élèves, et à l'aide du dossier et du graphique, réponds aux questions suivantes, puis prépare une présentation orale de 5 minutes.

\begin{enumerate}
\item À partir du graphique, repère le jour de l'\cle{ovulation} et délimite la \cle{période fertile} du cycle d'Amina. Justifie tes réponses en t'appuyant sur les durées de vie des gamètes.
\item Le rapport non protégé ayant eu lieu au 12\textsuperscript{ème} jour, évalue le \cle{risque de grossesse} encouru par Amina. Argumente.
\item Les signes décrits par Amina (retard de règles, nausées, tension des seins) sont-ils compatibles avec une grossesse ? Quel examen permet de confirmer le diagnostic ?
\item Dresse la liste des \cle{conséquences sanitaires, scolaires et sociales} possibles d'une grossesse à 16 ans, en distinguant les conséquences pour Amina, pour l'enfant à naître et pour la famille.
\item Quelles \cle{IST} Amina a-t-elle pu contracter lors de ce rapport non protégé ? Cite au moins quatre IST et précise l'agent responsable de chacune.
\item Justifie \cle{trois moyens de prévention} des IST qu'Amina et Ibrahim auraient dû (ou devraient) utiliser, en expliquant le mécanisme de protection de chacun.
\item Propose un \cle{plan d'action responsable} pour Amina à partir d'aujourd'hui : dialogue, consultation, dépistage, contraception. Précise le rôle de la \cle{double protection}.
\item Rédigez collectivement, pour toute la classe, une \cle{charte du comportement responsable} en matière de santé reproductive (5 à 8 engagements).
\end{enumerate}

\courssub{Ressources à mobiliser}

\begin{aretenir}[Ce qu'il faut avoir en tête]
\begin{itemize}
\item \textbf{Cycle menstruel (28 jours)} : l'ovulation a lieu vers le \cle{14\textsuperscript{ème} jour}, déclenchée par le pic de \cle{LH}. L'ovocyte n'est fécondable que \textbf{12 à 24 h}, mais les spermatozoïdes survivent \textbf{3 à 5 jours} dans les voies génitales féminines : la période fertile s'étend donc d'environ \textbf{5 jours avant à 1 jour après l'ovulation}.
\item \textbf{Signes de grossesse} : aménorrhée (absence de règles), nausées, tension mammaire ; confirmation par un \textbf{test de grossesse} (recherche de l'hormone \cle{hCG} dans les urines ou le sang).
\item \textbf{IST principales} : VIH/Sida (virus), blennorragie (gonocoque, bactérie), chlamydiose (bactérie), syphilis (tréponème, bactérie), hépatite B (virus), papillomavirus (HPV, virus), trichomonose (parasite).
\item \textbf{Prévention des IST} : abstinence, fidélité réciproque, \cle{préservatif} (masculin ou féminin, seul moyen protégeant à la fois des IST et de la grossesse), dépistage, vaccination (hépatite B, HPV).
\item \textbf{Contraception} : méthodes naturelles (abstinence périodique, peu fiables), barrières (préservatif), hormonales (pilule, injectable, implant), mécaniques (stérilet), d'urgence (pilule du lendemain, à usage exceptionnel).
\item \textbf{Double protection} : préservatif + méthode contraceptive efficace = protection contre les IST \emph{et} contre la grossesse.
\end{itemize}
\end{aretenir}

\courssub{Démarche et corrigé commenté}

\begin{exoresolu}[Étude du cas d'Amina : corrigé détaillé]
\textbf{Énoncé.} À partir du dossier et du graphique du cycle menstruel : (1) repérer la période fertile et évaluer le risque de grossesse ; (2) lister les conséquences possibles ; (3) identifier les IST encourues et justifier trois moyens de prévention ; (4) proposer un plan d'action responsable.

\tcblower

\textbf{Étape 1 — Repérage de la période fertile (question 1).}

Sur le graphique, le \cle{pic de LH} est observé au \textbf{14\textsuperscript{ème} jour} : c'est lui qui déclenche l'\cle{ovulation}, c'est-à-dire la libération de l'ovocyte par l'ovaire. L'ovocyte ne reste fécondable que \textbf{12 à 24 heures}. En revanche, les spermatozoïdes déposés dans les voies génitales peuvent survivre \textbf{3 à 5 jours}. Un rapport survenu plusieurs jours \emph{avant} l'ovulation peut donc aboutir à une fécondation. La période fertile s'étend ainsi approximativement du \textbf{9\textsuperscript{ème} au 16\textsuperscript{ème} jour} du cycle (zone colorée du graphique).

\textbf{Étape 2 — Évaluation du risque de grossesse (questions 2 et 3).}

Le rapport non protégé a eu lieu au \textbf{12\textsuperscript{ème} jour}, soit \textbf{2 jours avant l'ovulation}, en pleine période fertile. Des spermatozoïdes déposés à J12 peuvent encore être vivants à J14 et féconder l'ovocyte. Le \textbf{risque de grossesse est donc maximal} : c'est statistiquement l'un des jours les plus favorables à la conception du cycle.

Les signes cliniques d'Amina — \textbf{retard de règles} (aménorrhée), \textbf{nausées matinales}, \textbf{tension des seins} — sont les signes d'appel classiques d'un début de grossesse. Le diagnostic sera confirmé au centre de santé par un \textbf{test de grossesse} détectant l'hormone \cle{hCG} (gonadotrophine chorionique), sécrétée dès la nidation, dans les urines ou le sang.

\textbf{Étape 3 — Conséquences d'une grossesse précoce à 16 ans (question 4).}

\begin{center}
\begin{tabularx}{\linewidth}{|l|X|}\hline
\rowcolor{popblueL} \textbf{Niveau} & \textbf{Conséquences possibles} \\ \hline
Sanitaire (Amina) & Risques accrus d'hypertension gravidique, d'anémie, de complications de l'accouchement (bassin pas totalement mature), de fistule obstétricale ; risque d'avortement clandestin (hémorragies, infections, stérilité, décès). \\ \hline
Sanitaire (enfant) & Prématurité, faible poids de naissance, mortalité infantile plus élevée. \\ \hline
Scolaire & Interruption des études (déscolarisation), échec au Probatoire, perte des perspectives professionnelles. \\ \hline
Sociale et familiale & Rejet ou stigmatisation, mariage forcé, précarité économique, responsabilités précoces, conflits familiaux. \\ \hline
\end{tabularx}
\end{center}

\textbf{Étape 4 — IST encourues (question 5).}

Un rapport non protégé avec un partenaire aux antécédents inconnus expose à toutes les IST, notamment : le \textbf{VIH/Sida} (rétrovirus), la \textbf{blennorragie} (\emph{Neisseria gonorrhoeae}, bactérie), la \textbf{chlamydiose} (\emph{Chlamydia trachomatis}, bactérie), la \textbf{syphilis} (\emph{Treponema pallidum}, bactérie), l'\textbf{hépatite B} (virus) et les \textbf{papillomavirus} (HPV, impliqués dans le cancer du col de l'utérus). Beaucoup de ces infections peuvent être \textbf{asymptomatiques} au début : l'argument d'Ibrahim (« je n'ai jamais été malade ») ne prouve rien.

\textbf{Étape 5 — Trois moyens de prévention justifiés (question 6).}

\begin{enumerate}
\item \textbf{Le préservatif (masculin ou féminin)} : il constitue une \cle{barrière mécanique} empêchant le contact entre les sécrétions génitales et le sperme ; c'est le \textbf{seul} moyen qui protège à la fois des IST \emph{et} de la grossesse. Correctement utilisé à chaque rapport, son efficacité dépasse 95 \%.
\item \textbf{L'abstinence ou le report des premiers rapports} : en l'absence de rapport, aucun risque d'IST ni de grossesse ; c'est un choix légitime et respectable, surtout à 16 ans.
\item \textbf{Le dépistage et la fidélité réciproque} : un couple dont les deux partenaires sont dépistés négatifs et fidèles supprime le risque d'IST ; le dépistage permet aussi un traitement précoce (les IST bactériennes se guérissent par antibiotiques ; l'infection à VIH se contrôle par antirétroviraux).
\end{enumerate}

On peut ajouter la \textbf{vaccination} contre l'hépatite B et les HPV.

\textbf{Étape 6 — Plan d'action responsable pour Amina (question 7).}

\textbf{Un plan d'action en cinq points :}
\begin{enumerate}
\item \textbf{Dialoguer} : Amina parle avec sa sœur, ses parents ou un adulte de confiance, et exige d'Ibrahim le respect de son choix (aucun rapport sans préservatif).
\item \textbf{Consulter rapidement} le centre de santé : test de grossesse (hCG) et, le cas échéant, information sur la contraception d'urgence \emph{pour tout futur rapport à risque} (à prendre dans les 72 h, idéalement dans les 12 h). \emph{Nota : le rapport ayant eu lieu il y a trois semaines, la contraception d'urgence n'est plus possible aujourd'hui.}
\item \textbf{Se faire dépister} (VIH, syphilis, hépatite B, chlamydiose), elle et son partenaire ; le dépistage du VIH est gratuit et confidentiel dans de nombreuses formations sanitaires du Cameroun.
\item \textbf{Adopter la double protection} : préservatif systématique + méthode contraceptive adaptée (pilule, implant ou injectable) prescrite par un personnel de santé.
\item \textbf{Poursuivre sa scolarité} quoi qu'il arrive et s'appuyer sur les services de planification familiale du centre de santé.
\end{enumerate}

\textbf{Étape 7 — Charte du comportement responsable (question 8, exemple de production de classe).}

\begin{enumerate}
\item Je reporte mes premiers rapports ou j'utilise \textbf{systématiquement} un préservatif.
\item Je refuse tout rapport sous pression : mon corps m'appartient.
\item Je me fais dépister et je connais le statut de mon partenaire.
\item Je me renseigne sur la contraception auprès d'un personnel de santé qualifié.
\item Je respecte le choix de l'autre et je dialogue sans violence.
\item Je soutiens sans jugement toute personne confrontée à une grossesse ou à une IST.
\item Je privilégie mes études et mon avenir.
\end{enumerate}
\end{exoresolu}

\courssub{Bilan de l'activité}

Cette étude de cas a permis de mettre en évidence plusieurs points essentiels :

\begin{itemize}
\item La \cle{connaissance du cycle menstruel} est un outil d'évaluation du risque : un rapport non protégé en période fertile (ici J12, deux jours avant l'ovulation) expose à un risque \textbf{maximal} de grossesse ; les méthodes dites « naturelles » basées sur le calendrier sont trop incertaines pour être fiables, surtout avec des cycles irréguliers.
\item Une \cle{grossesse précoce} engage la santé de l'adolescente et de l'enfant, la scolarité et l'avenir social : c'est un problème de santé publique majeur au Cameroun.
\item Un rapport non protégé expose \textbf{simultanément} à la grossesse et aux IST, dont certaines sont silencieuses : l'absence de symptômes chez le partenaire ne garantit rien.
\item Le \cle{préservatif}, seul outil de \cle{double protection}, associé au dialogue, au dépistage et à une contraception adaptée, constitue la stratégie responsable.
\item Enfin, la construction collective de la charte montre que le comportement responsable repose sur des savoir-être : respect de soi et d'autrui, refus de la pression, recours aux services de santé, solidarité sans stigmatisation.
\end{itemize}

\begin{attention}[Erreurs fréquentes à éviter]
Ne confonds pas \textbf{contraception} (évite la grossesse) et \textbf{prévention des IST} : la pilule ne protège d'aucune IST. Ne considère jamais le retrait ou la méthode des températures comme fiables chez l'adolescente. Enfin, la contraception d'urgence (« pilule du lendemain ») est un rattrapage exceptionnel, pas une méthode régulière.
\end{attention}

\newpage

\coursec{Méthodes et exercices résolus}

\coursec{Physiologie de la reproduction humaine et puberté}

\courssub{Appareils reproducteurs mâle et femelle}

\begin{definition}[Appareil reproducteur féminin]
L'appareil reproducteur féminin assure la \cle{gamétogenèse} (production d'ovocytes), la \cle{fécondation}, la \cle{grossesse} et l'\cle{accouchement}. Il comprend : les \cle{ovaires} (gonades, production d'ovocytes et d'hormones : œstrogènes, progestérone), les \cle{trompes de Fallope} (site de la fécondation, transport de l'ovocyte/zygote), l'\cle{utérus} (organe musculaire creux, site de la nidation et du développement fœtal ; sa muqueuse interne est l'\cle{endomètre}), le \cle{col de l'utérus} (communication utérus-vagin, sécrète la \cle{glaire cervicale} dont la perméabilité varie au cours du cycle), le \cle{vagin} (organe de la copulation, voie d'accouchement) et la \cle{vulve} (organes génitaux externes).
\end{definition}

\begin{popfigure}
\begin{tikzpicture}[scale=0.85, every node/.style={font=\small}, >=Stealth]
% Ovary
\draw[poporange, thick, fill=poporangeL] (5.5,4.2) ellipse (0.7 and 0.4);
\node[poporange, right=0.2cm] at (6.2,4.2) {\textbf{Ovaire}};
% Fallopian tube (fimbriae to uterus)
\draw[popblue, thick] (4.8,4.0) .. controls (4.0,4.5) and (3.0,4.5) .. (2.5,3.5)
.. controls (2.0,2.5) .. (2.5,2.0);
\node[popblue, above] at (3.5,4.7) {\textbf{Trompe de Fallope}};
\node[popblue, above left] at (4.8,4.0) {Pavillon};
% Uterus (body)
\draw[poppurple, thick, fill=poppurpleL] (2.5,2.0) -- (3.5,2.0) -- (3.3,0.2) -- (2.7,0.2) -- cycle;
% Endometrium
\draw[poppink, thick, fill=poppinkL] (2.7,2.0) -- (3.3,2.0) -- (3.15,0.5) -- (2.85,0.5) -- cycle;
\node[poppurple, left=0.2cm] at (2.5,1.2) {\textbf{Utérus}};
\node[poppink, left=0.2cm] at (2.5,0.8) {Endomètre};
% Cervix
\draw[popgreen, thick, fill=popgreenL] (2.7,0.2) -- (3.3,0.2) -- (3.2,-0.5) -- (2.8,-0.5) -- cycle;
\node[popgreen, right=0.2cm] at (3.3,-0.15) {\textbf{Col de l'utérus}};
\node[popgreen, right=0.2cm] at (3.3,-0.5) {Glaire cervicale};
% Vagina
\draw[popdark, thick, fill=popcream] (2.5,-0.5) -- (3.5,-0.5) -- (3.3,-2.0) -- (2.7,-2.0) -- cycle;
\node[popdark, right=0.2cm] at (3.5,-1.2) {\textbf{Vagin}};
% Vulva
\draw[popdark, thick] (2.5,-2.0) -- (3.5,-2.0);
\node[popdark, right=0.2cm] at (3.5,-2.0) {\textbf{Vulve}};
% Annotations
\draw[popline, <->] (4.8,4.2) -- node[above, sloped, pos=0.5] {Ovulation / Capture ovocyte} (3.5,3.5);
\draw[popline, ->] (3.0,1.5) -- node[left] {Nidation} (2.9,1.0);
\end{tikzpicture}
\legende{Appareil reproducteur féminin (vue sagittale médiane). L'ovaire libère l'ovocyte capté par le pavillon de la trompe ; la fécondation a lieu dans l'ampoule de la trompe ; le zygote migre vers l'utérus pour la nidation dans l'endomètre.}
\end{popfigure}

\begin{definition}[Appareil reproducteur mâle]
L'appareil reproducteur mâle assure la \cle{spermatogenèse} (production de spermatozoïdes) et leur émission. Il comprend : les \cle{testicules} (gonades, logés dans le scrotum à \SI{34}{\celsius}, production de spermatozoïdes et de \cle{testostérone}), les \cle{épididymes} (maturation et stockage des spermatozoïdes), les \cle{canaux déférents} (transport), les \cle{vésicules séminales}, la \cle{prostate} et les \cle{glandes de Cowper} (glandes annexes sécrétant le \cle{liquide séminal} qui nourrit et protège les spermatozoïdes), l'\cle{urètre} (voie commune urine/sperme) et le \cle{pénis} (organe de la copulation).
\end{definition}

\begin{popfigure}
\begin{tikzpicture}[scale=0.85, every node/.style={font=\small}, >=Stealth]
% Testis
\draw[poporange, thick, fill=poporangeL] (1,3.5) ellipse (0.8 and 0.5);
\node[poporange, left=0.2cm] at (0.2,3.5) {\textbf{Testicule}};
\node[poporange, left=0.2cm] at (0.2,3.0) {(dans le scrotum)};
% Epididymis
\draw[popblue, thick, fill=popblueL] (1.8,3.5) -- (2.5,3.7) -- (2.5,3.3) -- cycle;
\node[popblue, above] at (2.2,3.8) {\textbf{Épididyme}};
% Vas deferens
\draw[popblue, thick] (2.5,3.5) -- (4.5,3.5) -- (5.0,2.5) -- (5.0,1.5);
\node[popblue, above] at (3.5,3.7) {\textbf{Canal déférent}};
% Seminal vesicle
\draw[poppurple, thick, fill=poppurpleL] (4.8,2.5) ellipse (0.4 and 0.6);
\node[poppurple, right=0.1cm] at (5.2,2.5) {\textbf{Vésicule séminale}};
% Prostate
\draw[popgreen, thick, fill=popgreenL] (4.5,1.0) circle (0.5);
\node[popgreen, below] at (4.5,0.3) {\textbf{Prostate}};
% Cowper glands
\draw[poppink, thick, fill=poppinkL] (4.0,0.0) circle (0.2);
\draw[poppink, thick, fill=poppinkL] (5.0,0.0) circle (0.2);
\node[poppink, below] at (4.5,-0.5) {Glandes de \textbf{Cowper}};
% Urethra / Penis
\draw[popdark, thick] (4.5,1.0) -- (4.5,-1.0) -- (4.0,-2.0) -- (5.0,-2.0) -- (4.5,-1.0);
\node[popdark, right=0.2cm] at (5.0,-1.0) {\textbf{Urètre}};
\node[popdark, right=0.2cm] at (5.0,-2.0) {\textbf{Pénis}};
% Bladder
\draw[popmuted, thick, dashed, fill=popmuted!20] (3.5,2.5) ellipse (0.8 and 0.4);
\node[popmuted] at (3.5,2.5) {Vessie};
\end{tikzpicture}
\legende{Appareil reproducteur mâle (vue sagittale médiane). Les spermatozoïdes produits dans les testicules mûrissent dans l'épididyme, remontent par les canaux déférents, reçoivent les sécrétions des vésicules séminales, de la prostate et des glandes de Cowper pour former le sperme, émis par l'urètre.}
\end{popfigure}

\courssub{Gamétogenèse : spermatogenèse et oogenèse}

\begin{propriete}[Spermatogenèse]
Processus continu à partir de la puberté dans les \cle{tubules séminifères} des testicules.
\begin{enumerate}
\item \cle{Phase de multiplication} : Cellules souches (spermatogonies, $2n$) $\xrightarrow{\text{mitoses}}$ spermatogonies.
\item \cle{Phase de croissance} : Spermatogonie $\rightarrow$ \cle{spermatocyte I} ($2n$, $2c$ ADN).
\item \cle{Phase de maturation (méiose)} :
\begin{itemize}
\item Méiose I : Spermatocyte I ($2n$) $\rightarrow$ 2 \cle{spermatocytes II} ($n$, $2c$ ADN).
\item Méiose II : 2 Spermatocytes II $\rightarrow$ 4 \cle{spermatides} ($n$, $1c$ ADN).
\end{itemize}
\item \cle{Spermiogenèse} : Différenciation des spermatides en \cle{spermatozoïdes} (flagelle, acrosome, noyau haploïde condensé). Durée totale : \textbf{$\approx$ 74 jours}.
\end{enumerate}
\end{propriete}

\begin{propriete}[Oogenèse]
Processus discontinu, débuté \emph{in utero} et achevé à la fécondation.
\begin{enumerate}
\item \emph{Vie fœtale} : Multiplication mitotique des \cle{ogonies} $\rightarrow$ \cle{ovocytes I} ($2n$) qui entrent en \cle{prophase I de méiose} et s'\cle{arrêtent} (stade dictyotène) jusqu'à la puberté. Constitution de la \cle{réserve ovarienne} (stock fini).
\item \emph{Puberté à ménopause} : À chaque cycle, quelques ovocytes I reprennent la méiose I. Un seul (généralement) aboutit à l'\cle{ovulation} : Ovocyte I $\xrightarrow{\text{Méiose I}}$ \cle{Ovocyte II} ($n$, $2c$ ADN) + \cle{globule polaire I}. L'ovocyte II commence la méiose II et s'\cle{arrête en métaphase II}.
\item \emph{Fécondation} : L'entrée du spermatozoïde déclenche la fin de la méiose II : Ovocyte II $\rightarrow$ \cle{Ovule} ($n$, $1c$ ADN) + \cle{globule polaire II}. Fusion des noyaux $\rightarrow$ \cle{Zygote} ($2n$).
\end{enumerate}
\end{propriete}

\begin{aretenir}[Comparaison spermatogenèse / oogenèse]
\begin{center}
\begin{tabularx}{\linewidth}{|l|X|X|}\hline
\rowcolor{popblueL} Critère & Spermatogenèse (Mâle) & Oogenèse (Femelle) \\ \hline
Lieu & Tubules séminifères (testicules) & Ovaires (follicules) \\ \hline
Début & Puberté (continue, vie entière) & Vie fœtale (arrêt prophase I) \\ \hline
Méiose I & Continue & Un par cycle (ovulation) \\ \hline
Méiose II & Immédiate après Méiose I & Arrêt en métaphase II ; fin \textbf{si fécondation} \\ \hline
Produits finaux & 4 spermatozoïdes fonctionnels & 1 ovule fonctionnel + 2-3 globules polaires \\ \hline
Durée & $\approx$ 74 jours & Années (arrêt prophase I) + quelques jours/cycle \\ \hline
Cytoplasme & Partagé équitablement & Majoritairement dans l'ovule (réserves) \\ \hline
\end{tabularx}
\end{center}
\end{aretenir}

\courssub{Cycle ovarien et cycle menstruel}

\begin{propriete}[Cycle ovarien (durée moyenne 28 jours)]
Régulé par l'axe \cle{hypothalamo-hypophysaire} (GnRH $\rightarrow$ FSH, LH) et les \cle{hormones ovariennes} (œstrogènes, progestérone).
\begin{itemize}
\item \textbf{Phase folliculaire (J1--J14)} : FSH stimule la croissance de follicules $\rightarrow$ sécrétion d'\textbf{œstrogènes} (estradiol). Œstrogènes $\rightarrow$ rétrocontrôle \emph{négatif} sur FSH (sélection follicule dominant) puis \emph{positif} sur LH (pic LH).
\item \textbf{Ovulation (J14)} : \cle{Pic de LH} $\rightarrow$ rupture du follicule de Graaf $\rightarrow$ libération de l'ovocyte II + formation du \cle{corps jaune}.
\item \textbf{Phase lutéale (J15--J28)} : Corps jaune sécrète \textbf{progestérone} (et œstrogènes). Rétrocontrôle \emph{négatif} sur FSH et LH $\rightarrow$ atrophie du corps jaune (J26) $\rightarrow$ chute brutale progestérone/œstrogènes $\rightarrow$ règles (si pas de grossesse). \textbf{Durée constante : 14 jours}.
\end{itemize}
\end{propriete}

\begin{popfigure}
\begin{tikzpicture}[scale=0.9]
\begin{axis}[
width=\linewidth, height=8cm,
axis lines=left,
xlabel={Jour du cycle}, ylabel={Concentration (UI/L) ou Épaisseur (mm)},
xmin=0, xmax=29, xtick={1,7,14,21,28}, xticklabels={1, 7, 14, 21, 28},
ymin=0, ymax=120,
legend style={at={(0.5,1.02)}, anchor=south, legend columns=3, font=\tiny, cells={anchor=west}},
domain=1:28, samples=200,
]
% FSH
\addplot[popblue, thick] {10 + 5*sin(deg((x-1)*360/28)) + 15*exp(-(x-14)^2/5)};
\addlegendentry{FSH}
% LH
\addplot[poporange, thick] {5 + 50*exp(-(x-14)^2/2)};
\addlegendentry{LH}
% Estradiol
\addplot[poppurple, thick] {20 + 40*exp(-(x-13)^2/10) + 15*exp(-(x-21)^2/15)};
\addlegendentry{Œstradiol}
% Progesterone
\addplot[popgreen, thick] {5 + 40*exp(-(x-21)^2/12)};
\addlegendentry{Progestérone}
% Ovulation line
\draw[popline, dashed, thick] (axis cs:14,0) -- (axis cs:14,120) node[above, pos=0.95, font=\tiny, popline] {Ovulation};
% Phases shading
\fill[popblueL!30] (axis cs:1,0) rectangle (axis cs:14,120);
\fill[poporangeL!30] (axis cs:14,0) rectangle (axis cs:28,120);
\node[font=\tiny, popblue] at (axis cs:7,110) {Phase folliculaire};
\node[font=\tiny, poporange] at (axis cs:21,110) {Phase lutéale};
\end{axis}
\end{tikzpicture}
\legende{Variations hormonales au cours d'un cycle ovarien de 28 jours. Le pic de LH (jour 14) déclenche l'ovulation. La phase lutéale (progestérone) dure 14 jours constants.}
\end{popfigure}

\begin{propriete}[Cycle utérin (menstruel) -- Variations de l'endomètre]
\begin{itemize}
\item \textbf{Phase menstruelle (J1--J5)} : Chute hormones $\rightarrow$ vasoconstriction artères spirales $\rightarrow$ nécrose/désquamation de l'endomètre fonctionnel $\rightarrow$ \cle{règles} (saignement).
\item \textbf{Phase proliférative (J6--J14)} : Œstrogènes $\rightarrow$ reconstruction de l'endomètre (prolifération glandes, vaisseaux). Épaisseur : 2 $\rightarrow$ \SI{10}{mm}.
\item \textbf{Phase sécrétoire (J15--J28)} : Progestérone $\rightarrow$ épaississement, vascularisation, sécrétion glycogène (nutrition embryon). Épaisseur : \SI{10}{mm} $\rightarrow$ \SI{14}{mm}. Si pas de grossesse $\rightarrow$ ischémie $\rightarrow$ règles.
\end{itemize}
\end{propriete}

\begin{definition}[Période fertile]
Fenêtre de fécondité maximale. L'ovule vit \textbf{12 à 24 h}. Les spermatozoïdes survivent \textbf{3 à 5 jours} dans la glaire cervicale (permébile sous effet œstrogènes).
$\Rightarrow$ Période fertile = \textbf{5 jours avant l'ovulation à 1 jour après}.
\textbf{Formule} : Jour ovulation $\approx$ Durée du cycle $-$ 14.
\end{definition}

\courssub{Fécondation, grossesse et accouchement}

\begin{propriete}[Fécondation et nidation]
\begin{enumerate}
\item \textbf{Fécondation} : Dans l'ampoule de la trompe. Un spermatozoïde franchit la \cle{zone pellucide} (réaction acrosomique) $\rightarrow$ fusion membranes $\rightarrow$ blocage polyspermie (réaction corticale) $\rightarrow$ achèvement méiose II ovocyte $\rightarrow$ formation pronucléus mâle et femelle $\rightarrow$ fusion $\rightarrow$ \cle{zygote} ($2n$, \SI{46}{chromosomes}).
\item \textbf{Segmentation \& Migration} : Divisions mitotiques (clivages) $\rightarrow$ morula (J3) $\rightarrow$ \cle{blastocyste} (J5, cavité blastocélique, masse cellulaire interne = embryoblast, trophoblaste). Migration vers utérus.
\item \textbf{Nidation (J6--J7)} : Éclosion (zona pellucida) $\rightarrow$ apposition/advasion trophoblaste dans endomètre sécrétoire $\rightarrow$ formation \cle{placenta} (échange mère-fœtus sans mélange sanguin).
\end{enumerate}
\end{propriete}

\begin{propriete}[Grossesse et accouchement]
\begin{itemize}
\item \textbf{Hormones placentaires} : \cle{hCG} (maintien corps jaune 1\textsuperscript{er} trimestre), \cle{progestérone} (quiescence utérine, immunotolérance), \cle{œstrogènes} (vascularisation, croissance utérus), \cle{lakrogène placentaire} (métabolisme maternel).
\item \textbf{Accouchement} : Déclenché par chute progestérone / montée œstrogènes $\rightarrow$ récepteurs ocytocine + prostaglandines $\rightarrow$ contractions utérines (travail) $\rightarrow$ expulsion fœtus $\rightarrow$ délivrance (placenta).
\end{itemize}
\end{propriete}

\courssub{La puberté}

\begin{definition}[Puberté]
Période de transition de l'enfance à l'âge adulte reproductif, marquée par la maturation de l'axe hypothalamo-hypophysaire-gonadique (réveil pulsatile de la \cle{GnRH}) et l'apparition des \cle{caractères sexuels secondaires}.
\end{definition}

\begin{aretenir}[Signes pubertaires principaux]
\begin{center}
\begin{tabularx}{\linewidth}{|l|X|X|}\hline
\rowcolor{popblueL} & \textbf{Fille (10--14 ans)} & \textbf{Garçon (12--16 ans)} \\ \hline
\emph{Premier signe} & \textbf{Thélarchée} (bourgeon mammaire) & \textbf{Augmentation volume testicules} ($> \SI{4}{ml}$) \\ \hline
\emph{Autres signes} & Pubarchée (poils pubiens), ménarchée (1\textsuperscript{ères} règles, $\approx$ 2 ans après thélarchée), pic de croissance, élargissement bassin. & Pubarchée, phallochée (croissance pénis), mue (voix), pilosité faciale/axillaire, pic de croissance (plus tardif), spermatogenèse active (spermarchée). \\ \hline
\emph{Hormones clés} & Œstrogènes (ovaires), LH/FSH. & Testostérone (testicules), LH/FSH. \\ \hline
\end{tabularx}
\end{center}
\end{aretenir}

\begin{savaistu}[Puberté précoce et tardive au Cameroun]
Au Cameroun, l'âge moyen de la ménarchée est d'environ \textbf{12,5--13 ans}. Une puberté est dite \emph{précoce} avant 8 ans (fille) / 9 ans (garçon) et \emph{tardive} absence de signes à 13 ans (fille) / 14 ans (garçon). Les causes peuvent être constitutionnelles, nutritionnelles (obésité $\rightarrow$ leptine), ou pathologiques (tumeurs, kystes ovariens). La prise en charge précoce au centre de santé (pédiatrie, endocrinologie) évite des conséquences sur la taille finale et la psychologie de l'adolescent.
\end{savaistu}

\coursec{Grossesses précoces et non désirées}

\courssub{Définitions et épidémiologie}

\begin{definition}[Grossesse précoce et non désirée]
Une \cle{grossesse précoce} survient chez une fille de \textbf{moins de 18 ans}, dont l'organisme (bassin, utérus, métabolisme) n'a pas achevé sa maturation. Elle est \cle{non désirée} lorsqu'elle survient sans projet de maternité, le plus souvent suite à des rapports non ou mal protégés. Au Cameroun, selon l'EDS-MICS, \textbf{plus de 25 \% des adolescentes de 15--19 ans ont déjà commencé leur vie reproductive} (grossesse en cours ou déjà mère), avec de fortes disparités régionales (Extrême-Nord, Nord, Adamaoua les plus touchées).
\end{definition}

\courssub{Causes}

\begin{aretenir}[Causes principales des grossesses précoces]
\begin{itemize}
\item \textbf{Manque d'information} : Méconnaissance du cycle menstruel, de la période fertile, de la physiologie de la reproduction.
\item \textbf{Accès limité à la contraception} : Ruptures de stock, coût, distance des centres de santé, tabous, refus parental/partenaire.
\item \textbf{Facteurs socio-économiques et culturels} : Pauvreté (sexes transactionnels), mariages précoces (malgré la loi fixant l'âge à 18 ans), pression familiale, violences sexuelles, non-scolarisation ou déscolarisation précoce.
\item \textbf{Facteurs psychologiques} : Recherche d'affection, manque d'estime de soi, influence des pairs, consommation d'alcool/drogues favorisant les rapports à risque.
\end{itemize}
\end{aretenir}

\courssub{Conséquences}

\begin{propriete}[Conséquences médicales (Mère et Enfant)]
L'immaturité physiologique de l'adolescente crée une \textbf{compétition mère-fœtus} et un \textbf{déséquilibre fœto-pelvien}.
\begin{itemize}
\item \textbf{Pour la mère} : \cle{Dystocie} (bassin étroit $\rightarrow$ césarienne, fistules obstétricales vésico-vaginales/recto-vaginales), \cle{hémorragies} de la délivrance (utérus immature), \cle{anémie} sévère (carence martiale + besoins fœtaux), \cle{éclampsie} (HTA gravidique plus fréquente), \cle{mortalité maternelle} (risque $\times 2$ à $\times 4$ vs 20--24 ans). Avortements clandestins $\rightarrow$ septicémie, péritonite, stérilité, décès.
\item \textbf{Pour l'enfant} : \cle{Prématurité}, \cle{hypotrophie} (poids $< \SI{2500}{g}$), \cle{mortalité néonatale} et infantile accrues, bas score d'Apgar.
\end{itemize}
\end{propriete}

\begin{propriete}[Conséquences sociales et psychologiques]
\begin{itemize}
\item \textbf{Scolaires} : Déscolarisation définitive (exclusion, honte, garde d'enfant), perte de qualification $\rightarrow$ précarité économique durable.
\item \textbf{Familiales} : Rejet par les parents, rupture avec le partenaire, mariage forcé précoce, isolement.
\item \textbf{Psychologiques} : Dépression post-partum, anxiété, sentiment d'échec, trouble de l'attachement mère-enfant, risque de récidive.
\end{itemize}
\end{propriete}

\coursec{Infections sexuellement transmissibles (IST)}

\courssub{Classification et agents pathogènes}

\begin{propriete}[Principales IST : agents, transmission, complications]
\begin{center}
\begin{tabularx}{\linewidth}{|l|l|X|X|}\hline
\rowcolor{popblueL} \textbf{Maladie} & \textbf{Agent (Nature)} & \textbf{Transmission} & \textbf{Complications majeures si non traitées} \\ \hline
\cle{VIH / Sida} & Virus (VIH-1, VIH-2) & Sexuelle, sanguine, mère-enfant & Immunodépression, infections opportunistes, cancers, décès \\ \hline
\cle{Hépatite B} & Virus (VHB) & Sexuelle, sanguine, mère-enfant & Hépatite chronique, cirrhose, carcinome hépatocellulaire \\ \hline
\cle{Infection HPV} & Virus (Papillomavirus, HPV 16/18) & Sexuelle (muqueuses/cutanée) & Condylomes, \textbf{cancer du col de l'utérus}, cancers ORL/anaux \\ \hline
\cle{Herpès génital} & Virus (HSV-1, HSV-2) & Sexuelle (contact lésions/muqueuses) & Récurrences, néonatale grave, cofacteur VIH \\ \hline
\cle{Gonococcie} & Bactérie (\emph{Neisseria gonorrhoeae}) & Sexuelle & \textbf{Salpingite} (stérilité), arthrite, ophtalmie néonatale \\ \hline
\cle{Chlamydiose} & Bactérie (\emph{Chlamydia trachomatis}) & Sexuelle & \textbf{Salpingite silencieuse} (stérilité tubaire), GEU, pneumonie néonatale \\ \hline
\cle{Syphilis} & Bactérie (\emph{Treponema pallidum}) & Sexuelle, mère-enfant (congénitale) & Atteintes cutanées, neurologiques, cardiovasculaires, mort fœtale \\ \hline
\cle{Trichomonase} & Parasite (\emph{Trichomonas vaginalis}) & Sexuelle & Vaginite/urétrite, facteur de risque VIH, prématurité \\ \hline
\end{tabularx}
\end{center}
\end{propriete}

\courssub{Prévention des IST}

\begin{propriete}[Niveaux de prévention]
\begin{itemize}
\item \textbf{Primaire} : \cle{Abstinence} (100 \% efficace), \cle{Fidélité mutuelle} (après dépistage couple), \cle{Préservatif masculin/féminin} (seule méthode \textbf{double protection} IST + grossesse), \cle{Vaccination} (VHB : 3 doses nourrisson/adolescent ; HPV : 2 doses filles 9--14 ans, rattrapage possible).
\item \textbf{Secondaire} : \cle{Dépistage volontaire} (TROD VIH, sérologies VHB/VDRL, PCR Chlamydia/Gonocoque sur urines/prélèvement), \cle{Traitement précoce} (antibiotiques pour bactéries/parasite ; ARV pour VIH), \cle{Traitement des partenaires} (PTN -- Partenaires Traités Notification).
\item \textbf{Tertiaire} : Prise en charge des complications (chirurgie fistules, FIV pour stérilité, chimiothérapie cancers), soutien psychologique.
\end{itemize}
\end{propriete}

\begin{savaistu}[Riposte camerounaise aux IST/VIH/Sida]
Le \textbf{Comité National de Lutte contre le Sida (CNLS)} coordonne la riposte. La \cle{Prise en Charge Globale (PEC)} du VIH est \textbf{gratuite} (ARV, biologiques) dans les centres agréés (hôpitaux, CPC, centres de santé). La \textbf{vaccination VHB} est intégrée au PEV (Programme Élargi de Vaccination) depuis 2005 (doses à la naissance, 6, 10, 14 semaines). La \textbf{vaccination HPV} (Gardasil/Cervarix) est déployée depuis 2020 dans les écoles et centres de santé pour les filles de 9--14 ans (stratégie OMS élimination cancer col utérus). Le \textbf{Préservatif} est distribué gratuitement dans les formations sanitaires et via les ONG (CAMNAFAW, ACMS, etc.).
\end{savaistu}

\coursec{Méthodes de contraception et planification familiale}

\courssub{Classification et mécanismes d'action}

\begin{propriete}[Grandes catégories de méthodes contraceptives]
\begin{center}
\begin{tabularx}{\linewidth}{|l|X|X|X|X|}\hline
\rowcolor{popblueL} \textbf{Catégorie} & \textbf{Méthodes} & \textbf{Mécanisme principal} & \textbf{Efficacité (Indice de Pearl)} & \textbf{Protection IST} \\ \hline
\cle{Naturelles} & Ogino (calendrier), Température, Billings (glaire), Coït interrompu, Allaitement (LAM) & Éviter rapports en période fertile / Pas d'éjaculation intravaginale / Prolactine inhibe ovulation & Faible à moyenne (15--25) & \textbf{Non} \\ \hline
\cle{Barrière} & \textbf{Préservatif masculin} (latex/polyuréthane), \textbf{Préservatif féminin} (nitrile), Diaphragme, Cape cervicale & \textbf{Obstacle mécanique} : empêche rencontre spermatozoïde-ovule + \textbf{barrière agents pathogènes} & Moyenne à élevée (2--15 selon usage) & \textbf{OUI (seules)} \\ \hline
\cle{Hormonales} & Pilule \textbf{œstroprogestative} (quotidienne), Pilule \textbf{progestative} (quotidienne), \textbf{Injectable} (DMPA, \SI{3}{mois}), \textbf{Implant} (sous-cutané, 3--5 ans) & \textbf{Blocage ovulation} (inhibition pic LH), \textbf{Épaississement glaire} (infranchissable), \textbf{Atrophie endomètre} (anti-nidation) & \textbf{Très élevée} ($< 1$) & Non \\ \hline
\cle{Intra-utérines (DIU)} & \textbf{DIU Cuivre} (5--10 ans), \textbf{DIU Hormonal} (LNG, 5 ans) & Cuivre : \textbf{Spermicide} + réaction inflammatoire endomètre anti-nidation. Hormonal : Épaississement glaire + atrophie endomètre (+/- ovulation) & \textbf{Très élevée} ($< 1$) & Non \\ \hline
\cle{Définitives} & \textbf{Ligature trompes} (femme), \textbf{Vasectomie} (homme) & Stérilisation chirurgicale (section/occlusion trompes ou canaux déférents) & \textbf{Quasi 100 \%} & Non \\ \hline
\cle{Urgence} & \textbf{Pilule du lendemain} (Lévonorgestrel \SI{1,5}{mg} ou Ulipristal \SI{30}{mg}), \textbf{DIU Cuivre} (jusqu'à 5 j post-rapport) & Lévonorgestrel : \textbf{Retarde/bloque ovulation} (inefficace si ovulation faite). Ulipristal : Efficace plus près ovulation. DIU : Anti-nidation. & Variable (décroît avec temps) & Non \\ \hline
\end{tabularx}
\end{center}
\end{propriete}

\begin{attention}[Choix de la méthode : critères médicaux et contextuels]
\begin{itemize}
\item \textbf{Adolescent / Couple instable / Multi-partenaires} $\rightarrow$ \textbf{Préservatif} (double protection). Association possible avec pilule/implant (double méthode).
\item \textbf{Femme allaitante} $\rightarrow$ Pilule progestative (micro-pilule), Injectable, Implant, DIU (pas d'œstrogènes : pas d'effet sur lactation).
\item \textbf{Contre-indications aux œstrogènes} (HTA, tabagisme $> 35$ ans, antécédents thrombo-emboliques, migraine avec aura, obésité sévère) $\rightarrow$ Progestatifs seuls (pilule, implant, injectable, DIU hormonal) ou Cuivre.
\item \textbf{Désir de grossesse rapide après arrêt} $\rightarrow$ Éviter l'injectable (retour fertilité $\approx$ 6--10 mois). Préférer pilule, implant, DIU (retour immédiat).
\item \textbf{Le choix final relève d'une consultation médicale} (bilan, examen clinique, conseil) dans un \textbf{Centre de Planning Familial} ou formation sanitaire.
\end{itemize}
\end{attention}

\courssub{Planification familiale}

\begin{definition}[Planification familiale (PF)]
Ensemble de moyens permettant aux individus et couples d'\textbf{anticiper et d'atteindre le nombre d'enfants désiré}, ainsi que l'espacement et le moment des naissances. Elle repose sur : l'\textbf{information} (conseil pré-conceptionnel), l'\textbf{accès} aux méthodes contraceptives variées, la \textbf{prise en charge de l'infertilité}, et la \textbf{prévention des IST/VIH}. C'est un \textbf{droit fondamental} (Santé reproductive) et un levier de développement (dividende démographique, scolarisation filles, autonomie femmes).
\end{definition}

\begin{savaistu}[Indicateurs de PF au Cameroun (EDS 2018 / MICS)]
Taux de prévalence contraceptive (TPC) méthodes modernes : \textbf{$\approx$ 19 \%} (faible). Besoins non satisfaits : \textbf{$\approx$ 18 \%}. Objectif \textbf{PDSR 2020--2024} : TPC moderne $> 30 \%$. Les \textbf{Relais Communautaires} et les \textbf{Infirmiers de Santé Scolaire} sont des acteurs clés pour l'information et la référence des adolescents vers les services adaptés (accueil sans jugement, confidentialité, gratuité ciblée).
\end{savaistu}

\coursec{Démarches d'investigation et résolution de problèmes (Savoir-faire)}

\courssub{Méthode 1 : Expliquer les risques liés aux grossesses précoces à partir d'un cas}

\begin{methode}[Analyser un cas de grossesse précoce et expliquer ses risques]
Face à un document (cas clinique, tableau statistique, témoignage) portant sur une grossesse chez une adolescente, adopte la démarche suivante :
\begin{enumerate}
\item \textbf{Identifier le problème} : relever l'âge de l'adolescente (généralement avant 18 ans, souvent 12--17 ans), le contexte (scolarisation, milieu familial, situation économique) et les données chiffrées du document.
\item \textbf{Dégager les causes} à partir des informations du cas : rapports sexuels précoces et non protégés, méconnaissance de la physiologie de la reproduction (période fertile), absence de contraception, pressions sociales ou économiques, mariages précoces, abus sexuels.
\item \textbf{Expliquer les conséquences} en les classant rigoureusement en trois catégories :
\begin{itemize}
\item \emph{Conséquences médicales} : l'organisme de l'adolescente n'est pas complètement mature (bassin osseux étroit, croissance non achevée) $\Rightarrow$ accouchements dystociques, fistules obstétricales, hémorragies, anémies, éclampsie, mortalité maternelle et infantile élevées ; risque d'avortements clandestins (infections, stérilité, décès) ;
\item \emph{Conséquences sociales} : déscolarisation, rejet familial, précarité économique, isolement ;
\item \emph{Conséquences psychologiques} : stress, dépression, sentiment de culpabilité.
\end{itemize}
\item \textbf{Appuyer chaque affirmation sur une donnée} du document ou sur un savoir physiologique précis (ex. : « le bassin n'atteint sa taille définitive que vers 18--20 ans »).
\item \textbf{Conclure} par une phrase de synthèse reliant le problème à la nécessité de la prévention (information, contraception, responsabilisation).
\end{enumerate}
\textbf{Réflexe de rédaction} : ne jamais écrire « c'est dangereux » sans préciser \emph{pourquoi} (mécanisme ou donnée). Chaque risque cité doit être justifié.
\end{methode}

\begin{exoresolu}[Étude d'un cas au Centre Hospitalier d'Essos (Yaoundé)]
\textbf{Énoncé.} Dans un service de maternité d'un hôpital de Yaoundé, on a relevé sur une année les données suivantes :

\begin{center}
\begin{tabularx}{\linewidth}{|l|X|X|}\hline
\rowcolor{popblueL} Tranche d'âge des parturientes & Proportion de complications à l'accouchement & Proportion de nouveau-nés de faible poids ($< \SI{2,5}{kg}$) \\ \hline
13--17 ans & 38 \% & 31 \% \\ \hline
20--30 ans & 12 \% & 9 \% \\ \hline
\end{tabularx}
\end{center}

Par ailleurs, Nadège, 15 ans, élève en classe de Troisième, enceinte de 7 mois, a dû interrompre sa scolarité. Son poids est de \SI{44}{kg} et elle souffre d'anémie.
\begin{enumerate}
\item Définir l'expression « grossesse précoce ».
\item À partir des données du tableau, comparer les risques d'accouchement chez les adolescentes et chez les femmes de 20--30 ans.
\item Expliquer, à partir de la physiologie de l'adolescente, l'origine de deux complications observées.
\item Dégager deux conséquences sociales de la grossesse de Nadège.
\end{enumerate}
\tcblower
\textbf{Solution détaillée.}

\textbf{1. Définition.} Une \cle{grossesse précoce} (ou grossesse chez l'adolescente) est une grossesse survenant chez une fille âgée de moins de 18 ans, dont l'organisme n'a pas encore achevé sa croissance et sa maturation physiologique, en particulier celle de l'appareil génital et du bassin osseux. Elle est dite \cle{non désirée} lorsqu'elle survient en l'absence de tout projet de maternité, le plus souvent à la suite de rapports sexuels non protégés.

\textbf{2. Comparaison des risques.} Le tableau montre que les complications à l'accouchement touchent $38\ \%$ des adolescentes de 13--17 ans contre seulement $12\ \%$ des femmes de 20--30 ans, soit un risque environ \textbf{trois fois plus élevé} ($38/12 \approx 3,2$). De même, la proportion de nouveau-nés de faible poids de naissance est de $31\ \%$ chez les adolescentes contre $9\ \%$ chez les femmes de 20--30 ans, soit un risque environ $3,4$ fois plus élevé. On conclut que la grossesse précoce aggrave nettement les risques pour la mère \emph{et} pour l'enfant.

\textbf{3. Explication physiologique de deux complications.}
\begin{itemize}
\item \emph{Dystocie (accouchement difficile) et fistules obstétricales} : chez l'adolescente, la croissance du squelette n'est pas terminée ; le \cle{bassin osseux} est encore étroit et n'atteint ses dimensions définitives que vers 18--20 ans. Le passage du fœtus lors de l'accouchement est donc difficile : le travail prolongé comprime les tissus mous (vessie, vagin) contre les os du bassin, ce qui peut provoquer leur nécrose et l'apparition de fistules vésico-vaginales (fuites permanentes d'urine).
\item \emph{Anémie et faible poids du nouveau-né} : l'adolescente est elle-même en pleine croissance ; le fœtus entre en \emph{compétition nutritionnelle} avec elle pour le fer, les protéines et le calcium. Nadège, qui ne pèse que \SI{44}{kg}, ne dispose pas de réserves suffisantes : il en résulte une anémie maternelle (carence martiale) et une hypotrophie fœtale (nouveau-né de moins de \SI{2,5}{kg}), ce qui explique les $31\ \%$ de faibles poids de naissance du tableau.
\end{itemize}

\textbf{4. Conséquences sociales.} Nadège a dû \emph{interrompre sa scolarité} (déscolarisation), ce qui compromet durablement son insertion professionnelle et l'enferme dans la précarité économique. S'y ajoutent fréquemment le rejet par la famille ou le père de l'enfant, la stigmatisation et l'isolement social.

\textbf{Conclusion.} Les données chiffrées et le cas de Nadège montrent que la grossesse précoce constitue un problème majeur de santé publique et de développement : sa prévention passe par l'information des adolescent(e)s, l'accès aux méthodes de contraception et le maintien des filles à l'école.
\end{exoresolu}

\courssub{Méthode 2 : Justifier l'utilisation des moyens de prévention des IST}

\begin{methode}[Justifier un moyen de prévention des IST]
Pour justifier scientifiquement l'usage d'un moyen de prévention (préservatif, dépistage, fidélité, abstinence, vaccination) :
\begin{enumerate}
\item \textbf{Rappeler la nature de l'IST} concernée : nom de la maladie, agent pathogène précis (virus : VIH, VHB, HPV, HSV ; bactérie : \emph{Neisseria gonorrhoeae}, \emph{Chlamydia trachomatis}, \emph{Treponema pallidum} ; parasite : \emph{Trichomonas vaginalis}).
\item \textbf{Identifier la voie de transmission} : sécrétions génitales, sperme, sang, contact de muqueuses lors de rapports sexuels non protégés ; transmission mère-enfant possible pour certaines (VIH, syphilis, hépatite B).
\item \textbf{Expliquer le mécanisme d'action du moyen de prévention} : le préservatif constitue une \emph{barrière mécanique} imperméable aux agents pathogènes ; le dépistage permet un traitement précoce qui rompt la chaîne de transmission ; la vaccination (hépatite B, HPV) confère une immunité spécifique.
\item \textbf{Citer les conséquences évitées} : stérilité (salpingites à \emph{Chlamydia}), cancers (col de l'utérus lié au HPV), sida (évolution du VIH en l'absence d'antirétroviraux), complications de la grossesse.
\item \textbf{Conclure} en reliant le moyen de prévention au comportement responsable.
\end{enumerate}
\textbf{Réflexe} : distinguer toujours \emph{prévention} (éviter la contamination) et \emph{traitement} (guérir ou contrôler une infection déjà contractée). Le préservatif est le seul moyen qui protège à la fois des IST \emph{et} des grossesses non désirées.
\end{methode}

\begin{exoresolu}[Campagne de dépistage dans un lycée de Douala]
\textbf{Énoncé.} Lors d'une campagne de santé scolaire organisée dans un lycée de Douala, une infirmière explique aux élèves de Première D que le district de santé a enregistré, l'année précédente, 240 nouveaux cas d'infection à \emph{Chlamydia trachomatis} chez des jeunes de 15 à 24 ans, dont $70\ \%$ ne présentaient aucun symptôme. Elle recommande l'usage systématique du préservatif et le dépistage régulier.
\begin{enumerate}
\item Nommer l'agent de la chlamydiose et préciser sa nature.
\item Expliquer pourquoi l'absence de symptômes rend cette IST particulièrement dangereuse.
\item Justifier scientifiquement les deux moyens de prévention recommandés par l'infirmière.
\end{enumerate}
\tcblower
\textbf{Solution détaillée.}

\textbf{1. Agent pathogène.} La chlamydiose est causée par \emph{Chlamydia trachomatis}, une \cle{bactérie} (et non un virus) : c'est une bactérie intracellulaire obligatoire qui infecte les cellules épithéliales des voies génitales (urètre, col de l'utérus). Étant bactérienne, cette IST est \emph{curable} par antibiothérapie (ex. : azithromycine), à condition d'être dépistée.

\textbf{2. Danger de l'infection asymptomatique.} Le document indique que $70\ \%$ des 240 jeunes infectés, soit $240 \times 0{,}70 = 168$ jeunes, ne présentaient aucun symptôme. Or une personne infectée asymptomatique :
\begin{itemize}
\item \emph{ignore qu'elle est porteuse} : elle ne consulte pas, n'est pas traitée et continue de transmettre la bactérie à ses partenaires lors de rapports non protégés, entretenant silencieusement l'épidémie ;
\item \emph{développe des complications tardives} : chez la jeune fille, la bactérie remonte vers les trompes de Fallope et provoque des \cle{salpingites} (inflammations des trompes) dont les cicatrices obstruent les trompes, cause majeure de \emph{stérilité} et de grossesses extra-utérines ; chez le jeune homme, des épididymites peuvent également compromettre la fertilité.
\end{itemize}
L'absence de symptômes retarde donc le diagnostic jusqu'à des lésions irréversibles.

\textbf{3. Justification des moyens de prévention.}
\begin{itemize}
\item \emph{Le préservatif} (masculin ou féminin) : il constitue une \textbf{barrière mécanique} en latex ou polyuréthane, imperméable aux bactéries et aux virus, qui empêche le contact entre les sécrétions génitales et les muqueuses des partenaires. Correctement utilisé à chaque rapport, il bloque la transmission de \emph{Chlamydia trachomatis}, mais aussi du VIH, de la gonococcie et de la syphilis. C'est le seul moyen assurant simultanément la prévention des IST et des grossesses non désirées.
\item \emph{Le dépistage régulier} : puisque $70\ \%$ des infections sont silencieuses, seul un test de dépistage (prélèvement urinaire ou cervical) permet de repérer les porteurs asymptomatiques. Le dépistage conduit au \emph{traitement antibiotique précoce} du patient et de son ou ses partenaires, ce qui guérit l'infection avant les complications (stérilité) et \textbf{rompt la chaîne de transmission} dans la population.
\end{itemize}

\textbf{Conclusion.} L'usage systématique du préservatif et le dépistage régulier sont deux stratégies complémentaires et scientifiquement fondées : l'une empêche la contamination, l'autre détecte et élimine les infections silencieuses. Les adopter relève d'un comportement responsable en matière de santé reproductive.
\end{exoresolu}

\courssub{Méthode 3 : Identifier la période fertile et raisonner sur le cycle ovarien}

\begin{methode}[Exploiter un cycle menstruel pour déterminer la période fertile]
\begin{enumerate}
\item \textbf{Repérer la durée du cycle} (en jours, du premier jour des règles au premier jour des règles suivantes ; durée moyenne : 28 jours, mais variable de 21 à 35 jours).
\item \textbf{Localiser l'ovulation} : elle se produit environ \textbf{14 jours avant les règles suivantes} (la phase lutéale, post-ovulatoire, est relativement constante : 14 jours). Pour un cycle de 28 jours : ovulation vers le 14\up{e} jour ; pour un cycle de 30 jours : vers le 16\up{e} jour ($30 - 14 = 16$).
\item \textbf{Délimiter la période fertile} : l'ovule n'est fécondable que 12 à 24 h, mais les spermatozoïdes survivent 3 à 5 jours dans les voies génitales féminines. La période fertile s'étend donc d'environ \textbf{5 jours avant l'ovulation à 1 jour après}.
\item \textbf{Relier aux hormones} : pic de LH déclenchant l'ovulation ; œstrogènes (phase folliculaire) puis progestérone (phase lutéale) agissant sur l'endomètre.
\item \textbf{Conclure prudemment} : la méthode du calendrier (Ogino) est \emph{très peu fiable} car la durée du cycle varie ; elle ne protège en outre d'aucune IST.
\end{enumerate}
\textbf{Formule à retenir} : jour de l'ovulation $\approx$ durée du cycle $-$ 14.
\end{methode}

\begin{exoresolu}[Calcul de la période fertile à partir de données de cycles]
\textbf{Énoncé.} Une jeune femme de 22 ans, mariée, note attentivement les dates de début de ses règles pendant trois cycles consécutifs : 3 janvier, 31 janvier et 2 mars. Elle souhaite comprendre sa période fertile.
\begin{enumerate}
\item Calculer la durée moyenne de son cycle.
\item Estimer la date d'ovulation du cycle commençant le 2 mars.
\item Délimiter la période fertile de ce cycle et justifier les bornes choisies.
\item Expliquer pourquoi la méthode du calendrier seule est déconseillée pour éviter une grossesse.
\end{enumerate}
\tcblower
\textbf{Solution détaillée.}

\textbf{1. Durée moyenne du cycle.} Le cycle est l'intervalle entre le premier jour de deux menstruations consécutives.
\begin{align*}
\text{Cycle 1} &: \text{3 janvier} \rightarrow \text{31 janvier} = 28 \text{ jours} \\
\text{Cycle 2} &: \text{31 janvier} \rightarrow \text{2 mars} = 30 \text{ jours} \quad (\text{anée non bissextile : 28 j fév})
\end{align*}
Durée moyenne : $\dfrac{28 + 30}{2} = 29$ jours.

\textbf{2. Date estimée de l'ovulation.} La phase lutéale (entre l'ovulation et les règles suivantes) est remarquablement constante et dure environ 14 jours, car elle correspond à la durée de vie du corps jaune. Pour un cycle de 29 jours :
\[
\text{jour de l'ovulation} \approx 29 - 14 = 15\textsuperscript{e} \text{ jour du cycle.}
\]
Le cycle commençant le 2 mars (jour 1), le 15\textsuperscript{e} jour correspond au \textbf{16 mars}. L'ovulation est donc estimée autour du 16 mars.

\textbf{3. Période fertile.} Deux durées de vie doivent être combinées :
\begin{itemize}
\item les \cle{spermatozoïdes} conservent leur pouvoir fécondant pendant 3 à 5 jours dans les voies génitales féminines (glaire cervicale favorable en période ovulatoire) ;
\item l'\cle{ovule} n'est fécondable que pendant 12 à 24 heures après son expulsion de l'ovaire.
\end{itemize}
Un rapport sexuel survenant jusqu'à 5 jours \emph{avant} l'ovulation peut donc aboutir à une fécondation, et un rapport le lendemain de l'ovulation également. La période fertile s'étend ainsi approximativement du \textbf{11 mars au 17 mars} (soit du 10\textsuperscript{e} au 16\textsuperscript{e} jour du cycle), avec un maximum de probabilité les 16 et 17 mars.

\textbf{4. Limites de la méthode du calendrier (méthode Ogino).} Cette méthode repose sur l'hypothèse d'un cycle régulier, or :
\begin{itemize}
\item la durée du cycle varie naturellement (ici déjà 28 puis 30 jours) sous l'influence du stress, de la fatigue, des maladies ou du voyage ; un décalage de l'ovulation de quelques jours rend le calcul faux ;
\item l'ovulation peut être imprévisible, voire survenir deux fois dans un même cycle.
\end{itemize}
Son taux d'échec est élevé (environ 20 à 25 grossesses pour 100 femmes sur un an d'utilisation) ; de plus, elle \emph{ne protège d'aucune IST}. Elle ne doit donc pas être utilisée seule comme méthode de contraception : on lui préfère des méthodes fiables (préservatif, pilule, implant, DIU).

\textbf{Conclusion.} La période fertile estimée s'étend du 11 au 17 mars, mais la variabilité naturelle du cycle rend la méthode du calendrier insuffisamment fiable pour prévenir une grossesse.
\end{exoresolu}

\courssub{Méthode 4 : Choisir et justifier une méthode de contraception adaptée à une situation}

\begin{methode}[Conseiller une méthode de contraception à partir d'un cas]
\begin{enumerate}
\item \textbf{Analyser la situation} : âge, statut (adolescent, couple stable, jeune mère qui allaite), objectif (contraception temporaire ou définitive), besoin de protection contre les IST.
\item \textbf{Passer en revue les grandes catégories de méthodes} :
\begin{itemize}
\item \emph{méthodes de barrière} : préservatifs masculin et féminin (seules à protéger aussi des IST) ;
\item \emph{méthodes hormonales} : pilule œstroprogestative ou progestative, injectables (ex. : DMPA, injection trimestrielle), implant sous-cutané (3 à 5 ans) ;
\item \emph{dispositif intra-utérin} (DIU ou « stérilet ») : au cuivre ou hormonal, efficace 5 à 10 ans ;
\item \emph{méthodes définitives} : ligature des trompes, vasectomie ;
\item \emph{contraception d'urgence} : pilule du lendemain, à prendre le plus tôt possible (dans les 72 h) après un rapport non protégé.
\end{itemize}
\item \textbf{Associer chaque méthode à son mécanisme d'action} : blocage de l'ovulation (hormones), épaississement de la glaire cervicale, réaction du corps étranger et altération de l'endomètre (DIU), obstacle mécanique (préservatif).
\item \textbf{Justifier le choix} en croisant la situation et les propriétés de la méthode (efficacité, durée, protection IST, réversibilité, contre-indications).
\item \textbf{Rappeler} que le choix définitif se fait avec un personnel de santé qualifié (centre de santé, planning familial).
\end{enumerate}
\textbf{Réflexe} : pour un adolescent ou un couple non stable, la réponse attendue est presque toujours le \textbf{préservatif}, car il répond au « double risque » (grossesse + IST).
\end{methode}

\begin{exoresolu}[Conseil contraceptif dans trois situations de la vie courante]
\textbf{Énoncé.} Un centre de planning familial de Bafoussam reçoit trois personnes :
\begin{itemize}
\item \textbf{Cas A} : Brice, 17 ans, élève de Première, qui débute sa vie sexuelle avec une partenaire dont il ne connaît pas le statut sérologique ;
\item \textbf{Cas B} : Mme Etoundi, 32 ans, mère de trois enfants, qui souhaite espacer ses naissances d'au moins 5 ans sans prise quotidienne de médicament ;
\item \textbf{Cas C} : Sandrine, 19 ans, dont le préservatif s'est déchiré la veille lors d'un rapport sexuel.
\end{itemize}
Pour chaque cas, proposer la méthode contraceptive la plus adaptée et la justifier scientifiquement en précisant son mécanisme d'action.
\tcblower
\textbf{Solution détaillée.}

\textbf{Cas A — Brice : le préservatif masculin.} Brice est exposé au \emph{double risque} caractéristique de l'adolescence : grossesse non désirée \emph{et} IST (il ignore le statut sérologique de sa partenaire). Le \cle{préservatif} est la seule méthode qui réponde aux deux risques simultanément. \emph{Mécanisme} : barrière mécanique en latex qui retient les spermatozoïdes et empêche tout contact entre les sécrétions génitales et les muqueuses, bloquant ainsi la transmission des agents pathogènes (VIH, gonocoques, chlamydiae...). Conditions d'efficacité : utilisation correcte et systématique à chaque rapport, vérification de la date de péremption et de l'intégrité de l'emballage.

\textbf{Cas B — Mme Etoundi : l'implant sous-cutané ou le DIU.} Elle souhaite unecontraception \emph{de longue durée} (au moins 5 ans), \emph{réversible} et sans contrainte quotidienne. Deux réponses sont recevables :
\begin{itemize}
\item l'\cle{implant} (bâtonnet progestatif inséré sous la peau du bras, efficace 3 à 5 ans) : il bloque l'\emph{ovulation} par inhibition de la sécrétion des gonadotrophines hypophysaires (FSH et LH) et épaissit la glaire cervicale, qui devient infranchissable pour les spermatozoïdes ;
\item le \cle{DIU au cuivre} (efficace jusqu'à 10 ans) : le cuivre exerce un effet spermicide et le dispositif provoque une réaction locale de l'endomètre empêchant la nidation.
\end{itemize}
Ces méthodes ne protègent pas des IST, ce qui est acceptable ici car il s'agit d'un couple stable et fidèle. L'implant est particulièrement adapté à son souhait de durée de 5 ans.

\textbf{Cas C — Sandrine : la contraception d'urgence.} L'échec du préservatif expose Sandrine à un risque de grossesse non désirée (et d'IST : un dépistage sera conseillé). Elle doit prendre au plus vite la \cle{pilule du lendemain} (lévonorgestrel), idéalement dans les 24 h et au plus tard dans les \textbf{72 heures} suivant le rapport : son efficacité diminue avec le temps. \emph{Mécanisme} : forte dose de progestatif qui \emph{retarde ou bloque l'ovulation} et modifie la glaire cervicale ; elle n'est efficace que si l'ovulation n'a pas encore eu lieu. Attention : c'est une méthode de \emph{rattrapage exceptionnel}, non une contraception régulière, et elle ne protège pas des IST.

\textbf{Conclusion.} Il n'existe pas de méthode contraceptive universelle : le choix dépend de l'âge, de la situation de couple, de la durée souhaitée et du risque d'IST. Le recours à un personnel de santé qualifié garantit un choix sûr et adapté.
\end{exoresolu}

\courssub{Méthode 5 : Construire un argumentaire de promotion d'un comportement responsable}

\begin{methode}[Adopter et promouvoir un comportement responsable en santé reproductive]
Pour rédiger un texte argumentatif, un message de sensibilisation ou répondre à une question ouverte sur les comportements responsables :
\begin{enumerate}
\item \textbf{Définir le comportement responsable} : comportement fondé sur l'information exacte, le respect de soi et d'autrui, l'anticipation des conséquences de ses actes et le recours aux services de santé.
\item \textbf{Structurer l'argumentation autour des piliers} :
\begin{itemize}
\item \emph{s'informer} : connaître la physiologie de la reproduction, les IST, les méthodes de contraception (sources fiables : centres de santé, personnels de santé, et non les rumeurs) ;
\item \emph{se protéger} : abstinence ou rapports protégés par le préservatif, contraception adaptée, vaccination (hépatite B, HPV) ;
\item \emph{se faire dépister} : dépistage volontaire du VIH et des IST, en couple avant tout rapport non protégé ;
\item \emph{respecter} : consentement mutuel, refus des relations forcées et des mariages précoces, fidélité réciproque dans le couple ;
\item \emph{consulter} : centres de planning familial, consultations prénatales précoces en cas de grossesse.
\end{itemize}
\item \textbf{Appuyer chaque argument sur un fait scientifique ou une donnée} (ex. : « le préservatif correctement utilisé réduit fortement le risque de transmission du VIH »).
\item \textbf{Adapter le message au public visé} (pairs adolescents, parents, communauté) et proposer des actions concrètes (clubs de santé scolaire, causeries éducatives, campagnes de dépistage).
\item \textbf{Conclure} par un appel à l'action clair et positif.
\end{enumerate}
\end{methode}

\begin{exoresolu}[Rédaction d'un message de sensibilisation pour une causerie éducative]
\textbf{Énoncé.} Dans le cadre de la semaine de la santé de ton lycée à Ngaoundéré, le club « Santé et Avenir » te charge de présenter aux élèves de Seconde un exposé de dix lignes intitulé : « Protéger sa santé reproductive, c'est préparer son avenir ». Rédige cet exposé en mobilisant au moins quatre arguments scientifiquement fondés et en proposant deux actions concrètes.
\tcblower
\textbf{Solution détaillée (texte rédigé).}

« Chers camarades, protéger sa santé reproductive, c'est préparer son avenir.

D'abord, \emph{connaître son corps} est une force : comprendre le cycle menstruel, la période fertile et le fonctionnement de la reproduction permet de prendre des décisions éclairées au lieu de subir les conséquences de l'ignorance.

Ensuite, les \emph{grossesses précoces} mettent en danger la vie des adolescentes : avant 18 ans, le bassin n'est pas complètement formé, ce qui multiplie par trois environ les complications de l'accouchement, et la déscolarisation qui s'ensuit ferme les portes de l'avenir.

De plus, les \emph{IST} comme la chlamydiose évoluent silencieusement — 7 infections sur 10 sans symptôme — et peuvent rendre stérile ; le VIH, lui, reste une infection grave qui exige un traitement à vie. Or le \emph{préservatif}, correctement utilisé à chaque rapport, protège à la fois des IST et des grossesses non désirées : c'est le seul moyen qui offre cette double protection.

Enfin, un comportement responsable, c'est aussi le \emph{respect} : le consentement mutuel, le refus des relations forcées, et le dépistage volontaire en couple avant tout engagement.

Agissons concrètement : \textbf{(1)} participons chaque trimestre à une causerie animée par l'infirmier du centre de santé de quartier, et \textbf{(2)} organisons avec le club une journée de dépistage volontaire et gratuit du VIH dans notre lycée.

Notre santé est notre premier capital : protégeons-la aujourd'hui pour construire demain. »

\textbf{Commentaire de correction.} Le texte respecte la consigne : quatre arguments fondés sur des faits scientifiques précis (physiologie du bassin, statistique des complications, infections asymptomatiques, double protection du préservatif), une dimension de savoir-être (consentement, respect), deux actions concrètes et réalisables dans le contexte scolaire camerounais, et une conclusion mobilisatrice. Chaque affirmation reste rigoureuse : aucune exagération, aucun chiffre inventé.
\end{exoresolu}

\courssub{Méthode 6 : Exploiter un document statistique ou un schéma physiologique}

\begin{methode}[Lire et interpréter un document en santé reproductive]
\begin{enumerate}
\item \textbf{Lire le titre et les légendes} : identifier la grandeur étudiée, les unités, les groupes comparés, la date et le lieu de l'étude.
\item \textbf{Dégager les faits} : relever les valeurs extrêmes, calculer les écarts ou les rapports (ex. : « le taux est multiplié par $3$ »), décrire les tendances (augmentation, diminution) \emph{sans encore interpréter}.
\item \textbf{Comparer} les groupes ou les périodes de façon chiffrée : toujours citer les valeurs exactes du document.
\item \textbf{Interpréter} en mobilisant les savoirs du cours : relier les faits observés à des mécanismes physiologiques ou à des comportements (ex. : baisse des grossesses précoces $\leftrightarrow$ campagnes d'information et accès à la contraception).
\item \textbf{Conclure} en répondant exactement à la question posée, puis éventuellement proposer une mesure de prévention.
\end{enumerate}
\textbf{Réflexe de rédaction} : bien séparer dans la copie « ce que montre le document » (constat) et « ce que j'en déduis » (interprétation). Le correcteur sanctionne la confusion entre les deux.
\end{methode}

\begin{exoresolu}[Analyse de données de santé scolaire dans la région du Centre]
\textbf{Énoncé.} Le tableau ci-dessous présente l'évolution du nombre de grossesses recensées parmi les élèves de trois lycées d'une ville du Cameroun, avant et après la mise en place d'un programme d'éducation à la santé reproductive (causeries trimestrielles, distribution d'information, relais-élèves).

\begin{center}
\begin{tabularx}{\linewidth}{|l|X|X|X|}\hline
\rowcolor{popblueL} Année scolaire & Lycée A (programme actif) & Lycée B (programme actif) & Lycée C (sans programme) \\ \hline
Avant programme & 24 & 19 & 21 \\ \hline
Deux ans après & 9 & 7 & 22 \\ \hline
\end{tabularx}
\end{center}

\begin{enumerate}
\item Calculer, pour chaque lycée, la variation en pourcentage du nombre de grossesses entre les deux périodes.
\item Comparer l'évolution des lycées A et B à celle du lycée C.
\item Interpréter ces résultats et dégager le rôle de l'éducation à la santé reproductive.
\end{enumerate}
\tcblower
\textbf{Solution détaillée.}

\textbf{1. Calcul des variations en pourcentage.} La variation relative se calcule par :
\[
\text{variation} = \frac{\text{valeur finale} - \text{valeur initiale}}{\text{valeur initiale}} \times 100.
\]
\begin{align*}
\text{Lycée A} &: \frac{9 - 24}{24} \times 100 = \frac{-15}{24} \times 100 = -62{,}5\ \% \\
\text{Lycée B} &: \frac{7 - 19}{19} \times 100 = \frac{-12}{19} \times 100 \approx -63{,}2\ \% \\
\text{Lycée C} &: \frac{22 - 21}{21} \times 100 = \frac{1}{21} \times 100 \approx +4{,}8\ \%
\end{align*}

\textbf{2. Comparaison.} Dans les lycées A et B, dotés du programme, le nombre de grossesses a \emph{chuté d'environ 62 à 63 \%} en deux ans (division par presque trois : de 24 à 9, et de 19 à 7). En revanche, dans le lycée C, resté sans programme, ce nombre est \emph{resté stable} ($+4{,}8\ \%$, variation négligeable). Les trois lycées partaient pourtant de niveaux comparables (19 à 24 grossesses), ce qui rend la comparaison valide.

\textbf{3. Interprétation.} La stabilité observée au lycée C montre que la baisse enregistrée dans les lycées A et B n'est pas due au hasard ni à une évolution générale de la ville, mais bien au \emph{programme d'éducation à la santé reproductive}. Mécanismes en jeu :
\begin{itemize}
\item les causeries ont levé les idées fausses (méconnaissance de la période fertile, croyances sur la « sécurité » du premier rapport) ;
\item l'information sur les méthodes de contraception et le rôle des relais-élèves ont favorisé l'adoption de comportements protecteurs (abstinence, usage du préservatif) ;
\item la sensibilisation aux conséquences des grossesses précoces (déscolarisation, risques médicaux) a renforcé la motivation à les éviter.
\end{itemize}

\textbf{Conclusion.} La comparaison chiffrée démontre que l'éducation à la santé reproductive est un outil de prévention efficace : informée et accompagnée, la jeunesse adopte des comportements responsables qui réduisent fortement les grossesses précoces. Généraliser de tels programmes à tous les établissements scolaires constitue donc une mesure de santé publique justifiée.
\end{exoresolu}

\begin{attention}[Les 5 erreurs qui coûtent des points au Probatoire]
\begin{enumerate}
\item \textbf{Confondre contraception et prévention des IST.} La pilule, l'implant et le DIU empêchent la grossesse mais \emph{ne protègent d'aucune IST} ; seul le préservatif assure la double protection. Écrire « la pilule protège du sida » est une erreur éliminatoire.
\item \textbf{Citer un risque sans l'expliquer.} Affirmer « la grossesse précoce est dangereuse » sans invoquer le bassin immature, la compétition nutritionnelle mère-fœtus ou les données du document ne rapporte qu'une fraction des points : chaque risque doit être \emph{justifié} par un mécanisme ou un chiffre.
\item \textbf{Se tromper dans le calcul de l'ovulation.} L'ovulation se situe 14 jours \emph{avant les règles suivantes} (jour $=$ durée du cycle $-14$), et non forcément le 14\up{e} jour du cycle : pour un cycle de 30 jours, c'est le 16\up{e} jour. Oublier la durée de survie des spermatozoïdes (3 à 5 jours) fausse aussi la délimitation de la période fertile.
\item \textbf{Nommer imprécisément les agents pathogènes.} Écrire « le microbe du sida » ou confondre virus et bactérie (la chlamydiose et la gonococcie sont \emph{bactériennes}, donc curables par antibiotiques ; le VIH est un \emph{virus}, non curable mais contrôlable par antirétroviraux). Le nom scientifique correct (\emph{Chlamydia trachomatis}, \emph{Treponema pallidum}) est valorisé.
\item \textbf{Décrire un document sans le chiffrer, ou interpréter sans décrire.} Dans l'exploitation d'un tableau, il faut d'abord citer les valeurs et calculer les écarts (« 38 \% contre 12 \%, soit un risque triple »), \emph{puis} interpréter. Mélanger constat et interprétation, ou tirer une conclusion non étayée par les données, fait perdre des points de méthode.
\end{enumerate}
\end{attention}

\newpage

\coursec{Exercices}

\courssub{Je vérifie mes connaissances}

\begin{exercice}[QCM : la puberté et la reproduction]
Pour chaque question, une seule réponse est correcte. Recopie la lettre correspondante.
\begin{enumerate}
\item La puberté chez la fille se manifeste notamment par :
\begin{enumerate}
\item[a)] l'apparition de la mue de la voix ;
\item[b)] l'apparition des premières règles (ménarche) ;
\item[c)] la chute des cheveux ;
\item[d)] l'arrêt de la croissance.
\end{enumerate}
\item L'ovulation a lieu en moyenne :
\begin{enumerate}
\item[a)] le 1\textsuperscript{er} jour du cycle ;
\item[b)] le 7\textsuperscript{e} jour du cycle ;
\item[c)] le 14\textsuperscript{e} jour d'un cycle de 28 jours ;
\item[d)] le 28\textsuperscript{e} jour du cycle.
\end{enumerate}
\item Le spermatozoïde et l'ovule se rencontrent normalement dans :
\begin{enumerate}
\item[a)] l'utérus ; \item[b)] le vagin ; \item[c)] la trompe de Fallope ; \item[d)] l'ovaire.
\end{enumerate}
\item L'hormone responsable de l'apparition des caractères sexuels secondaires chez le garçon est :
\begin{enumerate}
\item[a)] la progestérone ; \item[b)] la testostérone ; \item[c)] l'insuline ; \item[d)] la FSH.
\end{enumerate}
\end{enumerate}
\end{exercice}

\begin{exercice}[Vrai ou faux justifié]
Indique si chaque affirmation est vraie ou fausse, puis justifie ta réponse en une ou deux phrases.
\begin{enumerate}
\item « Une jeune fille ne peut pas tomber enceinte dès son premier rapport sexuel. »
\item « Le préservatif protège à la fois contre les grossesses non désirées et contre les IST. »
\item « La pilule contraceptive protège contre le VIH/Sida. »
\item « Le VIH se transmet par la piqûre de moustique. »
\item « Une IST non traitée peut entraîner la stérilité. »
\end{enumerate}
\end{exercice}

\begin{exercice}[Définitions]
Définis brièvement et précisément chacun des termes suivants : \cle{puberté}, \cle{grossesse précoce}, \cle{infection sexuellement transmissible (IST)}, \cle{contraception}, \cle{planification familiale}.
\end{exercice}

\begin{exercice}[Lecture d'un graphique du cycle menstruel]
Le graphique ci-dessous représente l'évolution de la concentration sanguine de deux hormones (LH et \oe{}strogènes) au cours d'un cycle menstruel de 28 jours chez une jeune fille de 17 ans.

\begin{popfigure}
\begin{tikzpicture}
\begin{axis}[width=0.9\linewidth,height=6cm,
xlabel={Jour du cycle},ylabel={Concentration (unités arbitraires)},
xmin=0,xmax=28,ymin=0,ymax=10,
xtick={0,7,14,21,28},ytick={0,2,4,6,8,10},
legend style={at={(0.02,0.98)},anchor=north west,font=\small,draw=popline,fill=popcream},
grid=major,grid style={popline,opacity=0.3}]
\addplot[popcurve,popblue,domain=0:28,samples=200]{2+6*exp(-((x-13.5)*(x-13.5))/2)};
\addlegendentry{LH}
\addplot[popcurve,poporange,domain=0:28,samples=200]{3+3.5*exp(-((x-12)*(x-12))/8)+2*exp(-((x-23)*(x-23))/18)};
\addlegendentry{\OE{}strogènes}
\end{axis}
\end{tikzpicture}
\legende{Variation des taux sanguins de LH (bleu) et d'œstrogènes (orange) au cours d'un cycle menstruel standard de 28 jours. Le pic de LH déclenche l'ovulation vers le 14\textsuperscript{e} jour.}
\end{popfigure}

\begin{enumerate}
\item À quel jour du cycle observe-t-on le pic de LH ? Que déclenche ce pic ?
\item En déduire la date approximative de l'ovulation.
\item Sachant que l'ovule reste fécondable environ 24 h et que les spermatozoïdes survivent 3 à 5 jours dans les voies génitales féminines, détermine la période fertile de ce cycle.
\item Explique pourquoi un rapport non protégé au 13\textsuperscript{e} jour présente un risque élevé de grossesse.
\end{enumerate}
\end{exercice}

\courssub{Je m'entraîne}

\begin{exercice}[Tableau des IST à compléter]
Recopie et complète le tableau suivant en associant à chaque IST son agent pathogène, un symptôme caractéristique et un moyen de prévention.

\begin{center}
\begin{tabularx}{\linewidth}{|l|X|X|X|}
\hline
\rowcolor{popblueL} \textbf{IST} & \textbf{Agent pathogène} & \textbf{Symptôme caractéristique} & \textbf{Prévention} \\
\hline
VIH/Sida & \dots & \dots & \dots \\
\hline
Gonococcie & \dots & \dots & \dots \\
\hline
Syphilis & \dots & \dots & \dots \\
\hline
HPV (papillomavirus) & \dots & \dots & \dots \\
\hline
\end{tabularx}
\end{center}
\end{exercice}

\begin{exercice}[La double protection]
Un élève de Première affirme : « Ma copine prend la pilule, donc nous sommes protégés contre tout. »
\begin{enumerate}
\item Explique scientifiquement pourquoi le préservatif est qualifié de moyen de \cle{double protection}.
\item Montre, en t'appuyant sur le mode d'action de la pilule, pourquoi cette seule méthode ne suffit pas pour un couple d'adolescents.
\item Rédige une phrase de conseil que tu pourrais adresser à cet élève.
\end{enumerate}
\end{exercice}

\begin{exercice}[Classer les méthodes de contraception]
On donne six méthodes de contraception : \textbf{calendrier (méthode Ogino), préservatif, pilule, stérilet (cuivre), spermicide, implant}.
\begin{enumerate}
\item Classe ces méthodes selon leur mode d'action : méthodes \cle{naturelles}, méthodes \cle{mécaniques (barrières)}, méthodes \cle{chimiques}, méthodes \cle{hormonales}.
\item Pour chaque méthode, indique en une phrase son principe de fonctionnement.
\item En t'aidant des indices de fiabilité (indice de Pearl) vus en cours, range ces méthodes de la moins efficace à la plus efficace en usage courant.
\end{enumerate}
\end{exercice}

\begin{exercice}[Sensibilisation à la planification familiale]
Le club santé de ton lycée organise une journée de sensibilisation à destination des élèves de Seconde. Rédige un court texte argumentatif (10 à 15 lignes) expliquant l'importance de la planification familiale et du comportement responsable en matière de santé reproductive. Ton texte devra comporter au moins deux arguments scientifiques et un appel à l'action.
\end{exercice}

\courssub{Je me prépare au Probatoire}

\begin{exercice}[Étude de cas : une grossesse précoce à Bafou]
Dans un village de la région de l'Ouest, Amina, 15 ans, élève en classe de Troisième, découvre qu'elle est enceinte de trois mois. Son ami, âgé de 17 ans, nie toute responsabilité. Les parents d'Amina, très pauvres, envisagent de la marier. Le médecin du centre de santé intégré signale que le bassin d'Amina n'est pas encore complètement développé.
\begin{enumerate}
\item Définis l'expression « grossesse précoce » et montre que le cas d'Amina y correspond.
\item Cite deux causes possibles de cette grossesse précoce.
\item Décris trois conséquences de cette grossesse : une conséquence sanitaire, une conséquence scolaire et une conséquence sociale.
\item Explique, à partir de tes connaissances en physiologie, pourquoi l'accouchement présente un risque accru chez Amina.
\item Propose deux mesures de prévention des grossesses précoces qui pourraient être mises en \oe{}uvre dans ce village.
\end{enumerate}
\end{exercice}

\begin{exercice}[Données épidémiologiques et prévention des IST]
Le tableau ci-dessous donne le nombre de nouveaux cas d'IST déclarés chez les 15--24 ans dans un district de santé du Cameroun entre 2020 et 2023.

\begin{center}
\begin{tabularx}{\linewidth}{|l|X|X|X|X|}
\hline
\rowcolor{popblueL} \textbf{Année} & \textbf{2020} & \textbf{2021} & \textbf{2022} & \textbf{2023} \\
\hline
Nouveaux cas d'IST & 120 & 145 & 190 & 260 \\
\hline
\end{tabularx}
\end{center}

\begin{enumerate}
\item Représente ces données par un diagramme en bâtons (1 cm pour 20 cas en ordonnée ; 1 cm pour une année en abscisse).
\item Calcule le pourcentage d'augmentation du nombre de cas entre 2020 et 2023.
\item Décris l'évolution observée et propose deux explications possibles.
\item Cite trois modes de transmission des IST et, pour chacun, le moyen de prévention correspondant.
\item Justifie l'intérêt du dépistage précoce et gratuit dans les centres de santé pour les adolescents.
\end{enumerate}
\end{exercice}

\begin{exercice}[Sujet de synthèse : santé reproductive et comportement responsable]
\textbf{Document 1.} Au Cameroun, environ 25 \% des adolescentes de 15 à 19 ans ont déjà commencé une vie de mère (enquêtes démographiques et de santé).

\textbf{Document 2.} Une enquête réalisée dans un lycée de Yaoundé montre que 60 \% des élèves de Première connaissent le préservatif, mais que seulement 20 \% savent citer une autre méthode de contraception.

\textbf{Document 3.} Le Ministère de la Santé publique mène des campagnes de dépistage gratuit du VIH et distribue des préservatifs dans les établissements scolaires partenaires.

\begin{enumerate}
\item À partir du document 1, dégage le problème de santé publique posé.
\item À partir du document 2, identifie le facteur qui favorise ce problème parmi les élèves.
\item Explique le rôle de chacune des actions décrites dans le document 3 dans la lutte contre les IST et les grossesses non désirées.
\item Compare le mode d'action de deux méthodes de contraception au choix (hormonale et mécanique) et précise leurs avantages respectifs pour un adolescent.
\item En t'appuyant sur l'ensemble des documents et sur tes connaissances, rédige un texte structuré (15 à 20 lignes) montrant comment un comportement responsable en matière de santé reproductive contribue à l'épanouissement de l'adolescent et au développement de la communauté.
\end{enumerate}
\end{exercice}

\newpage

\coursec{Corrigés des exercices}

\begin{corrige}[Exercice 1 — QCM : la puberté et la reproduction]
\textbf{1. Réponse b)} — Le \cle{ménarche} (apparition des premières règles) est le signe le plus caractéristique de la puberté chez la fille, vers 12--13 ans en moyenne. La mue de la voix (a) concerne le garçon ; la croissance staturale se poursuit encore après la puberté (d est donc faux).

\textbf{2. Réponse c)} — L'ovulation a lieu en moyenne le 14\textsuperscript{e} jour d'un cycle de 28 jours (plus exactement 14 jours \emph{avant} le début des règles suivantes). Le 1\textsuperscript{er} jour du cycle correspond au premier jour des règles.

\textbf{3. Réponse c)} — La fécondation a lieu normalement dans le tiers externe de la \cle{trompe utérine} (trompe de Fallope). L'ovocyte y est capté par les franges du pavillon après l'ovulation ; les spermatozoïdes y remontent depuis le vagin et l'utérus.

\textbf{4. Réponse b)} — La \cle{testostérone}, hormone stéroïde sécrétée par les testicules (cellules de Leydig, situées dans le tissu interstitiel) sous contrôle de la LH hypophysaire, est responsable de l'apparition des caractères sexuels secondaires masculins (mue de la voix, pilosité, développement musculaire). La progestérone est une hormone ovarienne, l'insuline régule la glycémie, et la FSH est une gonadostimuline hypophysaire qui agit sur les gonades mais n'est pas l'hormone des caractères sexuels.
\end{corrige}

\begin{corrige}[Exercice 2 — Vrai ou faux justifié]
\begin{enumerate}
\item \textbf{Faux.} Dès les premières règles, l'ovulation a commencé : une fille est fertile avant même son premier rapport. Un seul rapport non protégé pendant la période fertile peut donc entraîner une grossesse, même le tout premier.
\item \textbf{Vrai.} Le préservatif est une barrière mécanique qui empêche à la fois la rencontre entre spermatozoïdes et ovocyte (contraception) et le contact avec les sécrétions génitales et le sang porteurs d'agents pathogènes (prévention des IST, dont le VIH) : c'est la \cle{double protection}.
\item \textbf{Faux.} La pilule agit par voie hormonale (blocage de l'ovulation) : elle n'empêche aucunement le passage des virus ou des bactéries. Elle ne protège donc ni contre le VIH ni contre aucune autre IST.
\item \textbf{Faux.} Le VIH ne se transmet que par le sang, les sécrétions sexuelles et le lait maternel. Le moustique ne transmet pas le VIH (il transmet le paludisme) : le virus ne survit pas et ne se multiplie pas dans l'organisme de l'insecte, et la quantité de sang éventuellement présente sur sa trompe est infime.
\item \textbf{Vrai.} Une IST bactérienne non traitée (gonococcie, chlamydiose) peut remonter vers les trompes utérines (salpingite) ou les épididymes et provoquer des lésions définitives conduisant à la stérilité, chez la femme comme chez l'homme.
\end{enumerate}
\end{corrige}

\begin{corrige}[Exercice 3 — Définitions]
\begin{itemize}
\item \cle{Puberté} : période de transition entre l'enfance et l'âge adulte, marquée par la maturation des gonades, l'apparition des caractères sexuels secondaires et l'acquisition de la capacité de se reproduire, sous l'action des hormones sexuelles (testostérone chez le garçon, \oe{}strogènes et progestérone chez la fille).
\item \cle{Grossesse précoce} : grossesse survenant chez une adolescente (généralement entre 10 et 19 ans), dont l'organisme n'est pas encore complètement mature ; elle est dite non désirée lorsqu'elle n'est ni voulue ni planifiée.
\item \cle{Infection sexuellement transmissible (IST)} : infection causée par un agent pathogène (virus, bactérie, parasite) transmis principalement lors de rapports sexuels non protégés (ex. : VIH/Sida, gonococcie, syphilis, chlamydiose, papillomavirus HPV).
\item \cle{Contraception} : ensemble des méthodes utilisées pour éviter une grossesse de façon temporaire et réversible (préservatif, pilule, stérilet, implant, spermicides, méthodes naturelles).
\item \cle{Planification familiale} : démarche volontaire par laquelle un couple ou un individu choisit librement le nombre de ses enfants et l'espacement des naissances, grâce à l'information et aux méthodes de contraception.
\end{itemize}
\end{corrige}

\begin{corrige}[Exercice 4 — Lecture d'un graphique du cycle menstruel]
\begin{enumerate}
\item Le pic de LH est observé vers le \textbf{13\textsuperscript{e}--14\textsuperscript{e} jour} du cycle (courbe bleue, maximum autour du jour 13,5). Ce pic de LH déclenche l'\cle{ovulation}, c'est-à-dire la libération de l'ovocyte par l'ovaire (rupture du follicule de De Graaf).
\item L'ovulation suit le pic de LH de 24 à 36 h : elle a donc lieu approximativement le \textbf{14\textsuperscript{e} jour} du cycle.
\item Les spermatozoïdes survivent 3 à 5 jours dans les voies génitales féminines et l'ovocyte reste fécondable environ 24 h. La période fertile s'étend donc d'environ \textbf{5 jours avant l'ovulation à 1 jour après}, soit approximativement du \textbf{9\textsuperscript{e} au 15\textsuperscript{e} jour} du cycle.
\item Au 13\textsuperscript{e} jour, on se situe juste avant l'ovulation, au c\oe{}ur de la période fertile : les spermatozoïdes déposés lors du rapport survivent assez longtemps (3 à 5 jours) pour être encore présents lors de la libération de l'ovocyte le 14\textsuperscript{e} jour. La probabilité de fécondation est donc maximale : un rapport non protégé à cette date présente un risque élevé de grossesse.
\end{enumerate}
\end{corrige}

\begin{corrige}[Exercice 5 — Tableau des IST à compléter]
\begin{center}
\begin{tabularx}{\linewidth}{|l|X|X|X|}
\hline
\rowcolor{popblueL} \textbf{IST} & \textbf{Agent pathogène} & \textbf{Symptôme caractéristique} & \textbf{Prévention} \\
\hline
VIH/Sida & Virus VIH (rétrovirus) & Affaiblissement progressif des défenses immunitaires (infections opportunistes) & Préservatif, dépistage, fidélité, non-partage d'objets tranchants \\
\hline
Gonococcie & Bactérie \emph{Neisseria gonorrhoeae} (gonocoque) & Écoulement purulent (« chaude-pisse »), brûlures urinaires & Préservatif, dépistage et traitement des partenaires \\
\hline
Syphilis & Bactérie \emph{Treponema pallidum} (tréponème) & Chancre induré indolore au point d'entrée & Préservatif, dépistage, traitement antibiotique \\
\hline
HPV (papillomavirus) & Virus HPV & Verrues génitales (crêtes de coq) ; certains types cancérigènes (cancer du col de l'utérus) & Préservatif (protection partielle), vaccination des adolescentes \\
\hline
\end{tabularx}
\end{center}
\textbf{Point de vigilance :} ne pas confondre agent et maladie (le VIH est le virus, le Sida est le stade avancé de la maladie) ; le préservatif reste le moyen de prévention commun à toutes les IST.
\end{corrige}

\begin{corrige}[Exercice 6 — La double protection]
\begin{enumerate}
\item Le préservatif est dit moyen de \cle{double protection} car il assure simultanément deux fonctions : une fonction \textbf{contraceptive} (barrière mécanique empêchant les spermatozoïdes d'atteindre les voies génitales féminines, donc la fécondation) et une fonction \textbf{prophylactique} (barrière empêchant le contact avec le sang et les sécrétions génitales, donc la transmission des agents des IST, dont le VIH).
\item La pilule est une méthode \textbf{hormonale} : les \oe{}stroprogestatifs qu'elle contient bloquent l'ovulation en inhibant la sécrétion de FSH et de LH par l'hypophyse (rétrocontrôle négatif sur l'axe hypothalamo-hypophysaire). Elle n'a \textbf{aucune action barrière} : les virus et bactéries des IST passent librement lors du rapport. Un couple d'adolescents, exposé au risque d'IST, n'est donc protégé que contre la grossesse, pas contre les infections.
\item Exemple de conseil : « La pilule évite la grossesse mais pas les IST : utilise en plus un préservatif à chaque rapport, c'est la seule façon d'être protégé contre les deux risques à la fois. »
\end{enumerate}
\end{corrige}

\begin{corrige}[Exercice 7 — Classer les méthodes de contraception]
\begin{enumerate}
\item \textbf{Classement :}
\begin{itemize}
\item Méthode naturelle : calendrier (méthode d'Ogino) ;
\item Méthode mécanique (barrière) : préservatif ;
\item Méthode chimique : spermicide ;
\item Méthodes hormonales : pilule, implant ;
\item Le stérilet au cuivre est un dispositif intra-utérin à effet spermicide (le cuivre est toxique pour les spermatozoïdes) : on le classe à part.
\end{itemize}
\item \textbf{Principes de fonctionnement :}
\begin{itemize}
\item Calendrier (Ogino) : abstinence pendant la période fertile, repérée à partir de la date présumée de l'ovulation ;
\item Préservatif : barrière mécanique retenant les spermatozoïdes ;
\item Spermicide : substance chimique détruisant ou immobilisant les spermatozoïdes dans le vagin ;
\item Pilule : hormones (\oe{}stroprogestatifs) bloquant l'ovulation par rétrocontrôle sur l'axe hypothalamo-hypophysaire ;
\item Implant : diffusion lente et continue d'un progestatif bloquant l'ovulation pendant plusieurs années ;
\item Stérilet (au cuivre) : dispositif placé dans l'utérus, dont le cuivre exerce un effet spermicide et rend la muqueuse utérine impropre à la nidation.
\end{itemize}
\item \textbf{Classement de la moins efficace à la plus efficace en usage courant :} calendrier (Ogino) $<$ spermicide seul $<$ préservatif $<$ pilule $<$ stérilet $\approx$ implant. L'implant et le stérilet, méthodes de longue durée indépendantes de l'observance quotidienne, ont les indices de Pearl les plus faibles (moins de 1 grossesse pour 100 femmes-années).
\end{enumerate}
\end{corrige}

\begin{corrige}[Exercice 8 — Sensibilisation à la planification familiale]
\textbf{Exemple de texte attendu (10 à 15 lignes) :}

« Chers camarades, la santé reproductive nous concerne tous. La planification familiale consiste à choisir librement le moment et le nombre des naissances : c'est un droit et une responsabilité. Premier argument scientifique : une grossesse avant 18 ans expose la mère et l'enfant à des risques accrus (bassin immature, hémorragies, prématurité), car l'organisme de l'adolescente n'a pas terminé sa croissance. Deuxième argument : les rapports non protégés transmettent des IST comme le VIH, la gonococcie ou le HPV, qui peuvent entraîner la stérilité ou des cancers ; seul le préservatif protège à la fois des infections et des grossesses. Troisième argument : une naissance non désirée interrompt souvent la scolarité, surtout celle des filles, et appauvrit les familles. Grâce aux méthodes de contraception disponibles dans les centres de santé et à l'information, chacun peut décider en connaissance de cause. Alors agissons : informons-nous, faisons-nous dépister, utilisons le préservatif et parlons-en autour de nous. Un comportement responsable aujourd'hui, c'est un avenir réussi demain ! »

\textbf{Points de vigilance :} le texte doit contenir au moins deux arguments \emph{scientifiques} (physiologie, risques sanitaires), un appel à l'action clair, et rester respectueux et non moralisateur.
\end{corrige}

\begin{corrige}[Exercice 9 — Étude de cas : une grossesse précoce à Bafou]
\begin{enumerate}
\item Une \cle{grossesse précoce} est une grossesse survenant chez une adolescente (10 à 19 ans), avant la complète maturité de son organisme. Amina a 15 ans et est enceinte de trois mois : son cas correspond exactement à cette définition ; la grossesse est de plus non désirée (elle est élève, le père nie sa responsabilité).
\item Causes possibles (deux parmi) : rapports sexuels précoces non protégés ; méconnaissance ou non-utilisation des méthodes de contraception ; pauvreté favorisant les relations transactionnelles ; absence d'éducation à la santé reproductive ; pression sociale ou mariage précoce.
\item Conséquences :
\begin{itemize}
\item \textbf{Sanitaire :} risque accru de complications de la grossesse et de l'accouchement (dystocie, hémorragie, fistule obstétricale), prématurité et faible poids du nouveau-né ;
\item \textbf{Scolaire :} interruption ou abandon de la scolarité (Amina est en Troisième), perte de chances de formation et d'insertion professionnelle ;
\item \textbf{Sociale :} rejet par le père de l'enfant, stigmatisation, mariage forcé précoce, aggravation de la pauvreté de la famille.
\end{itemize}
\item À 15 ans, le bassin osseux d'Amina n'est pas complètement développé (la croissance du bassin se poursuit jusqu'à 18--20 ans). Le passage du f\oe{}tus lors de l'accouchement peut être impossible (dystocie mécanique), entraînant un travail prolongé, des déchirures, une fistule obstétricale, voire le décès de la mère ou de l'enfant : d'où la nécessité fréquente d'une césarienne.
\item Mesures de prévention (deux parmi) : éducation à la santé reproductive à l'école et dans la communauté ; accès gratuit ou subventionné aux contraceptifs (préservatifs) au centre de santé intégré ; lutte contre les mariages précoces et maintien des filles à l'école ; sensibilisation des garçons à la responsabilité partagée.
\end{enumerate}
\textbf{Barème indicatif (sur 10) :} définition et application au cas : 2 pts ; causes : 1 pt ; trois conséquences bien catégorisées : 3 pts ; justification physiologique (bassin immature $\Rightarrow$ dystocie) : 2 pts ; mesures de prévention : 2 pts.
\textbf{Points de vigilance :} bien distinguer les trois catégories de conséquences ; la justification de la question 4 doit citer explicitement l'immaturité du bassin, pas seulement « elle est trop jeune ».
\end{corrige}

\begin{corrige}[Exercice 10 — Données épidémiologiques et prévention des IST]
\begin{enumerate}
\item Diagramme en bâtons (échelle : 1 cm pour 20 cas) :

\begin{popfigure}
\begin{tikzpicture}
\begin{axis}[width=0.85\linewidth,height=7cm,
ybar,bar width=0.8cm,
xlabel={Année},ylabel={Nouveaux cas d'IST},
symbolic x coords={2020,2021,2022,2023},
xtick=data,ymin=0,ymax=300,
ytick={0,20,40,60,80,100,120,140,160,180,200,220,240,260,280,300},
grid=major,grid style={popline,opacity=0.3},
nodes near coords,nodes near coords style={font=\small,color=popdark}]
\addplot[fill=popblue,draw=popdark] coordinates {(2020,120) (2021,145) (2022,190) (2023,260)};
\end{axis}
\end{tikzpicture}
\legende{Évolution du nombre de nouveaux cas d'IST chez les 15--24 ans dans le district (2020--2023).}
\end{popfigure}

Hauteurs des bâtons à l'échelle demandée : 2020 : $120/20 = 6$ cm ; 2021 : $145/20 = 7,25$ cm ; 2022 : $190/20 = 9,5$ cm ; 2023 : $260/20 = 13$ cm.

\item Pourcentage d'augmentation entre 2020 et 2023 :
\[ \frac{260 - 120}{120} \times 100 = \frac{140}{120} \times 100 \approx 116,7\ \%. \]
Le nombre de cas a plus que doublé (augmentation d'environ \textbf{117 \%}).

\item \textbf{Évolution :} augmentation régulière et forte du nombre de nouveaux cas chaque année. \textbf{Explications possibles (deux parmi) :} multiplication des rapports non protégés chez les jeunes ; méconnaissance des moyens de prévention ; amélioration de l'accès au dépistage (plus de cas détectés et déclarés) ; urbanisation et mobilité accrues.

\item Modes de transmission et prévention :
\begin{itemize}
\item Rapports sexuels non protégés $\rightarrow$ préservatif (masculin ou féminin), abstinence, fidélité mutuelle ;
\item Sang contaminé (transfusion, partage de seringues ou d'objets tranchants : lames, aiguilles de tatouage) $\rightarrow$ utilisation de matériel stérile à usage unique, sang testé ;
\item Transmission mère--enfant (grossesse, accouchement, allaitement) $\rightarrow$ dépistage de la femme enceinte et traitement préventif (PTME).
\end{itemize}

\item Le dépistage précoce et gratuit permet de détecter une IST avant l'apparition des complications (stérilité, Sida), de traiter rapidement le porteur et ses partenaires, et de rompre la chaîne de transmission. La gratuité lève l'obstacle financier pour les adolescents, qui n'ont pas de revenus propres.
\end{enumerate}
\textbf{Barème indicatif (sur 10) :} diagramme correct avec échelle et légendes : 3 pts ; calcul du pourcentage : 1 pt ; description et explications : 2 pts ; trois couples transmission/prévention : 3 pts ; intérêt du dépistage : 1 pt.
\textbf{Points de vigilance :} respecter l'échelle et graduer les axes ; dans le calcul, diviser par la valeur \emph{initiale} (120), pas par 260.
\end{corrige}

\begin{corrige}[Exercice 11 — Sujet de synthèse : santé reproductive et comportement responsable]
\begin{enumerate}
\item Le document 1 révèle une \textbf{forte prévalence des grossesses précoces} : environ une adolescente camerounaise de 15 à 19 ans sur quatre a déjà commencé une vie de mère. C'est un problème de santé publique majeur, car ces grossesses exposent à des risques sanitaires élevés et compromettent la scolarisation des filles.
\item Le document 2 montre que la \textbf{méconnaissance des méthodes de contraception} est le facteur favorisant : si 60 \% des élèves connaissent le préservatif, seulement 20 \% savent citer une autre méthode. Ce déficit d'information conduit à des rapports non protégés.
\item Rôle des actions du document 3 :
\begin{itemize}
\item \textbf{Campagnes de dépistage gratuit du VIH :} permettent à chacun de connaître son statut sérologique, de se faire soigner tôt et d'éviter de transmettre le virus ;
\item \textbf{Distribution de préservatifs :} donne accès au seul moyen de double protection, réduisant à la fois les IST et les grossesses non désirées.
\end{itemize}
\item Comparaison (exemple) :
\begin{itemize}
\item \textbf{Pilule (hormonale) :} les hormones bloquent l'ovulation par rétrocontrôle sur l'axe hypothalamo-hypophysaire ; très efficace si prise régulièrement, mais sans protection contre les IST et exigeant une observance quotidienne ;
\item \textbf{Préservatif (mécanique) :} barrière retenant les spermatozoïdes ; efficace dès le premier usage, sans ordonnance, et seul à protéger aussi contre les IST — avantage décisif pour un adolescent.
\end{itemize}
\item \textbf{Texte structuré (exemple, 15 à 20 lignes) :}

« Adopter un comportement responsable en matière de santé reproductive, c'est d'abord se protéger soi-même. En évitant les rapports non protégés ou en utilisant le préservatif, l'adolescent se préserve des IST, dont certaines, non traitées, conduisent à la stérilité, et il évite une grossesse précoce qui menacerait sa santé : avant 18 ans, le bassin n'est pas mature et l'accouchement est risqué. C'est ensuite protéger son avenir : une adolescente mère à 15 ans quitte le plus souvent l'école, comme le montre la proportion élevée de mères adolescentes au Cameroun, et voit ses chances d'insertion professionnelle s'effondrer. Le comportement responsable passe aussi par l'information : connaître les méthodes de contraception, se faire dépister gratuitement, dialoguer avec son partenaire et les personnels de santé. À l'échelle de la communauté, des jeunes en bonne santé, instruits et maîtres de leur fécondité réduisent les dépenses de santé des familles et de l'État, limitent la propagation des IST et des abandons d'enfants, et contribuent à des naissances désirées et espacées, mieux prises en charge. La planification familiale devient ainsi un levier de développement. Enfin, la responsabilité est partagée : garçons et filles doivent assumer ensemble les conséquences de leurs choix. Se protéger, s'informer, se dépister, respecter l'autre : telles sont les clés d'une adolescence épanouie et d'une société plus saine et plus prospère. »
\end{enumerate}
\textbf{Barème indicatif (sur 20) :} question 1 : 2 pts ; question 2 : 2 pts ; question 3 : 4 pts ; question 4 : 4 pts ; texte de synthèse (structure, arguments scientifiques, lien individu/communauté, qualité de la langue) : 8 pts.
\textbf{Points de vigilance :} en question 4, citer le \emph{mode d'action} (pas seulement le nom) de chaque méthode ; dans le texte final, mobiliser explicitement les trois documents et des connaissances de physiologie ; éviter les généralités moralisatrices sans argument scientifique.
\end{corrige}

\newpage

\coursec{Fiche bilan}

\begin{popfigure}
\begin{tikzpicture}[mindmap, grow cyclic, every node/.style={concept, text=white, font=\small\bfseries},
  root concept/.append style={concept color=popdark, font=\small\bfseries},
  level 1/.append style={level distance=3.4cm, sibling angle=72},
  level 2/.append style={level distance=2.2cm, sibling angle=45, font=\scriptsize}]
\node[root concept]{Santé reproductive des adolescent(e)s}
  child[concept color=popblue]{ node{Puberté et reproduction}
    child{ node{Hormones : FSH, LH, testostérone, {\oe}strogènes, progestérone} }
    child{ node{Cycle menstruel (28 j)} }
  }
  child[concept color=poporange]{ node{Grossesses précoces}
    child{ node{Causes : rapports non protégés, pauvreté, ignorance} }
    child{ node{Conséquences : santé, scolarité, social} }
  }
  child[concept color=poppink]{ node{IST}
    child{ node{VIH/Sida, blennorragie, syphilis, chlamydiose} }
    child{ node{Prévention : abstinence, fidélité, préservatif} }
  }
  child[concept color=popgreen]{ node{Contraception}
    child{ node{Naturelles : Ogino, température} }
    child{ node{Artificielles : pilule, préservatif, DIU, implant} }
  }
  child[concept color=popgold]{ node{Planification familiale}
    child{ node{Espacement des naissances} }
    child{ node{Comportement responsable} }
  };
\end{tikzpicture}
\legende{Carte mentale de la leçon : les cinq grandes branches de la santé reproductive des adolescents.}
\end{popfigure}

\begin{bilan}
\textbf{1. Définitions clés}
\begin{itemize}
  \item \cle{Puberté} : période de transformation de l'organisme de l'enfant vers l'âge adulte, marquée par l'apparition des caractères sexuels secondaires et l'acquisition de la capacité de se reproduire (10--14 ans chez la fille, 12--16 ans chez le garçon).
  \item \cle{Cycle menstruel} : ensemble de transformations cycliques de l'ovaire et de l'utérus, d'une durée moyenne de 28 jours, sous contrôle hormonal ; l'\cle{ovulation} a lieu vers le 14\ieme{} jour.
  \item \cle{Grossesse précoce} : grossesse survenant chez une adolescente de moins de 19 ans, souvent non désirée.
  \item \cle{IST} : infection transmise principalement par les rapports sexuels non protégés (VIH/Sida, blennorragie, syphilis, chlamydiose, hépatite B, herpès génital).
  \item \cle{Contraception} : ensemble des méthodes empêchant la fécondation ou la nidation.
  \item \cle{Planification familiale} : ensemble des moyens permettant aux couples de choisir librement le nombre de leurs enfants et l'espacement des naissances.
\end{itemize}

\textbf{2. Hormones et organes à connaître par c{\oe}ur}
\begin{center}
\begin{tabularx}{\linewidth}{|l|X|X|}\hline
\rowcolor{popblueL} \textbf{Hormone} & \textbf{Origine} & \textbf{Rôle essentiel} \\ \hline
FSH & Hypophyse & Croissance des follicules (fille) ; spermatogenèse (garçon) \\ \hline
LH & Hypophyse & Déclenche l'ovulation ; stimule la sécrétion de testostérone \\ \hline
{\OE}strogènes & Ovaires & Caractères sexuels secondaires ; reprise de la prolifération utérine \\ \hline
Progestérone & Corps jaune & Prépare et maintient la muqueuse utérine (nidation) \\ \hline
Testostérone & Testicules & Caractères sexuels secondaires masculins ; spermatogenèse \\ \hline
\end{tabularx}
\end{center}

\textbf{3. Méthodes de contraception : tableau récapitulatif}
\begin{center}
\begin{tabularx}{\linewidth}{|l|X|X|}\hline
\rowcolor{popgreenL} \textbf{Catégorie} & \textbf{Exemples} & \textbf{Précision} \\ \hline
Naturelles & Méthode Ogino (calendrier), courbe de température, retrait & Peu fiables ; ne protègent pas des IST \\ \hline
Mécaniques / barrières & Préservatif masculin et féminin, diaphragme & Seul le préservatif protège aussi des IST \\ \hline
Chimiques & Pilule, implant, injectable, spermicides & Efficaces si bien utilisées ; sur prescription \\ \hline
Chirurgicales & Ligature des trompes, vasectomie & Définitives \\ \hline
\end{tabularx}
\end{center}

\textbf{4. Méthodes express}
\begin{itemize}
  \item \emph{Analyser un cas de grossesse précoce} : identifier la cause (rapport non protégé, méconnaissance, pression sociale) $\to$ citer les conséquences médicales (fistules, avortements clandestins, mortalité maternelle), scolaires (déscolarisation) et sociales (rejet, précarité) $\to$ proposer des solutions (information, contraception, soutien).
  \item \emph{Justifier la prévention des IST} : retenir la triade \cle{abstinence -- fidélité -- préservatif}, complétée par le dépistage et la vaccination (hépatite B, HPV).
  \item \emph{Retenir} : le préservatif est la seule méthode à double protection (grossesse + IST).
\end{itemize}
\end{bilan}

\begin{aretenir}[Checklist avant le Probatoire]
\begin{itemize}
  \item[$\square$] Je sais définir la puberté et citer ses manifestations chez la fille et chez le garçon.
  \item[$\square$] Je connais le rôle des hormones FSH, LH, {\oe}strogènes, progestérone et testostérone.
  \item[$\square$] Je sais décrire les phases du cycle menstruel et situer l'ovulation (vers le 14\ieme{} jour).
  \item[$\square$] Je sais définir une grossesse précoce et citer au moins trois causes.
  \item[$\square$] Je sais citer les conséquences médicales, scolaires et sociales des grossesses précoces.
  \item[$\square$] Je connais au moins quatre IST et leurs agents (VIH, blennorragie, syphilis, chlamydiose).
  \item[$\square$] Je sais justifier les trois moyens de prévention des IST : abstinence, fidélité, préservatif.
  \item[$\square$] Je sais que le préservatif est la seule méthode protégeant à la fois des grossesses et des IST.
  \item[$\square$] Je sais classer les méthodes de contraception (naturelles, barrières, chimiques, chirurgicales).
  \item[$\square$] Je sais définir la planification familiale et expliquer l'intérêt de l'espacement des naissances.
  \item[$\square$] Je sais analyser un cas concret et proposer un comportement responsable en santé reproductive.
\end{itemize}
\end{aretenir}

\findecours

\end{document}
